\documentclass[11pt]{article}
\usepackage[margin=1in]{geometry}
\usepackage{amsmath,amssymb,amsthm,booktabs,array}
\usepackage{microtype}
\usepackage{url}
\providecommand{\doi}[1]{\mbox{doi:}\,\url{https://doi.org/#1}}
\newtheorem{theorem}{Theorem}[section]
\newtheorem{proposition}[theorem]{Proposition}
\newtheorem{corollary}[theorem]{Corollary}
\newtheorem{lemma}[theorem]{Lemma}
\newtheorem{definition}[theorem]{Definition}
\theoremstyle{remark}
\newtheorem{remark}{Remark}[section]
\newcommand{\Gr}{\mathcal{G}}
\newcommand{\Le}{\Lambda_E}
\newcommand{\Cnorm}{\bigl\|\Le^{-1}H\bigr\|_2^2}
\newcommand{\minr}[1]{\sigma_{\min}(#1)}
\begin{document}
\title{\textbf{Identifiability of the correction term:\\
Worst-Case Leakage Is Governed by the Principal Angles of the Dictionary}}
\author{%
Xiangrui Meng\thanks{School of Mathematics and Physics, Shanghai Normal University, Shanghai,
China. Lead (human) author, ORCID 0009-0006-8902-0519.}%
\and
DeepSeek V4.1 Flash\thanks{AI co-author (Agent A). Role: experiments, certificate construction,
falsification of the manuscript's own claims.}%
\and
Qwen 3.8 Flash\thanks{AI co-author (Agent B). Role: derivation, theorems, numerics, literature
verification.}%
}
\date{\today}
\maketitle

\begin{abstract}
A scientific-ML model is written as a prior operator plus a learned correction,
$A^\star=\sum_i\theta_i^\star E_i+\sum_j\psi_j^\star F_j+\Delta_{\rm hidden}$.
Fitting the coefficients by least squares from $n$ input--output pairs answers a
question that is not about accuracy but about \emph{attribution}: how much of an
unmodeled operator $\Delta_{\rm hidden}$ is forced into the correction block?
We show this is a question in matrix geometry with an exact answer. Four identities
hold for the operator (Frobenius) Gram $\Gr$ of the dictionary: its block form, an
isometry, a square identity, and $\det\Gr=\det(\Le)\sin^{2q}\alpha$. They give a
two-sided bound on $\lambda_{\min}(\Gr)$ whose intrinsic quantity is the frame
invariant $\Cnorm$, and a sampling bound whose remainder is \emph{independent of the
principal angle}. The resulting certificate
\[
  \|\hat C-C^\star\|_F\;\le\;
  \frac{\|\Delta\|_F+\sigma\sqrt d}{\sin\alpha_{\min}}\,
  \sqrt{1+2\cos^2\alpha_{\min}\,\Cnorm}\,
  \Bigl(1-\tfrac{c}{2n}-\tfrac{0.88}{\sqrt n}\Bigr)^{-1}
\]
contains no fitted constant, is attained within $1\%$ in the worst sampled case, and
holds with the sample-size factor set to $1$ in $4320/4320$ adversarial trials spanning
$n=30\ldots1800$ and $\alpha=2^\circ\ldots85^\circ$. We also measure the identity
$\text{err}/\text{bound}=q_4q_1q_2q_3$ and show that the sample-size remainder
\emph{cancels structurally} ($q_1q_2\approx1$) whenever the perturbing operator is
realized through the same data matrix whose Gram it perturbs: the tightness of the
certificate is a property of $\Gr$ alone. Since $\lambda_{\min}(\Gr)$ attains the upper
branch sample by sample, the dictionary factor $q_4$ is a closed-form quantity, and with
the block-mass factor $q_3$ the observed looseness collapses to a single dictionary
scalar, $(1+2\cos^2\alpha\,\Cnorm)^{-1/2}$, which reproduces the fitted ratios of the
adversarial table to within $1\%$ and needs no calibration at any angle. The upper side is
then not a bound but a reduction: $\lambda_{\min}(\Gr)$ is the unique root of a scalar
equation in a $q\times q$ matrix pencil, and that root is the closed form above \emph{exactly}
when $p=q$. Section~\ref{sec:notclaims} states what we do not claim, including the real-data leg
of this project, which was retired by its own measurements rather than by argument.
\end{abstract}

\vspace{2pt}
\noindent\textbf{Keywords.} saddle-point matrix; arrowhead secular equation; principal angles
between subspaces; identifiability; attribution error; scientific machine learning; certified
a posteriori bound; Gauss quadrature.
\quad\textbf{AMS 2020 subject classifications.} 15A18, 65F15, 62J05, 68T07.

\section*{Introduction}

Hybrid models write a dynamical operator as something known plus something learned,
$A^\star=\sum_i\theta_i^\star E_i+\sum_j\psi_j^\star F_j+\Delta_{\rm hidden}$, fit both
blocks from the same $n$ input--output pairs, and are then validated by predictive
accuracy. Accuracy answers a different question from the one the model is asked to
address. When the correction block is reported as a \emph{finding} --- a closure term, a
missing forcing, a diagnosed coupling --- the claim is about attribution: what fraction
of an unmodeled operator $\Delta_{\rm hidden}$ was \emph{forced into} $\hat\psi$ by the
fit itself. A model can forecast perfectly and attribute nothing.

This note shows that attribution is a question in matrix geometry with an answer that is
computable \emph{before any data are collected}. For a dictionary whose correction
directions make principal angle $\alpha$ with the span of the prior block, four
identities hold for the operator (Frobenius) Gram $\Gr$: a block form, an isometry, a
square identity, and $\det\Gr=\det(\Le)\sin^{2q}\alpha$. They pin
$\lambda_{\min}(\Gr)$ between two branches whose intrinsic quantity is the frame
invariant $\Cnorm$, and the resulting leakage factor is $1/\sin\alpha_{\min}$ times a
correction that never exceeds $\sqrt{1+2\cos^2\alpha\,\Cnorm}$. Nothing in that chain
depends on $n$, on the noise level $\sigma$, or on a fitted constant.

The second point is methodological, and it is why we insist on the constant-free form.
The alternative in the literature is an empirically calibrated envelope
$c_1\|\Delta\|+c_2\sigma$. The two constants do not carry the same dependence on anything:
Corollary~\ref{cor:dilution} makes the model-form one a $\sqrt q/d$-type dilution factor
(the average case is a dimensional artifact, not physics; the $\sqrt q$ is block aggregation),
while the noise one is the
Gaussian quantile $z_{1-\delta/2}/\sqrt n$, with no $d$ at all and no mechanism in
between. Fitting a single pair $(c_1,c_2)$ averages two regimes whose scalings disagree
and tells a reader nothing about which regime their experiment sits in. The disagreement is
not a technicality: in the two-arm least-squares experiment behind this note (design
by Agent A, audit \texttt{b40}, \texttt{b41}, \texttt{b40b}, $q=3$), the measured model-form constant is
the random-projection fraction $\mathbb{E}\|P_{\mathcal F}\Delta\|/\|\Delta\|$ --- a Beta mean
$B(q/2+\tfrac12,(d^{2}-q)/2)/B(q/2,(d^{2}-q)/2)$, leading term $\kappa(q)\sqrt q/d$ with $\kappa$ the
$\chi_q$ mean factor --- matched to within $8\%$ over $d=6\ldots20$ and $6\%$ over $q=1\ldots6$, and
flat in $n$ to within $5\%$ over $n=150\ldots2400$; the noise
constant is instead $\sqrt{q/n}$ to within $13\%$. Their ratio is $d/\sqrt n$, measured from
$0.90$ at $n=150$ to $0.21$ at $n=2400$. One fitted constant cannot describe both, and a
table that reports one is reporting whichever regime its sample size happened to sit in.

\medskip
\noindent\textbf{Contributions, tagged by status.}
\begin{itemize}\itemsep1pt
\item Exact determinant law $\det\Gr=\det(\Le)\sin^{2q}\alpha$, and three companion
identities \hfill \emph{[proved; residuals $\le10^{-12}$]}
\item Two-sided bound on $\lambda_{\min}(\Gr)$, with the direction of the inequality
chosen so that a wrong dictionary always \emph{enlarges} the certificate
\hfill \emph{[proved; $6000$ draws over three scripts, $0$ violations on either side]}
\item The upper side sharpens to a closed form $\operatorname{cf}(t,\alpha)$, realised as the
Rayleigh minimum on an explicit plane; and $\lambda_{\min}(\Gr)$ reduces \emph{exactly} to a
one-dimensional root find over a $q\times q$ matrix pencil
\hfill \emph{[both proved; the exact formula to $1.8\cdot10^{-11}$ on $6480$ dictionaries, the
closed form \emph{identically} $\lambda_{\min}(\Gr)$ when $p=q$; the root find survives an arbitrary
principal-cosine spectrum, the bound \emph{direction} of the closed form does not]}
\item The conjecture that the plane direction is the top singular direction is
\emph{false} for $q<p$, and the counterexample was found only by replacing random sampling with
an exhaustive circle search \hfill \emph{[disproved; defect $\le10^{-6}$ relative]}
\item Sampling remainder $O(\sqrt{(p+q)/n})$ that does not know about $\alpha$, and a
refined one-sided $1/n$ bias term computable from the dictionary \hfill \emph{[proved in
shape; the constants are calibrated, not derived]}
\item Worst-case leakage bound with constant \emph{exactly $1$}, plus a construction
attaining it \hfill \emph{[proved, subject to $n>d$]}
\item Certificate $(\star)$ and its $4320$-trial adversarial test; the analytic form of
its remaining slack \hfill \emph{[measured; the closed form matches the exact dictionary
factor over $4400$ dictionaries and eleven angles, worst per-sample deviation $0.004$]}
\item The achieved ratio factorizes as $q_4q_1q_2q_3$ (dictionary, sampling, realization,
block mass), and with $q_3$ included the observed looseness of $(\star)$ collapses to the
single dictionary scalar $(1+2\cos^2\alpha\,\Cnorm)^{-1/2}$ \hfill \emph{[identity verified
per trial; the collapse is measured, not proved]}
\item Three structural negative results: the sample-size margin cancels and must not be
tabulated; the $\alpha$-uniformity of the bias is \emph{not} proved although it is
flat over the scanned grid; and the $q_1q_2$ cancellation is only partial under a
correlated (trajectory) design, where it reaches $0.975$ at $\rho=0.95$ --- always in the
safe, looser direction \hfill \emph{[measured, with the grid stated]}
\end{itemize}

\medskip
Sections~\ref{sec:setting}--\ref{sec:identities} hold the geometry,
\ref{sec:twosided}--\ref{sec:sampling} the two-sided bound and the sampling bridge,
\ref{sec:leakage}--\ref{sec:cert} the worst case and the certificate,
\ref{sec:cancellation} the measured factorization, and \ref{sec:notclaims} the
limitations, including seventeen claims we retracted ourselves and two open items.
Section~\ref{sec:position}
states what is not new here.

\section{Setting, notation, and the one geometric quantity}\label{sec:setting}

Let $A^\star\in\mathbb{R}^{d\times d}$ be the true linear operator and let a
\emph{dictionary} be a pair of finite families of operators,
\begin{equation}\label{eq:decomp}
  A^\star=\underbrace{\sum_{i=1}^{p}\theta_i^\star E_i}_{\text{prior block }P}
        +\underbrace{\sum_{j=1}^{q}\psi_j^\star F_j}_{\text{correction block }C^\star}
        +\Delta ,
\end{equation}
where $\Delta$ is the \emph{model-form error}: the part of the dynamics that the chosen
dictionary does not span. We observe
\begin{equation}\label{eq:data}
  Y=A^\star U+\Sigma,\qquad
  U\in\mathbb{R}^{d\times n}\ \text{i.i.d.\ standard Gaussian},\quad
  \Sigma\in\mathbb{R}^{d\times n}\ \text{i.i.d.\ } \mathcal N(0,\sigma^2),
\end{equation}
and fit $(\hat\theta,\hat\psi)$ by least squares over the design matrix
\begin{equation}\label{eq:design}
  G=\bigl[\,\operatorname{vec}(E_1U)\ \cdots\ \operatorname{vec}(E_pU)\ \big|\
          \operatorname{vec}(F_1U)\ \cdots\ \operatorname{vec}(F_qU)\,\bigr]
   \in\mathbb{R}^{dn\times(p+q)} .
\end{equation}

All geometry below lives in the vectorized operator space
$\mathbb{R}^{d^2}$ with the Frobenius inner product. Define the \emph{operator Gram}
\begin{equation}\label{eq:opgram}
  \Gr_{ab}=\langle M_a,M_b\rangle_F,\qquad \{M_a\}_{a=1}^{p+q}=\{E_i\}\cup\{F_j\},
\end{equation}
the prior Gram $\Le=(\langle E_i,E_j\rangle_F)\in\mathbb{R}^{p\times p}$, and the
cross-Gram $H=(\langle E_i,a_j\rangle_F)\in\mathbb{R}^{p\times q}$, where $a_j$ is the
component of $F_j$ inside $\operatorname{span}\{E_i\}$ (made precise in
\eqref{eq:angle}). Two abbreviations recur:
$M:=\Le^{-1/2}H$ and $C:=\Le^{-1}H$.

\begin{definition}[dictionary angle]\label{def:angle}
The construction used throughout (and in the experiments of Agent A) is
\begin{equation}\label{eq:angle}
  F_j=\cos\alpha\, a_j+\sin\alpha\, b_j,\qquad
  \{a_j\}_{j=1}^{q}\ \text{orthonormal}\subset\operatorname{span}\{E_i\},\quad
  \{b_j\}_{j=1}^{q}\ \text{orthonormal}\subset\operatorname{span}\{E_i\}^{\perp},
\end{equation}
with all orthonormality in the Frobenius norm of $\mathbb{R}^{d^2}$.
Then $\|F_j\|_F=1$, $\langle F_i,F_j\rangle_F=\delta_{ij}$, and
$\langle E_i,F_j\rangle_F=\cos\alpha\,\langle E_i,a_j\rangle_F$; $\alpha$ is the
principal angle between the two blocks, and $\alpha\to0$ means the correction is
\emph{aligned} with the prior.
\end{definition}

\begin{remark}[where exactness lives]\label{rem:exact}
\eqref{eq:angle} is imposed in the \emph{vectorized} space, hence
$\|F_j\|_F=1$ and $\langle F_i,F_j\rangle_F=\delta_{ij}$ hold sample by sample, not in
expectation. What is \emph{not} exact is $\Le$: the prior operators are random and
mutually correlated, with off-diagonal entries of size $O(\sqrt d)$ relative to the
diagonal. Every looseness reported in this chapter traces back to $\kappa(\Le)$, not to
the correction block. This is a correction of an earlier draft of ours, which assumed
the opposite.
\end{remark}

The single scalar that controls the certificate is
\begin{equation}\label{eq:frameinv}
  \Cnorm=\|\Le^{-1}H\|_2^2=\|M\|_2^2 : \text{the frame condition functional of the dictionary}.
\end{equation}
It is computable from the dictionary alone --- no data, and (Theorem~\ref{thm:twosided}
remarks) no $\alpha$. Replacing it by the quantity $\rho_E^{-2}=\|E_i\|_F^{-2}$ that an
earlier version of this chapter used is wrong except when $\Le\approx\rho_E^2 I$; the
measured gap is $\Cnorm\cdot d=1.005$--$1.073$ for $d\ge16$ but $1.5$--$3.0$ at $d=4$.
A computation pitfall worth recording: by Proposition~\ref{prop:isometry},
$\|\Le^{-1/2}H\|_2^2\equiv1$ identically, so the invariant must be formed as a
\emph{solve}, $\Le^{-1}H=$ \texttt{solve(Lambda\_E, H)}, never as a Cholesky
half-inverse times $H$. The latter measures $\|R^{-\top}H\|_2^2$ and returns
$1.05$--$1.07$ at $d=8$, where the true value is $0.17$.

\section{Four exact identities}\label{sec:identities}

\begin{proposition}[block Gram]\label{prop:block}
With \eqref{eq:angle},
$\Gr=\begin{pmatrix}\Le & \cos\alpha\,H\\ \cos\alpha\,H^{\top} & I_q\end{pmatrix}.$
\end{proposition}

\begin{proposition}[isometry]\label{prop:isometry}
$H^{\top}\Le^{-1}H=M^{\top}M=I_q$; consequently $K:=MM^{\top}$ is the orthogonal
projection of rank $q$ onto $\operatorname{ran}\Le^{-1/2}H$.
\end{proposition}

\begin{proposition}[square identity]\label{prop:square}
For every coefficient pair $(\theta,w)$,
\begin{equation}\label{eq:square}
  \theta^{\top}\Le\theta+2\cos\alpha\,\theta^{\top}Hw+\|w\|^2
  =\bigl\|\Le^{1/2}\theta+\cos\alpha\,Mw\bigr\|_2^2+\sin^2\alpha\,\|w\|_2^2 .
\end{equation}
The left-hand side is $\|\sum_i\theta_iE_i+\sum_jw_jF_j\|_F^2$.
\end{proposition}

\begin{proof}
Expand the right-hand side: $\theta^{\top}\Le\theta+2\cos\alpha\,\theta^{\top}\Le^{1/2}
Mw+\cos^2\alpha\,\|Mw\|^2+\sin^2\alpha\|w\|^2$. Now $\Le^{1/2}M=H$ and, by
Proposition~\ref{prop:isometry}, $\|Mw\|^2=w^{\top}M^{\top}Mw=\|w\|^2$.
\end{proof}

\begin{theorem}[determinant law]\label{thm:det}
$\det\Gr=\det(\Le)\,\sin^{2q}\alpha$.
\end{theorem}

\begin{proof}
Take the Schur complement of the $I_q$ block in Proposition~\ref{prop:block}:
$I_q-\cos^2\alpha\,H^{\top}\Le^{-1}H=I_q-\cos^2\alpha\,M^{\top}M=\sin^2\alpha\,I_q$ by
Proposition~\ref{prop:isometry}. Hence
$\det\Gr=\det(\Le)\det(\sin^2\alpha\,I_q)=\det(\Le)\sin^{2q}\alpha$.
\end{proof}

\begin{corollary}\label{cor:volume}
Each correction direction contributes one factor $\sin^2\alpha$ of identifiability loss;
the generalized volume of a $q$-dimensional correction block collapses as
$\sin^{2q}\alpha$.
\end{corollary}

\noindent\textbf{Numerical status.} Measured against a directly computed Frobenius Gram
of $\{E_i\}\cup\{F_j\}$, the block form \eqref{eq:opgram} agrees to
$4.0\!\times\!10^{-16}$, $5.8\!\times\!10^{-16}$, $7.4\!\times\!10^{-16}$ at $d=8,16,32$;
the isometry residual is $\le4.8\!\times\!10^{-16}$; and the determinant law satisfies
$\bigl|\det\Gr/(\det\Le\,\sin^{2q}\alpha)-1\bigr|\le1.6\!\times\!10^{-12}$ over
$d=3\ldots20$, $\alpha=1^\circ\ldots85^\circ$. The four identities are exact, not
asymptotic; the numbers above are round-off. (Script: \texttt{b11\_theorem\_remainder.py}.)

\section{The dictionary side: a two-sided bound for $\lambda_{\min}(\Gr)$}
\label{sec:twosided}

Parametrize the quadratic form of \eqref{eq:square} by the free variable
$z=\Le^{1/2}\theta+\cos\alpha\,Mw$, so that $\theta=\Le^{-1/2}z-\cos\alpha\,\Le^{-1/2}Mw
=\Le^{-1/2}z-\cos\alpha\,Cw$. Then $\lambda_{\min}(\Gr)$ is the Rayleigh quotient
\begin{equation}\label{eq:rayleigh}
  \lambda_{\min}(\Gr)=\min_{(z,w)\ne0}
  \frac{\|z\|^2+\sin^2\alpha\,\|w\|^2}
       {\|\Le^{-1/2}z-\cos\alpha\,Cw\|^2+\|w\|^2}.
\end{equation}

\begin{theorem}[two-sided bound]\label{thm:twosided}
For every $d,p,q$ and every $\alpha\in(0,\pi/2]$,
\begin{equation}\label{eq:twosided}
  \begin{gathered}
  \min\!\Bigl(\tfrac{\lambda_{\min}(\Le)}{2}\;,\;
             \tfrac{\sin^2\alpha}{1+2\cos^2\alpha\,\|\Le^{-1}H\|_2^2}\Bigr)
  \ \le\ \lambda_{\min}(\Gr)\ \le\\
  \le\ \frac{\sin^2\alpha}{1+\cos^2\alpha\,\|\Le^{-1}H\|_2^2}.
  \end{gathered}
\end{equation}
\end{theorem}

\begin{proof}
\emph{Upper bound.} In \eqref{eq:rayleigh} restrict to the feasible choice $z=0$ and any
unit $w$: the quotient becomes
$\sin^2\alpha/(\cos^2\alpha\,\|Cw\|^2+1)$, whose minimum over $\|w\|=1$ is the
right-hand side of \eqref{eq:twosided} at $w$ aligned with the top singular direction of
$C$. \emph{Lower bound.} For the denominator, Young's inequality with $\varepsilon=1$
gives $\|\Le^{-1/2}z-\cos\alpha\,Cw\|^2\le2\|\Le^{-1/2}z\|^2+2\cos^2\alpha\|Cw\|^2
\le2\lambda_{\min}(\Le)^{-1}\|z\|^2+2\cos^2\alpha\,\Cnorm\|w\|^2$, hence with
$c_1:=2/\lambda_{\min}(\Le)$ and $d_2:=1+2\cos^2\alpha\,\Cnorm$ the quotient in
\eqref{eq:rayleigh} is at least
$(a+\sin^2\alpha\,b)/(c_1a+d_2b)$ where $a=\|z\|^2$, $b=\|w\|^2$. Writing this as
$(c_1a)\cdot(1/c_1)+(d_2b)\cdot(\sin^2\alpha/d_2)$ over $c_1a+d_2b$ exhibits it as a
weighted average of $1/c_1$ and $\sin^2\alpha/d_2$, so it is bounded below by
$\min(1/c_1,\sin^2\alpha/d_2)$, which is the left-hand side of \eqref{eq:twosided}.
\end{proof}

\begin{remark}[the inequality direction is the safety-critical part]
\label{rem:direction}
An earlier version of this chapter stated
$\minr{\Gr^{1/2}}=\sin\alpha/\sqrt{1+\cos^2\alpha\,\|C\|_2^2}$ as an \emph{equality}.
The right-hand side of Theorem~\ref{thm:twosided} is an \textbf{upper} bound on
$\lambda_{\min}(\Gr)$, i.e.\ an upper bound on $\minr{\Gr^{1/2}}$. Substituting an upper
bound on a smallest singular value into
$\|\hat C-C^\star\|\le\|r\|/\minr G$ \emph{underestimates} the error --- the unsafe
direction. Quantitatively, the ratio of true to this surrogate is $0.112$ at the corner
$(d=3,\alpha=85^\circ)$, i.e.\ a factor-of-nine underestimation, while in the working
region $\alpha\le20^\circ$, $d\ge6$ it is $\ge0.97$. Certificates must use the
\emph{lower}-bound multiplier
$\sqrt{1+2\cos^2\alpha\,\Cnorm}/\sin\alpha$, which at $d=10$ costs a factor $1.08$:
essentially free.
\end{remark}

\begin{remark}[scaling invariance]
$E_i\mapsto sE_i$ gives $H\mapsto sH$, $\Le\mapsto s^2\Le$, hence
$\Cnorm\mapsto\Cnorm/s^2$ and the multiplier $\to1$. Enlarging the prior dictionary can
therefore manufacture identifiability \emph{in coefficient space} while leaving the
operator-space error untouched. This is the analytic explanation of the
reparameterization invariance observed numerically (error invariant for
$s=0.01\ldots10$).
\end{remark}

\begin{remark}[case distinction in the small-$\alpha$ regime]
The branch $\lambda_{\min}(\Le)/2$ does not shrink with $\alpha$, so for small $\alpha$
the $\sin^2\alpha$ branch controls. The lower bound is tight in the regime that matters
($\lambda_{\min}(\Le)/2$-versus-true ratio $0.84$--$1.00$ at $d=8$, $0.94$--$1.00$ at
$d=20$) and loose only at $d=3$ ($0.50$--$0.58$ at $\alpha=85^\circ$), consistent with
Remark~\ref{rem:exact}: looseness is carried by $\kappa(\Le)$.
\end{remark}

\noindent\textbf{Numerical status.} $d\in\{3,4,6,8,12,20\}\times
\alpha\in\{1,5,20,45,70,85\}^\circ\times25$ samples $=900$ draws: zero violations of
either side. (Script: \texttt{b13\_two\_sided\_bound.py}.)

\noindent\textbf{The upper side is attained, sample by sample.} A sharper test
(2400 further draws, $d=8$, $p=4$, $q=3$, $\alpha\in\{1,2,5,10,20,45,70,85\}^\circ$)
compares each draw's $\lambda_{\min}(\Gr)$ with \emph{its own} upper branch:
zero violations, and the ratio $\lambda_{\min}(\Gr)/\mathrm{UB}$ has mean
$0.99278$--$0.99999$ and minimum $0.98489$ over the eight cells
(script \texttt{b29\_upper\_bound\_check.py}). The $\sin^2\alpha$ branch is therefore
not merely a bound: it is the right-hand side of an identity up to $1.5\%$, with the
worst case at $\alpha=45^\circ$. Since $\lambda_{\min}(\Gr)$ and $\|\Le^{-1}H\|_2^2$ are
functions of the same draw, the mean-versus-mean comparison is not a test of the bound;
only the per-sample ratio above is, and it is what makes the slack of
\eqref{eq:star} analytic (\S\ref{sec:cancellation}).

\begin{proposition}[the upper side, restricted to a plane]\label{prop:plane}
Write $t:=\|\Le^{-1}H\|_2^2$ and, for $t>0$ and $\alpha\in(0,\pi/2)$,
\begin{equation}\label{eq:cf}
  \operatorname{cf}(t,\alpha)\;:=\;\frac{(1+t)-\sqrt{(1+t)^2-4t\sin^2\alpha}}{2t}\,,
\end{equation}
the smaller root of $Q(\lambda):=t\lambda^2-(1+t)\lambda+\sin^2\alpha$; by continuity
$\operatorname{cf}(0,\alpha)=\sin^2\alpha$. Then
\begin{equation}\label{eq:chain}
  \lambda_{\min}(\Gr)\ \le\ \operatorname{cf}(t,\alpha)\ <\
  \frac{\sin^2\alpha}{1+\cos^2\alpha\,t}\,,
\end{equation}
so $\operatorname{cf}$ is a closed-form upper bound for $\lambda_{\min}(\Gr)$, depending on the
dictionary through $t$ alone, that is strictly sharper than the upper branch of
Theorem~\ref{thm:twosided}.
\end{proposition}

\begin{proof}
In the coordinates $x=\Le^{-1/2}z$ of \eqref{eq:rayleigh}, with $C:=\Le^{-1}H$,
$\lambda_{\min}(\Gr)=\min_{(x,w)\ne0}(x^{\top}\Le x+\sin^2\alpha\|w\|^2)
/(\|x-\cos\alpha Cw\|^2+\|w\|^2)$.
Let $w_0$ be a unit top right-singular vector of $C$, put $u:=\cos\alpha\,Cw_0$, and restrict the
minimiser to the two-dimensional family $x=c\,u$, $w=w_0$: the minimum over a subset is at least the
global one, so it suffices to evaluate it. The isometry $C^{\top}\Le C=I_q$ gives
$u^{\top}\Le u=\cos^2\alpha$ and $\|u\|^2=\cos^2\alpha\,t$, whence the quotient on the family is
$\varphi(c)=(c^2\cos^2\alpha+\sin^2\alpha)/(c^2\|u\|^2-2c\|u\|^2+\|u\|^2+1)$, whose
sublevel condition $\varphi(c)\ge\lambda$ for all $c$ reads
$Q(\lambda)\ge0$ after division by $\cos^2\alpha$ (the leading coefficient of $\varphi-\lambda$ is
positive exactly when $\lambda<1/t$, and $Q(1/t)=\sin^2\alpha-1<0$ places $1/t$ between the roots).
The largest admissible $\lambda$ is therefore the smaller root of $Q$, which is
$\operatorname{cf}(t,\alpha)$. For the second inequality of \eqref{eq:chain}, substituting the upper
branch gives $Q\bigl(\sin^2\alpha/(1+\cos^2\alpha t)\bigr)
=-\cos^2\alpha\sin^2\alpha\,t^2/(1+\cos^2\alpha t)^2<0$, and a quadratic which is negative at a point
puts that point strictly between its roots.
\end{proof}

\begin{remark}[what $\operatorname{cf}$ buys, and what it does not]\label{rem:cf}
The gain over Theorem~\ref{thm:twosided}'s branch is small in $\alpha$-units but systematic: over
$42$ cells $\times\,150$ dictionaries ($d=3\ldots16$, $\alpha=1^\circ\ldots89^\circ$) the ratio
$\operatorname{cf}/\mathrm{UB}$ is $\le1$ on every draw, with worst-case slack
$|\lambda_{\min}(\Gr)/\operatorname{cf}-1|\le2.2\times10^{-3}$ against
$1.5\times10^{-2}$ for $\mathrm{UB}$, both at $\alpha=45^\circ$
(scripts \texttt{b33\_lambda\_min\_sharp.py}, \texttt{b36\_cf\_chain.py}). The identity of
$\operatorname{cf}$ with the plane minimum was checked against brute-force one-dimensional
minimisation on all $6300$ draws, agreeing to $\le4.5\times10^{-12}$: the reduction is exact, not
fitted. Like the branch it replaces, \eqref{eq:cf} bounds $\lambda_{\min}(\Gr)$ from \emph{above}, so
it is an \emph{unsafe} direction for a certificate (Remark~\ref{rem:direction}); no downstream
estimate in this chapter uses $\operatorname{cf}$, which appears only as the comparison point of
Corollary~\ref{cor:square}. Its role is comparative: with Theorem~\ref{thm:exact} available, it says
how much of Theorem~\ref{thm:twosided}'s upper branch is an artifact of the restriction $z=0$ rather
than of the geometry.
\end{remark}

\begin{theorem}[exact scalar formula for $\lambda_{\min}(\Gr)$]\label{thm:exact}
Put $C:=\Le^{-1}H$ and, for $\lambda\in(0,\lambda_{\min}(\Le))$,
$B(\lambda):=C^{\top}\Le(\Le-\lambda I)^{-1}C\in\mathbb{R}^{q\times q}$. Then
$\lambda_{\min}(\Gr)=\min\bigl(\lambda_{\min}(\Le),\lambda^{\star}\bigr)$, where
$\lambda^{\star}$ is the unique root in $(0,\lambda_{\min}(\Le))$ of
\begin{equation}\label{eq:exact}
  \lambda\;+\;\cos^{2}\alpha\cdot\lambda\cdot\lambda_{\max}\!\bigl(B(\lambda)\bigr)
  \;=\;\sin^{2}\alpha\,,
\end{equation}
taken as $\lambda_{\min}(\Le)$ when no root exists in that interval. Thus
$\lambda_{\min}(\Gr)$ is a property of the dictionary, obtainable to \emph{absolute} round-off by a
monotone one-dimensional root find on a $q\times q$ eigenvalue, with no eigenvalue of the
$(p+q)$-dimensional Gram needed; Remark~\ref{rem:multi}(iii) quantifies what the word
\emph{round-off} can and cannot mean here.
\end{theorem}

\begin{proof}
In the coordinates $x=\Le^{-1/2}z$ the quotient \eqref{eq:rayleigh} is invariant under
$(x,w)\mapsto r(x,w)$, so the unit sphere splits into the face $w=0$, where it equals
$x^{\top}\Le x/\|x\|^{2}$ with minimum $\lambda_{\min}(\Le)$, and the slice $\|w\|=1$, on which we may
minimise over $x$ first. Fix unit $w$, write $y=Cw$ and $u=\cos\alpha\,y$, and let
$\mu(w):=\min_x(x^{\top}\Le x+\sin^2\alpha)/(\|x-u\|^2+1)$. A number $\lambda$ satisfies
$\lambda\le\mu(w)$ iff the form $x^{\top}\Le x+\sin^2\alpha-\lambda(\|x-u\|^2+1)$ is nonnegative for
all $x$; for $\lambda<\lambda_{\min}(\Le)$ that form is a strictly convex quadratic, minimised at
$x=-\lambda(\Le-\lambda I)^{-1}u$ with minimal value $-\Phi_w(\lambda)$, where, using
$I+\lambda(\Le-\lambda I)^{-1}=\Le(\Le-\lambda I)^{-1}$ to collapse the two $u$-terms,
\begin{equation}\label{eq:phiw}
  \Phi_w(\lambda)\;:=\;\lambda+\cos^2\alpha\cdot\lambda\cdot
  y^{\top}\Le(\Le-\lambda I)^{-1}y\;-\;\sin^2\alpha\,.
\end{equation}
So $\mu(w)$ is the largest $\lambda<\lambda_{\min}(\Le)$ with $\Phi_w(\lambda)\le0$. Because
$\Le(\Le-\lambda I)^{-1}=(I-\lambda\Le^{-1})^{-1}$ is Loewner-increasing in $\lambda$, each $\Phi_w$
is strictly increasing, and its root is unique; for $\lambda\ge\lambda_{\min}(\Le)$ the form is
unbounded below, so $\mu(w)\le\lambda_{\min}(\Le)$ and
$\lambda_{\min}(\Gr)=\min(\lambda_{\min}(\Le),\min_{\|w\|=1}\mu(w))$. It remains to perform the
$w$-minimisation. By \eqref{eq:phiw}, $\Phi_w$ depends on $w$ only through
$y^{\top}\Le(\Le-\lambda I)^{-1}y$, so
$\Psi(\lambda):=\max_{\|w\|=1}\Phi_w(\lambda)=\lambda+\cos^2\alpha\cdot\lambda\cdot
\lambda_{\max}(B(\lambda))-\sin^2\alpha$, with equality at a top eigenvector $w^{\star}(\lambda)$ of
$B(\lambda)$. $\Psi$ is strictly increasing by the same argument and $\Psi(0^+)=-\sin^2\alpha<0$.
If $\Psi$ has no root in the interval, set $\lambda^{\star}:=\lambda_{\min}(\Le)$ and every $\mu(w)$
equals $\lambda_{\min}(\Le)$, which is the claimed value; otherwise let $\lambda^{\star}$ be its root.
Then $\Phi_w(\lambda^{\star})\le\Psi(\lambda^{\star})=0$ gives $\mu(w)\ge\lambda^{\star}$ for every
unit $w$, while $w^{\star}:=w^{\star}(\lambda^{\star})$ satisfies
$\Phi_{w^{\star}}(\lambda^{\star})=\Psi(\lambda^{\star})=0$, i.e.\ $\mu(w^{\star})=\lambda^{\star}$.
Hence $\min_{\|w\|=1}\mu(w)=\lambda^{\star}$.
\end{proof}

\begin{corollary}[the square case: $\operatorname{cf}$ is an identity, not a bound]
\label{cor:square}
If $p=q$ and $C$ is invertible, then
$\lambda_{\min}(\Gr)=\operatorname{cf}(\|\Le^{-1}H\|_2^2,\alpha)$ exactly.
\end{corollary}

\begin{proof}
$C^{\top}\Le C=I_q$ with $C$ square and invertible gives $\Le=C^{-\top}C^{-1}$, so
$\Le^{-1}=CC^{\top}$, $\lambda_{\min}(\Le)=1/\sigma_{\max}^2(C)=1/t$, and with $S:=C^{\top}C$,
$B(\lambda)=C^{\top}(I-\lambda CC^{\top})^{-1}C=(I-\lambda S)^{-1}S=f(S)$, $f(x)=x/(1-\lambda x)$,
using the intertwining $C^{\top}g(CC^{\top})=g(C^{\top}C)C^{\top}$. On the admissible interval
$\lambda<1/t$ every eigenvalue of $S$ is $<1/\lambda$ and $f$ is increasing, so
$\lambda_{\max}(B(\lambda))=f(\lambda_{\max}(S))=t/(1-\lambda t)$, attained at $w_0$; and
\eqref{eq:exact} is $\lambda+\cos^2\alpha\cdot\lambda\,t/(1-\lambda t)=\sin^2\alpha$, i.e.\
after multiplying by $1-\lambda t$, exactly $t\lambda^2-(1+t)\lambda+\sin^2\alpha=0$. Its larger
root is at least $1/t$, because $\Delta:=(1+t)^2-4t\sin^2\alpha\ge(1-t)^2$ (as
$4t\cos^2\alpha\ge0$) gives $(1+t+\sqrt\Delta)/(2t)\ge(1+t+\lvert1-t\rvert)/(2t)\ge1/t$; since
$\lambda<\lambda_{\min}(\Le)=1/t$, the only admissible root is the smaller one,
$\operatorname{cf}(t,\alpha)$.
\end{proof}

\begin{remark}[$w_0$ is not optimal when $q<p$; what the closed form then costs]
\label{rem:w0}
The direction selected by Theorem~\ref{thm:exact} is a top eigenvector of $B(\lambda^{\star})$, not
of $C^{\top}C$; for $q<p$ the two differ, and $w_0$-optimality is \emph{false}. A dense search over
$S^1$ ($4000$ points) beats $\mu(w_0)$ in $157$ of the $840$ dictionaries tested in \texttt{b37}, and the
breakdown is the point: of the $180$ draws in the two non-square $q=2$ cells with
$\alpha\ge20^\circ$ --- where a circle search is \emph{exhaustive}, so the count is a measurement and
not a sampling rate --- $143$ admit a direction with $\mu(w)<\mu(w_0)$, i.e.\ the failure is the norm
($79\%$), not an exception, with angular separation from $w_0$ up to $0.30$ rad. The remaining $14$
come from random sampling in the $q=3$ cells and are therefore a lower bound on what a sphere search
would find. An earlier version of this remark reported that ``no random unit $w$ beats $w_0$'' in
$3200$ draws; that was an artifact of sampling --- $4000$ random points do
not resolve $S^2$, and the claim was only ever testable at $q=2$, where a circle search is exhaustive.
The same script contains the exception that identifies the true theorem: its two square cells,
$(d,p,q)=(5,2,2)$ and $(6,3,3)$, report $0$ violations and
$\mu(w_0)/\lambda_{\min}(\Gr)-1$ at the $10^{-13}$ level. Corollary~\ref{cor:square} is what explains
them --- those two rows had been logged as an unexplained anomaly in the round that produced them.
The cost is nevertheless small: the gap between $\mu(w_0)$ and the sphere minimum is
$\le9.6\cdot10^{-7}$ in $\lambda_{\min}(\Gr)$ units (the $\mu(w_0)$ restriction itself sits
$\le1.3\cdot10^{-5}$ above $\lambda_{\min}(\Gr)$), it grows
with $\alpha$ (the $w$-dependence beyond $\|Cw\|^2$ enters with weights $\lambda^j$), and
\eqref{eq:cf} stays a valid upper bound, short of the exact value by $\le1.9\times10^{-3}$ at
$(d,p,q)=(8,4,3)$ and $\le6.9\times10^{-3}$ at $(6,4,3)$, both at $\alpha=45^\circ$
(script \texttt{b39\_square\_case\_cf\_exact.py}, $200$ draws per cell, which also confirms
$\lambda_{\min}(\Gr)/\operatorname{cf}-1=0$ to $1.4\times10^{-11}$ on all seven square configurations
$\times$ five angles). Theorem~\ref{thm:exact} itself was checked against the exact Gram eigenvalue
on $6480$ dictionaries, nine $(d,p,q)$ configurations and $\alpha=1^\circ\ldots89^\circ$: worst
relative deviation $1.8\times10^{-11}$, i.e.\ round-off, no exception (script
\texttt{b38\_exact\_scalar\_formula.py}).
\end{remark}

\begin{proposition}[the equal-angle hypothesis is an artifact of the construction]\label{prop:multi}
Let $\Le\in\mathbb R^{p\times p}$ be positive definite, let $K\in\mathbb R^{p\times q}$, and let
$\Gr=\bigl(\begin{smallmatrix}\Le&K\\ K^{\top}&I_q\end{smallmatrix}\bigr)$ be positive semidefinite,
so that the singular values $c_1\ge\ldots\ge c_q\ge0$ of $\Le^{-1/2}K$ are principal cosines. On
$(0,\bar\lambda)$ with $\bar\lambda:=\min(\lambda_{\min}(\Le),1)$, put
\begin{equation}\label{eq:multi}
  h(\lambda)\;:=\;(1-\lambda)-\lambda_{\max}\!\bigl(K^{\top}(\Le-\lambda I)^{-1}K\bigr)\,.
\end{equation}
Then $h$ is
strictly decreasing and $\lambda_{\min}(\Gr)=\min(\bar\lambda,\lambda^{\star})$, where
$\lambda^{\star}$ is the root of $h$ if one lies in the interval, $\bar\lambda$ if $h>0$ throughout,
and $0$ if $h(0^{+})\le0$.
\end{proposition}

\begin{proof}
For $\lambda<\bar\lambda$ the block $\Le-\lambda I$ is positive definite, so $\Gr-\lambda I$ is
positive semidefinite iff its Schur complement in that block,
$(1-\lambda)I_q-K^{\top}(\Le-\lambda I)^{-1}K$, is. That complement is strictly decreasing in the
Loewner order on the interval: $\lambda\mapsto1-\lambda$ decreases and
$\lambda\mapsto(\Le-\lambda I)^{-1}$ increases. Hence its positivity set is an interval
$[0,\lambda^{\star}]$ and $\lambda_{\min}(\Gr)$, its supremum, is as claimed. At $\lambda=0$ the
complement is $I_q-K^{\top}\Le^{-1}K$, positive semidefinite because $\Gr$ is, and singular exactly
when $c_1=1$; that is the third case, the two subspaces meeting.
\end{proof}

\begin{corollary}[the isotropic locus: $\operatorname{cf}$ is exact with no angle hypothesis]
\label{cor:iso}
Let $\Le=\gamma I_p+Z$ with $\gamma>0$, $Z\in\mathbb R^{p\times p}$ symmetric, $ZK=0$, and let
$K\in\mathbb R^{p\times q}$ be arbitrary, with $\Gr$ positive semidefinite; write
$c^{2}:=\lambda_{\max}(K^{\top}\Le^{-1}K)\in[0,1]$. Then, with
$\bar\lambda:=\min(\lambda_{\min}(\Le),1)$,
\begin{equation}\label{eq:iso}
  \lambda_{\min}(\Gr)\;=\;\min\!\bigl(\bar\lambda,\;\lambda_-\bigr)\,,\qquad
  \lambda_-\;=\;\frac{(\gamma+1)-\sqrt{(\gamma-1)^{2}+4\gamma c^{2}}}{2}
  \;=\;\operatorname{cf}(\gamma^{-1},\alpha)\,,\quad \cos^{2}\alpha:=c^{2}\,.
\end{equation}
In particular $\lambda_{\min}(\Gr)=\lambda_-$ \emph{and $\operatorname{cf}$ is exact} whenever
$\lambda_{\min}(\Le)\ge\lambda_-$, which holds automatically for $Z\succeq0$; for $\gamma=1$ and
$Z\preceq0$ the condition reads $c\ge1-\lambda_{\min}(\Le)$. Under the chapter's normalisation
$\operatorname{diag}\Le=1$ isotropy forces $\gamma=1$, \emph{i.e.}\ orthonormal atoms, and
\eqref{eq:iso} reads $\lambda_{\min}(\Gr)=\min(\bar\lambda,1-c)$. The hypothesis $ZK=0$ is
non-vacuous exactly when $K$ has a non-trivial left kernel, \emph{i.e.}\ $p>q$ or a rank-deficient
dictionary: over-completeness is what makes a whole direction of $\Le$ invisible to
$\lambda_{\min}(\Gr)$.
\end{corollary}

\begin{proof}
$\Le-\lambda I=(\gamma-\lambda)I+Z$ leaves $\operatorname{ran}Z$ and its orthogonal complement
invariant, is $(\gamma-\lambda)I$ on that complement, and $K$ maps into it because $ZK=0$; hence
$(\Le-\lambda I)^{-1}K=(\gamma-\lambda)^{-1}K$ \emph{whether or not $Z$ is zero}. So
$\lambda_{\max}(K^{\top}(\Le-\lambda I)^{-1}K)=\gamma c^{2}/(\gamma-\lambda)$ on $(0,\bar\lambda)$ and
\eqref{eq:multi} becomes $h(\lambda)=(1-\lambda)-\gamma c^{2}/(\gamma-\lambda)$, whose zero set is
the quadratic $\lambda^{2}-(\gamma+1)\lambda+\gamma(1-c^{2})=0$. Its smaller root is the $\lambda_-$
displayed, and $\lambda_-\le\gamma$ and $\lambda_-\le1$ are each equivalent to
$|\gamma-1|\le\sqrt{(\gamma-1)^{2}+4\gamma c^{2}}$, so $\lambda_-\le\min(\gamma,1)$; for $Z=0$ that is
$\lambda_-\le\bar\lambda$ and Proposition~\ref{prop:multi} gives $\lambda_{\min}(\Gr)=\lambda_-$
(the quadratic is non-negative at $0$ and equals $-\gamma c^{2}\le0$ at $\bar\lambda$; at $c=0$ it
returns $\bar\lambda$, which is $\lambda_{\min}$ of $\gamma I_p\oplus I_q$). With $Z\ne0$ the same
argument runs on $(0,\bar\lambda)$ and gives the announced $\min(\bar\lambda,\lambda_-)$, the cap
being taken at $\lambda_{\min}(\Le)=\min(\gamma,\gamma+\lambda_{\min}(Z))$ rather than at $\gamma$.
At $\gamma=1$ with $Z\preceq0$ the criterion $\lambda_-\le\lambda_{\min}(\Le)=1-|\lambda_{\min}(Z)|$
is $h(\lambda_{\min}(\Le))\le0$, \emph{i.e.}\ $c^{2}\ge|\lambda_{\min}(Z)|^{2}$, the stated
$c\ge1-\lambda_{\min}(\Le)$.
On the other hand $\Le^{-1}K=\gamma^{-1}K$ again by $ZK=0$, so
$t=\|\Le^{-1}K\|_{2}^{2}/c^{2}=\gamma^{-1}$ whatever $Z$, and clearing the factor $1/\gamma$ in
$\operatorname{cf}(t,\alpha)$ of \eqref{eq:cf} turns its discriminant into
$(\gamma+1)^{2}-4\gamma(1-c^{2})=(\gamma-1)^{2}+4\gamma c^{2}$, which is the second equality.
Script \texttt{b44} checks \eqref{eq:iso} on $18$ cells ($\le1.7\cdot10^{-15}$ absolute): $K$ a plain
Gaussian block rescaled to prescribed top cosine $0.3/0.9/0.999$, full and rank-$1$, $(p,q)$ equal to
$(4,3)$, $(3,3)$ and $(3,5)$, and $\gamma\in\{1,1.5,0.6\}$. Script \texttt{b45} checks the $Z\ne0$
clause: $\Le=I_4-\delta ww^{\top}$ with $w\in\ker K^{\top}$ ($p=4>q=3$), $\delta=10^{-3},10^{-6}$,
one spread and one equal-angle spectrum, $400$ draws --- $|t-1|\le1.3\cdot10^{-15}$, median
$|\lambda_{\min}(\Gr)-\lambda_-|\le3.0\cdot10^{-16}$ and $\max|D|\le1.3\cdot10^{-15}$ absolute.
\end{proof}

\begin{remark}[the two exactness loci are two fibres of one mechanism]\label{rem:loci}
Corollary~\ref{cor:square} ($p=q$, isometry) and Corollary~\ref{cor:iso} ($\Le$ scalar, any $K$) look
alike because in both cases the pencil $\lambda\mapsto B(\lambda)=C^{\top}\Le(\Le-\lambda I)^{-1}C$ is
an \emph{increasing scalar function applied to one fixed Hermitian matrix} ---
$B(\lambda)=(I-\lambda C^{\top}C)^{-1}C^{\top}C$ in the square case, $B(\lambda)=\bigl(\gamma/(\gamma-\lambda)\bigr)C^{\top}C$
here --- so its top eigenvalue is determined by one number and \eqref{eq:exact} collapses to a
quadratic. Neither hypothesis is the other: the square-isometric locus has non-scalar $\Le$
($\Le=C^{-\top}C^{-1}$), and the isotropic locus needs no isometry at all. What they share is the
conclusion, and it is the honest reading of $\operatorname{cf}$: the scalar reduction is exact when
$\Le$ carries no directional information that $t$ could lose. At scalar $\Le$ the resolvent is a
scalar multiple of the identity, so \eqref{eq:multi} decouples into $q$ quadratics, one per principal
cosine --- the classical arrowhead case rather than a new formula. What it does for this chapter is
name the hypothesis that $\operatorname{cf}$ was really using.
\end{remark}

\begin{proposition}[the residual is a moment defect of one scalar measure, and its sign is decidable]
\label{prop:side}
Write $\Le=\sum_{i=1}^{p}\ell_i u_iu_i^{\top}$, $k_i:=K^{\top}u_i$, and for a unit $v\in\mathbb R^{q}$
let $\nu_v:=\sum_i(k_i\cdot v)^2\delta_{\ell_i}$, so that
$\int x^{-j}d\nu_v=v^{\top}K^{\top}\Le^{-j}Kv$ for every $j\ge1$. Let $v$ maximise
$K^{\top}\Le^{-1}K$, so that $\int x^{-1}d\nu_v=c^{2}$, and on $[0,\bar\lambda)$ put
\[
  f(\lambda):=\lambda_{\max}\!\bigl(K^{\top}(\Le-\lambda I)^{-1}K\bigr),\qquad
  g(\lambda):=\frac{c^{2}}{1-t\lambda}\,,\qquad
  D(\lambda):=f(\lambda)-g(\lambda)\,.
\]
Then $g$ is $\int d\nu/(x-\lambda)$ for the single atom $\hat\nu:=(c^{2}/t)\,\delta_{1/t}$, which
matches $\nu_v$ in $\int x^{-1}$ and over-states it in $\int x^{-2}$, and
\emph{(i)} in the uncapped case $\lambda_{\min}(\Gr)<\bar\lambda$, with
$\operatorname{cf}\in(0,\bar\lambda)$,
\begin{equation}\label{eq:side}
  \begin{gathered}
  D(\operatorname{cf})\;=\;\bigl(\operatorname{cf}-\lambda_{\min}(\Gr)\bigr)
  +\bigl(f(\operatorname{cf})-f(\lambda_{\min}(\Gr))\bigr),\\[2pt]
  \operatorname{sgn}D(\operatorname{cf})=\operatorname{sgn}\bigl(\operatorname{cf}-\lambda_{\min}(\Gr)\bigr),
  \qquad
  \bigl|\operatorname{cf}-\lambda_{\min}(\Gr)\bigr|\le\bigl|D(\operatorname{cf})\bigr|;
  \end{gathered}
\end{equation}
if $f$ is differentiable between the two points, $\operatorname{cf}-\lambda_{\min}(\Gr)=
D(\operatorname{cf})/(1+f'(\xi))$ for some $\xi$. \emph{(ii)} As $\lambda\downarrow0$,
\begin{equation}\label{eq:slack}
  D(\lambda)\;=\;-\lambda\,\operatorname{slack}+o(\lambda)\,,\qquad
  \operatorname{slack}:=c^{2}t-\lambda_{\max}(E_0^{\top}K^{\top}\Le^{-2}KE_0)\;\ge\;0\,,
\end{equation}
where $E_0$ is the top eigenspace of $K^{\top}\Le^{-1}K$: so the leading term of $D$ is always on the
\emph{flip} side ($\operatorname{cf}<\lambda_{\min}(\Gr)$), it vanishes exactly when some maximiser
of $K^{\top}\Le^{-1}K$ also maximises $K^{\top}\Le^{-2}K$, and where it vanishes the sign is decided
by the remainder. \emph{(iii)} Three named loci have $\operatorname{slack}=0$: $q=1$ (scalars),
equal angles ($K^{\top}\Le^{-1}K=c^{2}I_q$, so $E_0=\mathbb R^{q}$), and
Corollary~\ref{cor:iso} (there $D\equiv0$, not merely to first order).
\end{proposition}

\begin{proof}
$\hat\nu$ gives $\int d\hat\nu/(x-\lambda)=(c^{2}/t)/(1/t-\lambda)=c^{2}/(1-t\lambda)=g(\lambda)$,
$\int x^{-1}d\hat\nu=c^{2}=\int x^{-1}d\nu_v$, and
$\int x^{-2}d\hat\nu=c^{2}t=\|\Le^{-1}K\|_2^{2}\ge v^{\top}K^{\top}\Le^{-2}Kv=\int x^{-2}d\nu_v$.
Both $f$ and $g$ are increasing on $[0,\bar\lambda)$ ($\Le-\lambda I$ decreases in the Loewner order,
and $t>0$). At $\lambda_{\min}(\Gr)$, which is the root $\lambda^{\star}$ of \eqref{eq:multi} in the
uncapped case, $(1-\lambda_{\min})=f(\lambda_{\min})$; at $\operatorname{cf}$, clearing $q$ in
\eqref{eq:cf} gives $(1-\operatorname{cf})=g(\operatorname{cf})$, and $\operatorname{cf}<1/t$ is
strict because $q(1/t)=-c^{2}\le0$ while $\operatorname{cf}$ is the smaller root, so
$\operatorname{cf}$ is inside the interval where $g$ is finite. Subtracting the two scalar identities
yields the first display of (i); its second factor has the sign of
$\operatorname{cf}-\lambda_{\min}(\Gr)$ because $f$ is increasing, giving the sign equality and, since
the two summands then share a sign, the bound $|\operatorname{cf}-\lambda_{\min}(\Gr)|\le|D|$. The
mean-value reading is the same identity with
$f(\operatorname{cf})-f(\lambda_{\min})=f'(\xi)(\operatorname{cf}-\lambda_{\min})$ and $1+f'(\xi)\ge1$.
For (ii), $g(\lambda)=c^{2}+\lambda c^{2}t+O(\lambda^{2})$, while the right derivative at $0$ of
$\lambda\mapsto\lambda_{\max}(A+\lambda B)$ at a multiple top eigenvalue of $A$ is
$\lambda_{\max}(E_0^{\top}BE_0)$ (Lidskii--Danskin), here with
$A=K^{\top}\Le^{-1}K$, $B=K^{\top}\Le^{-2}K$; the displayed $\ge0$ is
$\lambda_{\max}(E_0^{\top}BE_0)\le\lambda_{\max}(B)=c^{2}t$, with equality iff a maximiser of $A$
attains the top of $B$. Hence $D(\lambda)=-\lambda\operatorname{slack}+o(\lambda)$, and if
$\operatorname{slack}>0$ then $D<0$ on some $(0,\delta)$: a $\operatorname{cf}$ landing there sits
below $\lambda_{\min}(\Gr)$. (iii) is immediate: $q=1$ and equal angles make $E_0^{\top}BE_0=B$ a
scalar or give $E_0=\mathbb R^{q}$, and at scalar $\Le$ the resolvent identity used in
Corollary~\ref{cor:iso} gives $f=g$ on the whole interval.
\end{proof}

\begin{remark}[what the side-bound test measured, and what it did not]\label{rem:side}
Proposition~\ref{prop:side} exists because the size of the residual had been attributed by \texttt{b44}
but its \emph{direction} unexplained. Script \texttt{b45} tests (i)--(iii) on ${\sim}6650$ draws,
$\operatorname{eigvalsh}$-based floors included, and the pre-registered reading is the measured one:
(i) gives $\operatorname{sgn}D(\operatorname{cf})\ne\operatorname{sgn}(\operatorname{cf}-\lambda_{\min})$
in $0$ of $500$ resolved draws, and the magnitude reading is a \emph{quantitative} oracle, not merely a
sign oracle --- $\operatorname{gap}\big/\bigl(D(\operatorname{cf})/(1+f'(\lambda_{\min}))\bigr)$ has
median deviation $\le1.1\cdot10^{-4}$ from $1$ in every cell (best $2.6\cdot10^{-8}$). (ii) holds where
it was claimed: $\operatorname{slack}\ge-1.8\cdot10^{-15}$ in all cells, and in every flipping cell
$|\operatorname{gap}/(\operatorname{cf}\cdot\operatorname{slack})|=0.92$--$1.02$, while the two families
that do \emph{not} flip sit at $13$ and $126$--$394$: the linear term is the whole story exactly when it
wins, which is the content of the $o(\lambda)$ and not of a fitted law. At $\Le=\gamma I_p$ the
diagnostic collapses as (iii) demands ($\max|D|\le1.3\cdot10^{-15}$ over $15$ cells,
$\gamma\in\{1,1.5,0.6\}$). What the script does \emph{not} buy: a criterion for flipping in terms of
the cosine spectrum alone --- see retraction~(m) of Section~\ref{sec:notclaims} --- and any statement about
$\lambda_{\min}(\Gr)$ at capped $\operatorname{cf}\ge\bar\lambda$, which the draws record as skipped
rather than read ($43$ of the $500$ $q=1$ draws, all at the smallest cosine $\cos\alpha=0.2$).
\end{remark}

\begin{proposition}[the chord surrogate: an upper bound with no angle hypothesis]
\label{prop:chord}
Keep the notation of Proposition~\ref{prop:side} ($\Le\succ0$, $K\in\mathbb R^{p\times q}$ with
$\Gr$ positive semidefinite, $c^{2}:=\lambda_{\max}(K^{\top}\Le^{-1}K)\in(0,1]$, $E_0$ the top
eigenspace of $K^{\top}\Le^{-1}K$) and put
\begin{equation}\label{eq:chord}
  t_1\;:=\;\frac{\lambda_{\max}(E_0^{\top}K^{\top}\Le^{-2}KE_0)}{c^{2}}\;\le\;t\,,\qquad
  \operatorname{cf}_1\;:=\;\operatorname{cf}(t_1,\alpha),\qquad\cos^{2}\alpha:=c^{2}\,.
\end{equation}
Then, with no hypothesis on the angles, on $\Le$, and no case split at $\bar\lambda$,
\begin{equation}\label{eq:chain1}
  \lambda_{\min}(\Gr)\;\le\;\operatorname{cf}_1\,,\qquad
  \operatorname{cf}\;\le\;\operatorname{cf}_1\,,
\end{equation}
the second being equality exactly where $\operatorname{slack}=0$ in \eqref{eq:slack}. More generally,
for every unit $v\in\mathbb R^{q}$ with $B(v):=v^{\top}K^{\top}\Le^{-1}Kv$ and
$C(v):=v^{\top}K^{\top}\Le^{-2}Kv$, the smaller root $\lambda(v)$ of
\begin{equation}\label{eq:family}
  C(v)\,\lambda^{2}-\bigl(B(v)+C(v)\bigr)\lambda+B(v)\bigl(1-B(v)\bigr)\;=\;0
\end{equation}
satisfies $\lambda_{\min}(\Gr)\le\lambda(v)$, and $\operatorname{cf}_1$ is the member of that family
with the largest intercept ($B=c^{2}$, which forces $v\in E_0$) and, among those, the largest slope.
\end{proposition}

\begin{proof}
Fix a unit $v$ with $B(v)>0$, write $a_i:=(k_i\cdot v)^{2}$, and let
$\rho_v:=\sum_i(a_i/\ell_i)\,\delta_{1/\ell_i}$, a positive measure of mass $B(v)$ and first moment
$C(v)$. For $\lambda\in[0,\bar\lambda)$ every atom has $\lambda/\ell_i<1$ and
$y\mapsto(1-\lambda y)^{-1}$ is convex there, so Jensen's inequality
$\int h\,d\rho_v\ge\rho_v(\mathbb R)\,h(\int y\,d\rho_v/\rho_v(\mathbb R))$ in the variable $y=1/x$
turns the measure $\nu_v$ of Proposition~\ref{prop:side} into the \emph{chord}
\[
  \phi_v(\lambda):=\sum_i\frac{a_i}{\ell_i-\lambda}\;\ge\;\frac{B(v)}{1-\lambda C(v)/B(v)}
  \;=:\;s_v(\lambda)\,,\qquad\text{hence}\qquad f\;\ge\;s_v\ \text{on }[0,\bar\lambda)
\]
for \emph{every} $v$, because $f(\lambda)=\max_{\|v\|=1}\phi_v(\lambda)$. Let
$F(\lambda):=(1-\lambda)(B(v)-\lambda C(v))-B(v)^{2}$, so that $s_v=1-\lambda$ reads $F=0$, i.e.\
\eqref{eq:family}: $F(0)=B(v)(1-B(v))\ge0$, $F$ opens upward in $\lambda$, and
$F(B(v)/C(v))=-B(v)^{2}<0$ where $s_v$ blows up, so the smaller root $\lambda(v)$ is the first
contact and obeys $\lambda(v)\le B(v)/C(v)$, where $s_v$ is finite and increasing. If
$\lambda(v)<\lambda_{\min}(\Le)$, then
$f(\lambda(v))\ge s_v(\lambda(v))=1-\lambda(v)$, and as
$\lambda\mapsto f(\lambda)-(1-\lambda)$ is strictly increasing on $[0,\bar\lambda)$ with unique zero
$\lambda^{\star}$ (Proposition~\ref{prop:multi}), this gives $\lambda(v)\ge\lambda^{\star}\ge
\min(\bar\lambda,\lambda^{\star})=\lambda_{\min}(\Gr)$. If instead
$\lambda(v)\ge\lambda_{\min}(\Le)$, then
$\lambda(v)\ge\lambda_{\min}(\Le)\ge\bar\lambda\ge\lambda_{\min}(\Gr)$, the last step being the second
clause of Proposition~\ref{prop:multi}. That is \eqref{eq:chain1} for the family. Take now $v_1\in E_0$
maximising $C$ over the unit sphere of $E_0$: $K^{\top}\Le^{-1}K=c^{2}I$ on $E_0$ gives $B(v_1)=c^{2}$,
and $C(v_1)=c^{2}t_1$ by the definition of $t_1$, so $s_{v_1}(\lambda)=c^{2}/(1-t_1\lambda)$ and
$F(\lambda)=0$ is exactly $(1-\lambda)(1-t_1\lambda)=c^{2}$, i.e.\ $Q(\lambda,t_1)=0$ of \eqref{eq:cf},
whose smaller root is $\operatorname{cf}_1$. Compressing cannot raise a maximum, so $t_1\le t$;
differentiating $Q(\lambda,t)=0$ at the smaller root gives
$\partial\lambda/\partial t=\lambda(1-\lambda)/(2t\lambda-1-t)<0$ because $2t\lambda<1+t$ there, so
$\operatorname{cf}$ decreases in $t$, whence $\operatorname{cf}\le\operatorname{cf}_1$ with equality
iff $t_1=t$, which is $\operatorname{slack}=0$.
\end{proof}

\begin{remark}[what the chord-surrogate tests measured]\label{rem:chord}
$\operatorname{cf}_1$ was derived \emph{before} it was run, as the surrogate that replaces $g$ by the
chord through the two moments it should have used; the pre-registered tests were four, and
a fifth test (\texttt{b48}) searches rather than samples.
\emph{One:} the flips found by the two-cosine scan (\texttt{b45}) must go. On the same cells $\operatorname{cf}$ lands below
$\lambda_{\min}(\Gr)$ in $100/100$, $100/100$ and $98/100$ draws while $\operatorname{cf}_1$ does so in
$0/100$ in every cell, uncapped and capped alike ($1600$ further draws in the capped rescan \texttt{b47}: $0$ resolved
sign flips, $\min\operatorname{gap}_1=1.1\cdot10^{-11}$ over the thirteen cells whose gap clears the
floor; the three isotropy cells report $0/0$ because there the gap is identically zero). \emph{Two:} the correction should
cancel the first-order term and leave the remainder. The ratio
$\operatorname{med}|\operatorname{gap}_1|/\operatorname{med}|\operatorname{gap}|$ is $0.001$--$0.31$
in every cell where $\operatorname{slack}$ dominates, and $1.00$--$1.10$ in the ten where the slope
inflation is absent or negligible --- the five $q=1$ controls, the two equal-angle cells, and the three
cells (the $30$-$88$-$89$ family of that scan) that run left unexplained; among them the cell of
retraction~(m), which $\operatorname{cf}_1$ leaves untouched, as it must: the correction removes only
what the slope inflation put there. Along that scan's top-cosine ladder the log-log slope of
$|\operatorname{gap}_1|$ against $\operatorname{cf}_1$ is $2.081$, the predicted value being $2$ and
the falsifier $1$. \emph{Three:} the $q=1$ and isotropy nulls must be unchanged, for $t_1=t$ there:
$\max|\operatorname{cf}_1-\operatorname{cf}|=1.3\cdot10^{-15}$ over five $q=1$ cells, and at
$\Le=\gamma I_p$, $\gamma\in\{1,1.5,0.6\}$, $|t_1-1/\gamma|\le2.0\cdot10^{-15}$,
$|\operatorname{cf}_1-\operatorname{cf}|\le2.5\cdot10^{-16}$, $\max|D_1|\le1.3\cdot10^{-15}$ --- which
independently re-verifies that the Lidskii--Danskin slope of Proposition~\ref{prop:side}(ii) really is
the right derivative at $0$. \emph{Four:} the inequality itself, attacked directly: the Jensen residual
$\phi_{v_1}-s_{v_1}$ is $\ge0$ at $\lambda=0.3,0.7,0.95$ of $\operatorname{cf}_1$ in every cell (worst
$-5.6\cdot10^{-16}$, in the isotropy cells where the gap is identically $0$), and the whole family
$\lambda(v)$ of \eqref{eq:family} stayed above $\lambda_{\min}(\Gr)$ for all $200$ random directions of
all $1600$ draws --- $320{,}000$ tests, $0$ violations. That scan also seemed to show nothing beats
$\operatorname{cf}_1$, and the polar-net search (\texttt{b48}, $96\times240$ directions) instead
found that it does: the first variation of $\lambda(v)$ at $v_1$ is
$-2\,v_1^{\top}K^{\top}\Le^{-2}Ku\cdot\lambda(1-\lambda)/|\partial_\lambda F|$ along
$v(t)=\cos t\,v_1+\sin t\,u$, $u\perp v_1$, generally nonzero because $t_1$ maximises $C$ over $E_0$
alone --- so $\operatorname{cf}_1$ is a saddle of the family, stationary iff
$K^{\top}\Le^{-2}Kv_1\parallel v_1$, a condition that holds exactly on the $\operatorname{slack}=0$
loci ($\dim E_0=q$; there the net finds $0$ competitors) and fails generically ($84/84$ draws). The
propagation is limited: the best ray removes $0.07$--$0.50$ of $\operatorname{cf}_1$'s over-estimate at
the median ($0.58$ worst), every competitor buys slope with intercept
($C/(c^{2}t_1)=1.000005\ldots1.000499$ against $B/c^{2}\le1-10^{-6}$), and no family member drops below
$\lambda_{\min}(\Gr)$ on the net or on $144{,}000$ ray points, so the proposition stands
(retraction~(n) records the two artifacts of that run, a mis-chosen $v_1$ and a normalised noise
direction). The closed-form sizing of that descent (\texttt{b49}) needs no optimiser. On the steepest
ray, $u\parallel P_{v_1^{\perp}}K^{\top}\Le^{-2}Kv_1$, the two jets are
$B(t)=c^{2}+b_{2}t^{2}+O(t^{4})$ --- no linear term, since $B$ is an exact trigonometric quadratic on a
great circle when $A_0v_1=c^{2}v_1$ --- and $C(t)=C(v_1)+c_{1}t+c_{2}t^{2}+O(t^{3})$ with
$c_{1}=2\|P_{v_1^{\perp}}K^{\top}\Le^{-2}Kv_1\|$, and implicit differentiation of \eqref{eq:family} at
the smaller root gives $\lambda(t)=\operatorname{cf}_1+L_1t+L_2t^{2}$,
$L_1=-F_Cc_1/\partial_\lambda F<0$ and
$L_2=-(\partial_{\lambda\lambda}F\,L_1^{2}+2\partial_{\lambda C}F\,L_1c_1+2F_Bb_2+2F_Cc_2)/(2\partial_\lambda F)$,
where all coefficients of $F(\lambda;B,C)=C\lambda^{2}-(B+C)\lambda+B-B^{2}$ are evaluated at
$(\operatorname{cf}_1,c^{2},C(v_1))$: $F_C=\lambda(\lambda-1)$, $F_B=1-\lambda-2B$,
$\partial_\lambda F=2C\lambda-B-C<0$, $\partial_{\lambda\lambda}F=2C$, $\partial_{\lambda C}F=2\lambda-1$,
so the ray has an interior minimum at $t^{*}=-L_1/(2L_2)$ of depth
$\operatorname{dip}=-L_1^{2}/(4L_2)$:
\emph{quadratic} in the off-$E_0$ gradient although the descent itself is linear in $t$. Measured over
$280$ saddle draws, $\operatorname{dip}_{\mathrm{meas}}/\operatorname{dip}=1$ to within
$6\cdot10^{-6}\ldots1.0\cdot10^{-3}$ per cell and $t_{\mathrm{meas}}/t^{*}\in[0.993,1.008]$, with
$L_2>0$ in $280/280$; the single ray step removes $0.0004$--$0.50$ of $\operatorname{cf}_1$'s
over-estimate at the median ($0.76$ worst), reproducing the net's range on a different seed by
a different method. \emph{$\operatorname{cf}_1+\operatorname{dip}$ is not offered as a bound} ---
one-sidedness would need a signed third jet, and the measured drift with $|c_1|$ has no consistent sign,
which is exactly the audit I ask of other people's closed forms. What is a bound is the family member
itself, $\operatorname{cf}_2:=\lambda(\cos t^{*}\,v_1+\sin t^{*}\,u)$: Proposition~\ref{prop:chord}
applied to that explicit unit vector gives $\lambda_{\min}(\Gr)\le\operatorname{cf}_2$ with no new proof,
at the price of one $2\times2$ root, and the two-jet supplies only where to step and how much is there.
A $64\times64$ net on the same draws gets at most $\sim10\%$ deeper than the ray (median
$\operatorname{ray}/\operatorname{net}=0.906$--$1.009$, and that net is not a superset of the ray at
$5.6^{\circ}$ resolution, so the statement is ``sampled, none found'', not exhaustive); the absolute size
is $|\operatorname{dip}|\le1.2\cdot10^{-7}$, so this leg closes here --- no closed form for the family
minimum, and no marginal left worth chasing. Two caveats are recorded in the statement rather than hidden: the proposition uses
$\lambda_{\max}$ \emph{over} $E_0$, not ``any top eigenvector'', because at a multiple top eigenvalue
the vector an eigensolver returns need not maximise $C$ --- the equal-angle cells have $\dim E_0=3$ and
$|C(v_1)-c^{2}t_1|$ up to $6.6\cdot10^{-1}$ --- and \eqref{eq:chain1} is an upper bound, the
\emph{unsafe} side for a certificate (Remark~\ref{rem:direction}), which no downstream estimate in this
chapter uses. What changes is the status of the closed form: an upper estimate of
$\lambda_{\min}(\Gr)$ now exists with no equal-angle, no isometry and no simplicity hypothesis, at the
price of one extra $q\times q$ compression; and the flip of Proposition~\ref{prop:side} is explained as
the one unlicensed step in \eqref{eq:cf} --- pairing the intercept $c^{2}$ with a slope measured in a
direction where that intercept was never attained.
\end{remark}

\begin{remark}[what Proposition~\ref{prop:multi} gives, and what it takes away]\label{rem:multi}
Put $K=\cos\alpha\,H$, $C:=\Le^{-1}H$, and assume the isometry $C^{\top}\Le C=I_q$ of
Theorem~\ref{thm:exact}. Then $\Le(\Le-\lambda I)^{-1}=I+\lambda(\Le-\lambda I)^{-1}$ gives
$K^{\top}(\Le-\lambda I)^{-1}K=\cos^{2}\alpha\,(I_q+\lambda B(\lambda))$, so $h(\lambda)=0$ reads
\eqref{eq:exact}: Theorem~\ref{thm:exact} is the \emph{equal-angle fibre} of
Proposition~\ref{prop:multi}, and the cost statement transfers verbatim --- a monotone one-root find
on a $q\times q$ eigenvalue, no eigenvalue of the $(p+q)$-dimensional Gram. Three consequences, the
first two of them against my own interest.
\emph{(i) The general statement is not new, and \eqref{eq:multi} is not offered as a contribution.}
$h(\lambda)=0$ is the secular equation of an arrowhead (bordered diagonal) matrix: the eigenvalues of
a diagonal-plus-rank-one matrix are the roots of one monotone scalar equation, each computable
separately, which is the basis of the rank-one modification and divide-and-conquer solvers of
numerical linear algebra \cite{bunch1978,cuppen1980,gu-eisenstat1995} and of the shift-and-invert
arrowhead solvers \cite{stor2013}, whose rank-one secular equation is
the case $q=1$, $K$ a single column, of \eqref{eq:multi}. Proposition~\ref{prop:multi} is the
rank-$q$ reading of that device, with the scalar equation replaced by the top eigenvalue of a
$q\times q$ pencil; we state it because the chapter's own theorem is the $q\ge2$ case of it and not
because the reduction is new. What this note owns is the identification: the dictionary factor of
$(\star)$ \emph{is} an arrowhead secular equation. \emph{(The classical records above are
verified bibliographically --- DOI, authors, title, venue --- and the connection was drawn at
abstract level, not against those papers' theorems; if the rank-$q$ case is already spelled out
somewhere in that literature, Proposition~\ref{prop:multi} should be replaced by a citation to it.)}
\emph{(ii) Relaxing the angle buys almost nothing for $\lambda_{\min}$, and the residual
$\operatorname{cf}$ leaves behind is not an angle effect at all.} The prescribed-spectrum run \texttt{b42} prescribes
cosine spectra in place of the constructed single angle ($p=4$, $q=3$, $q=2$ cells, $80$--$120$ draws
per cell).
Proposition~\ref{prop:multi} reproduces $\lambda_{\min}(\Gr)$ to $\le1.6\cdot10^{-12}$ relative over
spread spectra, and to $\le1.3\cdot10^{-15}$ against the independent closed answer
$\lambda_{\min}(\Gr)=1-c_1$ available at $\Le=I_p$; on the equal-angle cells it agrees with
\eqref{eq:exact} to $\le4.8\cdot10^{-13}$, so it is a generalisation and not a correction. Yet
$\operatorname{cf}(t,\alpha)$, given the most generous angle available (the top cosine) and the
rescaled $t$, is still within $3.4\%$ of $\lambda_{\min}(\Gr)$ --- and the $2\times5$
design (\texttt{b44}) settles \emph{which} hypothesis that residual is paying for. It is not the angle distribution,
as a first reading of these numbers said (retraction~(l) of Section~\ref{sec:notclaims}): with the
cosine spectrum fixed, replacing a generic Gram by $\Le=\gamma I_p$ cuts the median relative gap by
eight to twelve orders of magnitude ($1.3\cdot10^{-4}\!\to\!5.5\cdot10^{-15}$ at $20^\circ$,
$1.8\cdot10^{-3}\!\to\!4.0\cdot10^{-13}$ at $2$--$45$--$88^\circ$), exactly the collapse
Corollary~\ref{cor:iso} asserts; with $\Le$ fixed, replacing the equal angle by a spread spectrum
raises it by only thirteen to three hundred. The anisotropy of $\Le$ is the variable, and
the two-cosine scan says of what: by Proposition~\ref{prop:side} the residual is a \emph{moment defect}, the
gap between the measure $\nu_v$ that $f$ integrates and the single atom $\hat\nu=(c^{2}/t)\delta_{1/t}$
that $\operatorname{cf}$ integrates. The correct reading is therefore not ``the angles are spread'' but
``$\nu_v$ carries more than one atom, and the leading term \eqref{eq:slack} beats the remainder'' ---
in which a second cosine is neither necessary nor sufficient. I wrote that it was necessary
(retraction~(m) of Section~\ref{sec:notclaims}): with the geometry held multi-atom, the same scan
measures the flip in $100/100$ draws at a second cosine of $0.087$ and again in $100/100$ at
$\operatorname{cf}=0.133$, while the \emph{same} spectra placed on a single-atom geometry leave every
gap on the $\operatorname{eigvalsh}$ floor ($12$ cells, $\operatorname{slack}=0$ to round-off). So
That design's tally is a statement about size, not about which hypothesis licenses a flip:
$\operatorname{cf}-\lambda_{\min}$ is $\ge0$ on all $240$ equal-angle draws and $\le0$ on $100/100$
($5$--$40$--$75^\circ$) and $99/100$ ($2$--$45$--$88^\circ$) of the spread draws, whereas
$30$--$88$--$89^\circ$, spread but with secondary cosines $\approx0.03$, does not flip ($0/100$) ---
there too $\operatorname{slack}>0$, but $\operatorname{cf}\cdot\operatorname{slack}$ is smaller than
the gap by a factor $126$--$394$, which is the remainder winning, exactly as \eqref{eq:slack}'s $o(\lambda)$
allows. The anisotropy design \texttt{b44} also closes the audit
the prescribed-spectrum run left open: the smallest violating gap in the two flipping cells is $9.1\cdot10^{5}$ and
$1.1\cdot10^{4}$ times the per-draw $\operatorname{eigvalsh}$ floor $100\epsilon\|\Gr\|_2$, with no
draw below its own floor, so the flip is algebraic and not round-off --- while at $\Le=I_p$, where the
true gap is identically zero, the same tally comes out as a coin flip ($35$--$50$ of $100$), which is
what a null should produce.
The upper-side reading of \eqref{eq:cf} is an equal-angle statement too --- which is consistent with
Remark~\ref{rem:cf}: no estimate in this chapter uses $\operatorname{cf}$, and a future one that
called it a bound for a learned dictionary would be overstepping the hypotheses of
Proposition~\ref{prop:plane}.
\emph{(iii) The root find is accurate in absolute terms, and the conditioning scan measures where that stops
being enough.} Over $60$ draws per cell at $p=q=3$ and $\alpha=45^\circ$ down to $0.001^\circ$, the
relative deviation of the bisection from $\operatorname{eigvalsh}(\Gr)$ obeys
$\mathrm{relerr}\times\lambda_{\min}(\Gr)\in[1.6,1.8]\cdot10^{-15}$ with no trend across six decades of
$\lambda_{\min}(\Gr)$: the method pins $\lambda_{\min}(\Gr)$ to a few units of the pencil's absolute
round-off, so its \emph{relative} accuracy inherits the very $1/\sin^{2}\alpha$ amplification this
chapter proves for the error ($\lambda_{\min}(\Gr)\approx\sin^{2}\alpha/(1+\cos^{2}\alpha\,t)$).
Two operational consequences. The bracket endpoint must sit a distance $\gg\epsilon\|\Le\|$ below the
pole $\lambda_{\min}(\Le)$: a relative gap of $10^{-12}$ makes the resolvent return garbage of order
$1$, and the ``no root, report the cap'' branch then fires silently --- that, not \eqref{eq:multi}, was
the source of the order-$100$ errors of the first version of that scan, and the guarded endpoint restores
the absolute floor down to $\lambda_{\min}(\Gr)\approx10^{-9}$. Below
$\lambda_{\min}(\Gr)\approx\epsilon\|\Le\|$ no tuning of the bracket helps, which is exactly the regime
the high-relative-accuracy arrowhead solvers of \cite{stor2013} are built for. None of this
touches $(\star)$, which uses the closed-form lower side $\mathrm{LB}$ and not $\lambda^{\star}$
(Remark~\ref{rem:exactcost}); what it constrains is any future use of $\lambda^{\star}$ as a
\emph{reported number} at small $\alpha$.
\end{remark}

\begin{proposition}[the chord surrogate is the first member of a certified moment hierarchy]
\label{prop:gauss}
Keep the notation of Proposition~\ref{prop:chord} and put, for a unit $v$ and $j\ge0$,
$m_j:=v^{\top}K^{\top}\Le^{-(j+1)}Kv=\int y^{j}\,d\rho_v(y)$ (so $m_0=B(v)$, $m_1=C(v)$).
Run the Stieltjes procedure on $\rho_v$: $h_0:=m_0$, $\pi_0:=1$, $\pi_{j+1}:=(y-a_j)\pi_j-b_j\pi_{j-1}$
with $a_j:=\int y\pi_j^{2}\,d\rho_v/h_j$ and $b_j:=h_j/h_{j-1}$, $h_j:=\int\pi_j^{2}\,d\rho_v$, and let
$J_n(v)$ be the resulting symmetric tridiagonal matrix with diagonal $a_0,\ldots,a_{n-1}$ and
off-diagonals $\sqrt{b_1},\ldots,\sqrt{b_{n-1}}$. Assume $h_0,\ldots,h_{n-1}>0$, \emph{i.e.}\
$\rho_v$ carries at least $n$ atoms, and let $x_1,\ldots,x_n$ be the eigenvalues of $J_n(v)$ with
weights $w_j:=m_0\,(q_j)_1^{2}$ ($(q_j)_1$ the first component of the $j$-th normalised eigenvector).
Then $G_n^v(\lambda):=\sum_j w_j/(1-\lambda x_j)$ is the $n$-point Gauss rule for
$y\mapsto(1-\lambda y)^{-1}$, it is finite and strictly increasing on $[0,\bar\lambda)$, satisfies
$G_n^v(0)=B(v)$ and $G_n^v(\lambda)\to\infty$ as $\lambda\uparrow\max_j 1/x_j\ (\ge\bar\lambda)$, and
its unique crossing $\lambda_n(v)$ with $1-\lambda$ obeys
\begin{equation}\label{eq:hier}
  \lambda_{\min}(\Gr)\;\le\;\lambda_n(v)\qquad
  \text{for every unit }v\text{ and every such }n\ge1,
\end{equation}
and $\lambda_1(v)$ is the smaller root of \eqref{eq:family}, so that
$\operatorname{cf}_1=\lambda_1(v_1)$ is the member $n=1$.
Explicitly, $\lambda_2(v)$ is the smallest non-negative root of
\begin{equation}\label{eq:cubic}
  p_2\lambda^{3}-(p_2+s_1)\lambda^{2}+\bigl[(s_1+1)-(m_0s_1-m_1)\bigr]\lambda+(m_0-1)=0,
\end{equation}
where $a_0=m_1/m_0$, $b_1=m_2/m_0-a_0^{2}$, $a_1=(m_3-2a_0m_2+a_0^{2}m_1)/(m_0b_1)$,
$s_1=a_0+a_1$, $p_2=a_0a_1-b_1$, and $G_2^v(\lambda)=[m_0-(m_0s_1-m_1)\lambda]/(1-s_1\lambda+p_2\lambda^{2})$.
\end{proposition}

\begin{proof}
For $\lambda\in[0,\bar\lambda)$ every atom has $\lambda y<1$ and
$h_\lambda(y):=(1-\lambda y)^{-1}\in C^{\infty}$ of the support of $\rho_v$, with all even derivatives
$h_\lambda^{(2n)}(y)=(2n)!\,\lambda^{2n}(1-\lambda y)^{-2n-1}\ge0$. The classical Gauss error formula
$\int h_\lambda\,d\rho_v-G_n^v(\lambda)=h_\lambda^{(2n)}(\xi)/(2n)!\cdot h_n\ge0$ is the
Gauss--Markov sign statement for rules whose integrand has $2n$ nonnegative even derivative; it is
recorded in this form in the standard treatments of Gauss quadrature and orthogonal polynomials
\cite{gautschi1996,gautschi2004,szego1939}, and we cite it as a classical theorem rather than
reproducing it with a section number, because we have verified those records bibliographically but
not read the displayed formula in them. Nothing in what follows rests on that gap: the consequence
actually used --- $G_n^v\le\phi_v$ --- is checked instead by the moment identities
$\sum_jw_jx_j^k=m_k$, $k<2n$, on every draw (Remark~\ref{rem:gauss}), so the inequality is verified
inside the construction that uses it
$G_n^v\le\phi_v$ on $[0,\bar\lambda)$. The nodes are in the convex hull of the support,
$x_j\le1/\lambda_{\min}(\Le)$, hence $\max_j1/x_j\ge\bar\lambda$ and $G_n^v$ is finite and strictly
increasing there; $G_n^v(0)=\sum_jw_j=m_0=B(v)\le1$ and $G_n^v\to\infty$ at the pole give a unique
crossing $\lambda_n(v)$ with $1-\lambda$, and $\lambda_n(v)\ge\bar\lambda$ already implies
$\lambda_n(v)\ge\bar\lambda\ge\lambda_{\min}(\Gr)$ by the second inequality of
Proposition~\ref{prop:side}. Otherwise $\lambda_n(v)\in[0,\bar\lambda)$ and
$f(\lambda_n(v))\ge\phi_v(\lambda_n(v))\ge G_n^v(\lambda_n(v))=1-\lambda_n(v)$, which as in
Proposition~\ref{prop:chord} places $\lambda_n(v)\ge\lambda^{\star}\ge\lambda_{\min}(\Gr)$ --- the only
step replaced is $s_v\le\phi_v$ by $G_n^v\le\phi_v$, so the case split is the same one.
For $n=1$: $J_1=[m_1/m_0]$ and $w_1=m_0$, so
$G_1^v(\lambda)=m_0/(1-\lambda m_1/m_0)=B(v)^2/(B(v)-\lambda C(v))=s_v$, whose crossing is
\eqref{eq:family}; taking $v=v_1\in E_0$ with $B=\cos^2\alpha$ gives $\operatorname{cf}_1$.
For $n=2$: $\sum_jw_j=m_0$ and $\sum_jw_jx_j=m_0\,e_1^{\top}J_2e_1=m_0a_0$ are exact (degree $\le3$
is integrated exactly), $\sum_jw_jx_j^{2}=m_0(a_0^{2}+b_1)=m_2$ identifies $b_1$, and
$\operatorname{tr}J_2=a_0+a_1=s_1$, $\det J_2=a_0a_1-b_1=p_2$ give $\{x_j\}$; the numerator of $G_2^v$
is $m_0-\lambda(w_1x_2+w_2x_1)$ with $w_1x_2+w_2x_1=m_0s_1-m_1$ from
$(w_1+w_2)(x_1+x_2)=\sum w_jx_j+\sum w_jx_{j'}$. Clearing the denominator, which is
$(1-\lambda x_1)(1-\lambda x_2)=1-s_1\lambda+p_2\lambda^2>0$ at the crossing, gives \eqref{eq:cubic}.
\end{proof}

\begin{proposition}[the fibre bound is exact at the minimal eigenvector, and that is all it is]
\label{prop:tight}
Write $\Gr[x;y]=\lambda[x;y]$ in the blocks $x\in\mathbb R^{p}$, $y\in\mathbb R^{q}$, and let
$\lambda:=\lambda_{\min}(\Gr)$. If $\lambda<\lambda_{\min}(\Le)$ --- which holds whenever
$\lambda_{\min}(\Le)>1$, since $\lambda_{\min}(\Gr)\le\bar\lambda\le1$ by Proposition~\ref{prop:side}
--- then $y\ne0$ and
\begin{equation}\label{eq:fibre}
  K^{\top}(\Le-\lambda I)^{-1}K\,y\;=\;(1-\lambda)\,y\,,
  \qquad\text{so for }v_\ast:=y/\|y\|:\quad
  \phi_{v_\ast}(\lambda)=1-\lambda\,.
\end{equation}
The crossing $\lambda_{\rm ex}(v_\ast)$ of $\phi_{v_\ast}$ with $1-\lambda$ therefore obeys
$\lambda_{\rm ex}(v_\ast)=\lambda_{\min}(\Gr)$: the family $\{\lambda(v)\}$ of \eqref{eq:family}, and
the hierarchy of Proposition~\ref{prop:gauss} once $n$ reaches the atom count of $\rho_{v_\ast}$,
attain the certificate. Nothing here makes them \emph{computable}: $v_\ast$ is the $q$-block of the
eigenvector whose eigenvalue is being bounded, and the member $n=1$ does not attain it --- measured on
$320$ draws, $(\lambda_1(v_\ast)-\lambda_{\min}(\Gr))/(\operatorname{cf}_1-\lambda_{\min}(\Gr))$ has
median $0.50$--$1.01$ in all eight cells (Remark~\ref{rem:tight}).
\end{proposition}

\begin{proof}
The two block rows of $\Gr[x;y]=\lambda[x;y]$ read $(\Le-\lambda I)x=-Ky$ and $K^{\top}x=(\lambda-1)y$.
If $y=0$ the first gives $\Le x=\lambda x$ with $x\ne0$, contradicting $\lambda<\lambda_{\min}(\Le)$, so
$y\neq0$; and $\Le-\lambda I\succ0$, so $x=-(\Le-\lambda I)^{-1}Ky$ and the second row becomes
\eqref{eq:fibre}. Divide by $\|y\|^{2}$ and use $a_i=(k_i\cdot v_\ast)^{2}$: the left side of
\eqref{eq:fibre} is $\sum_ia_i/(\ell_i-\lambda)=\phi_{v_\ast}(\lambda)$, the right side is
$1-\lambda$. On $[0,\lambda_{\min}(\Le))$ the map $\mu\mapsto\phi_{v_\ast}(\mu)-(1-\mu)$ is strictly
increasing, is $\le0$ at $\mu=0$ (where it equals $B(v_\ast)-1$) and diverges at the pole
$\mu=\ell_{\min}$, so it has exactly one zero there; $\lambda$ is a zero, hence
$\lambda=\lambda_{\rm ex}(v_\ast)$. The displayed inequality is Proposition~\ref{prop:gauss} at
$n=\#\operatorname{supp}\rho_{v_\ast}$, where $G_n^{v_\ast}=\phi_{v_\ast}$ exactly.
\end{proof}

\begin{lemma}[ordering two members of the family: the margin]\label{lem:margin}
For a unit $v$ with $B:=B(v)\in(0,1)$ put $t:=C(v)/B(v)$, and let $r(B,t)$ be the smaller root of
$(1-x)(1-tx)=B$, so that $\lambda_1(v)=r(B(v),t(v))$ is the $n=1$ member of \eqref{eq:family}
and $\operatorname{cf}_1=r(c^{2},t_1)$. Then on the smaller-root branch
\begin{equation}\label{eq:partials}
  \frac{\partial r}{\partial t}=\frac{r(1-r)}{2tr-1-t}\;<\;0\,,\qquad
  \frac{\partial r}{\partial B}=\frac{1}{2tr-1-t}\;<\;0\,,
\end{equation}
and for every $\mu\in\bigl[0,(1+t)/(2t)\bigr]$ the comparison of a family member with a candidate value
is an \emph{exact} criterion,
\begin{equation}\label{eq:margin}
  r(B,t)\;\le\;\mu\quad\Longleftrightarrow\quad(1-\mu)(1-t\mu)\;\le\;B
  \quad\Longleftrightarrow\quad
  t\;\ge\;t_1+\frac{c^{2}-B}{c^{2}}\cdot\frac{1-t_1\mu}{\mu}\,,\qquad\mu=\operatorname{cf}_1 .
\end{equation}
Moreover the two knobs collapse to \emph{one} at first order: since at $\mu=\operatorname{cf}_1$ one has
$\partial r/\partial t=\mu(1-\mu)\,\partial r/\partial B$ and $(1-\mu)(1-t_1\mu)=c^{2}$, the intercept
term equals $\delta:=c^{2}-B$ times $\partial r/\partial B$ exactly, and
\begin{equation}\label{eq:firstorder}
  \lambda_1(v)-\operatorname{cf}_1\;=\;\frac{\mu(1-\mu)}{2t_1\mu-1-t_1}
  \Bigl(t-t_1-\frac{\delta}{c^{2}}\frac{1-t_1\mu}{\mu}\Bigr)\bigl(1+o(1)\bigr):
\end{equation}
the sign of the bracket --- the \emph{margin} --- decides which of the two members is the better bound.
\end{lemma}

\begin{proof}
Both partials are implicit differentiation of $Q(x;t,B):=tx^{2}-(1+t)x+1-B$ along its smaller root:
$Q_x=2tx-1-t$, and the smaller root lies \emph{left} of the vertex $(1+t)/(2t)$, so $Q_x<0$ there;
$Q_t=x^{2}-x=-x(1-x)$ and $Q_B=-1$ give \eqref{eq:partials} with $0<r<1$. Since $Q$ opens upward with
$Q(0)=1-B>0$, and is strictly decreasing on $[0,(1+t)/(2t)]$, for $\mu$ in that interval $r\le\mu$
reads $Q(\mu)\le0$, i.e.\ $(1-\mu)(1-t\mu)\le B$. Taking $\mu=\operatorname{cf}_1$, dividing by
$1-\mu>0$ and substituting $(1-\mu)(1-t_1\mu)=c^{2}$ --- the root equation at $(c^{2},t_1)$ --- turns
that into $1-t\mu\le(1-\delta/c^{2})(1-t_1\mu)$, which is the right hand side of \eqref{eq:margin};
every step is reversible, so the only side conditions are $\mu\le(1+t)/(2t)$, $\mu<1$ and
$1-t_1\mu>0$, and the ray scan \texttt{b55} verifies each on every draw. Finally
$\frac{\partial r}{\partial t}\cdot\frac{\delta}{c^{2}}\frac{1-t_1\mu}{\mu}
=\delta\,\frac{\partial r}{\partial B}$ by the two displayed identities, which is the collapse in
\eqref{eq:firstorder}.
\end{proof}

\begin{corollary}[the direction question is closed when the angles are equal]\label{cor:equal}
If $M:=\Le^{-1/2}K$ satisfies $M^{\top}M=c^{2}I_q$ --- an equal cosine spectrum does, since then
$M=Q_1\operatorname{diag}(c,\ldots,c)Q_2^{\top}$ --- then $B(v)\equiv c^{2}$ on the unit
sphere, $E_0=\mathbb R^{q}$, and
\[
  \min_{\|v\|=1}\lambda_1(v)\;=\;\operatorname{cf}_1\,,
\]
attained exactly at the unit maximisers of $C(v)$; every other direction, $v_\ast$ included, gives a
strictly larger member. So in those cells the certificate is \emph{not} attained by any $n=1$ member,
however the direction is chosen --- and $\delta=0$ makes \eqref{eq:margin} read
$\lambda_1(v)\le\operatorname{cf}_1\iff t(v)\ge t_1$, which no direction satisfies.
\end{corollary}

\begin{proof}
$B(v)=v^{\top}M^{\top}Mv=c^{2}\|v\|^{2}$, so $\delta=0$ and every unit vector saturates the intercept,
which is $E_0=\mathbb R^{q}$; then $t_1=\max_{\|v\|=1}C(v)/c^{2}$ by \eqref{eq:chord}, and
$\lambda_1(v)=r(c^{2},C(v)/c^{2})$, which \eqref{eq:partials} minimises exactly where $C$ is maximal,
with value $r(c^{2},t_1)=\operatorname{cf}_1$. Strictness is the strict negativity of
$\partial r/\partial t$.
\end{proof}

\begin{lemma}[the descent set on a ray is a cubic]\label{lem:cubic}
Fix a unit $v_1\in E_0$ maximising $C$ over $E_0$, so $B(v_1)=c^{2}$, $t(v_1)=t_1$ and
$\lambda_1(v_1)=\operatorname{cf}_1=:\mu$, and let $u$ be any unit vector with $u\perp v_1$. Put
\[
  a:=B(u),\qquad b:=C(u),\qquad g:=v_1^{\top}K^{\top}\Le^{-2}Ku,\qquad
  D:=c^{2}-a,\qquad P:=\frac{Dc^{2}}{\mu(1-\mu)}-b+t_1a ,
\]
in which the coefficient of the cubic below uses only $\mu$:
$\sigma:=\frac{1-t_1\mu}{c^{2}\mu}=\frac1{\mu(1-\mu)}$, since $(1-\mu)(1-t_1\mu)=c^{2}$.
For $\theta\in(0,\pi/2)$ and $v_\theta:=\cos\theta\,v_1+\sin\theta\,u$,
\[
  B(v_\theta)=c^{2}-D\sin^{2}\theta,\qquad
  C(v_\theta)-t_1B(v_\theta)=2g\sin\theta\cos\theta+(b-t_1a)\sin^{2}\theta,
\]
and with $\Phi(x):=(1+x^{2})(Px-2g)-\sigma D^{2}x^{3}$,
\begin{equation}\label{eq:raycubic}
  \lambda_1(v_\theta)\;\le\;\operatorname{cf}_1\qquad\Longleftrightarrow\qquad
  \Phi(\tan\theta)\;\le\;0 .
\end{equation}
So the boundary of the descent set on the ray is a positive root of a cubic in $\tan\theta$ with four
coefficients built from five scalars; leaving $v_1$ descends at all iff $g>0$, since $\Phi(0)=-2g$;
dropping $\sigma D^{2}$ makes the boundary exactly $\tan\theta=2g/P$, which is $t_{\rm hi}$ of
\texttt{b52} read analytically; and if $P>0$ that two-jet value pertains to the true root as
$x_{\rm hi}=x\bigl(1+\frac{\sigma D^{2}x^{2}}{P(1+x^{2})}+O((\sigma D^{2})^{2})\bigr)$, $x=2g/P$.
\end{lemma}

\begin{proof}
$K^{\top}\Le^{-1}Kv_1=c^{2}v_1$ gives $v_1^{\top}K^{\top}\Le^{-1}Ku=c^{2}(v_1\cdot u)=0$ for every
$u\perp v_1$, so the cross term in $B(v_\theta)$ vanishes and $B=v_1$-term $\cos^{2}\theta+a\sin^{2}
\theta=c^{2}-D\sin^{2}\theta$; the same expansion with $\Le^{-2}$ keeps the cross term, which is $g$,
and subtracting $t_1B$ cancels the $c^{2}\cos^{2}\theta$ against $t_1c^{2}\cos^{2}\theta$ to leave the
displayed $C-t_1B$. Now \eqref{eq:margin} is $\lambda_1(v_\theta)\le\mu\iff T\ge0$ for
$T:=t-t_1-\sigma(c^{2}-B)$, and $B>0$ lets us multiply it by $B$:
$0\le C-t_1B-\sigma B(c^{2}-B)=2g\,sc+bs^{2}-t_1as^{2}-\sigma s^{2}(c^{2}-Ds^{2})
=s^{2}\bigl[2g\,c/s-P+\sigma D^{2}s^{2}\bigr]$ with $s:=\sin\theta$, $c:=\cos\theta$. Dividing by
$s^{2}>0$ and multiplying by $x(1+x^{2})>0$ at $x=\tan\theta$ (so $c/s=1/x$, $s^{2}=x^{2}/(1+x^{2})$)
turns the bracket into $-\Phi(x)$, which is \eqref{eq:raycubic}. Setting $\sigma D^{2}=0$ in $\Phi$ leaves
$(1+x^{2})(Px-2g)$, whose only real root is $2g/P$; and at that root the $\sigma D^{2}$-free part of
$\Phi$ has derivative $P(1+x^{2})$, while $\Phi$ itself changes by $-\sigma D^{2}x^{3}$, so the root
shifts by $\sigma D^{2}x^{3}/(P(1+x^{2}))$ to first order in $\sigma D^{2}$ --- the stated
perturbation, and it is $O(x^{3})$ in the layer width, which is why the two-jet value is the right one
to quote when $x$ is small.
\end{proof}

\begin{remark}[What the direction leg did not buy]
\label{rem:dirleg}
The direction leg was run to look for a cheap way to close the residual gap
$\operatorname{cf}_1-\lambda_{\min}(\Gr)$ by optimising the direction $v$ alone, and it bought
nothing: the $n=1$ chord loses $0.50$--$1.01$ of that gap \emph{even at} $v_\ast$, the direction
that attains the exact fibre crossing of Proposition~\ref{prop:tight}, whereas one additional
scalar moment removes $1.00$ of it on the same tables (median \emph{and} max, $20/20$ cells).
The residual of $\operatorname{cf}_1$ is therefore truncational, not directional, and the
angle $\theta=\angle(v_\ast,v_b)$ that a reader might otherwise take for a convergence error is
a model-discrepancy angle: $v_\ast$ is not a critical point of $\lambda_1$ at all, so no
stationary-point story explains $\operatorname{cf}_1$. The measurement is recorded in full in
Appendix~\ref{app:direction} --- direction nets and their resolution limits, the two independent
routes to the descent-set boundary, the cubic-share attribution of the residual, and the moment
hierarchy --- because a negative result is only as credible as the resolutions at which it was
tested.
\end{remark}

\begin{remark}[what the exact formula does for the certificate]
\label{rem:exactcost}
$\lambda_{\min}(\Gr)$ is now available as an \emph{equality}, so the dictionary factor $q_4$ of
\S\ref{sec:cancellation} measures the price of the closed form that $(\star)$ uses, not an
intrinsic limitation: replacing $\mathrm{LB}$ in $(\star)$ by $\lambda^{\star}$ of
\eqref{eq:exact} removes the $5\%$ at $\alpha=2^\circ$ entirely and leaves only the sampling
factors. We keep the closed form, because it needs no root find, no eigenvalue, and no extra
calibration --- and because $q_4$ is the quantity that says what that choice costs.
\end{remark}

\section{The sampling side: a remainder that does not know about $\alpha$}
\label{sec:sampling}

The certificate needs $\minr G$, but only $\Gr$ is a property of the dictionary. The
bridge is
$\mathbb{E}[G^{\top}G]=n\Gr$ \emph{exactly} (the expectation identity is analytic; the
finite-sample version was verified numerically (script \texttt{b11}),
so all that can go wrong is the concentration of
$G^{\top}G/n$ around $\Gr$ in the small directions.

\begin{lemma}[per-direction quadratic form]\label{lem:quad}
For a fixed unit coefficient vector $x$ put $M_x=\sum_a x_aM_a$ and
$r(M):=\|MM^{\top}\|_F/\|M\|_F^2\le1$. Then
\begin{equation}\label{eq:quad}
  \frac{\|M_xU\|_F^2}{n\|M_x\|_F^2}=1+O\!\Bigl(r(M_x)\sqrt{\tfrac1n}\Bigr)
  \quad\text{with}\quad
  \operatorname{Var}\bigl(\|M_xu_i\|^2\bigr)=2\|M_xM_x^{\top}\|_F^2\le2\|M_x\|_F^4 .
\end{equation}
\end{lemma}

\begin{proof}
$\|M_xU\|_F^2=\sum_{i=1}^{n}\|M_xu_i\|^2$, and for standard Gaussian $u$,
$\mathbb{E}\|M_xu\|^2=\|M_x\|_F^2$ while $\|M_xu\|^2$ is a Wishart quadratic form with
variance $2\|M_xM_x^{\top}\|_F^2$. Averaging $n$ independent copies and dividing by
$n\|M_x\|_F^2$ gives the relative fluctuation
$2^{1/2}r(M_x)/\sqrt n$. Covering the unit sphere of the $(p+q)$-dimensional
coefficient space by a net converts the fixed-direction statement into a uniform one at
the price of a factor $\sqrt{p+q}$.
\end{proof}

\begin{theorem}[angle-free sampling remainder]\label{thm:sampling}
With the notation above, uniformly over coefficient directions,
\begin{equation}\label{eq:anglefree}
  \Bigl|\frac{\|M_xU\|_F^2}{n\|M_x\|_F^2}-1\Bigr|
  \le \kappa\, r(M_x)\sqrt{\tfrac{p+q}{n}},
  \qquad\text{hence}\qquad
  \minr G=\sqrt n\,\minr{\Gr^{1/2}}\Bigl(1\pm\kappa\sqrt{\tfrac{p+q}{n}}\Bigr),
\end{equation}
and the right-hand side contains \textbf{no} $\alpha$ and \textbf{no} $d$. The
statement is valid when $\kappa\sqrt{(p+q)/n}\ll1$, i.e.\ $n\gg p+q$, independently of
$\alpha\to0$.
\end{theorem}

\begin{remark}[what this buys, and what it costs]
Theorem~\ref{thm:sampling} is the hinge of the whole chapter: it says the $1/\sin\alpha$
law survives going from the exact dictionary to finite samples. An earlier attempt at
this proof used an operator-norm perturbation bound and produced $O(d/(\sqrt n\sin^2\alpha))$,
which at $\alpha=0.3^\circ$, $d=24$, $n=200$ predicts a relative deviation of $6.2\!\times\!10^{4}$
against a measured $0.0125$ --- wrong by seven orders of magnitude, and structurally
so, because a smallest-eigenvalue perturbation is quadratic, not linear. Singular-value
lower bounds must not be obtained from global operator-norm perturbations.
\end{remark}

\begin{remark}[the displayed rate is provable but loose]\label{rem:loose}
The net argument in Lemma~\ref{lem:quad} gives $\sqrt{p+q}$; the measured fluctuation does not
scale that way. Theorem~\ref{thm:bias} below replaces it by the dimension of the
near-null subspace, $D_{\rm eff}=O(1)$. We state both because the first is what we can
prove and the second is what the data say; publishing only the first would misreport the
size of the remainder, publishing only the second would overclaim a theorem.
\end{remark}

\noindent\textbf{Numerical status of Theorem~\ref{thm:sampling}.}
The mechanism ($r$ flat in $\alpha$, $\sim1.4/\sqrt d$):
\begin{center}\begin{tabular}{lccc}
\toprule & $\alpha=1^\circ$ & $\alpha=30^\circ$ & isotropic $1/\sqrt d$\\
\midrule
$d=10$ & $0.4515$ & $0.4500$ & $0.316$\\
$d=24$ & $0.2910$ & $0.2921$ & $0.204$\\
\bottomrule
\end{tabular}\end{center}
The rate ($\alpha$ across $90\times$ at fixed $n$ shows no trend; $n:100\to3200$ decays
$\approx6\times$ against the $\sqrt n$ prediction $5.7\times$), with $|\varepsilon|$
averaged over 12 trials at $d=10$, $p+q=6$:
\begin{center}\begin{tabular}{lcccc}
\toprule $n$ & $\alpha=1^\circ$ & $5^\circ$ & $30^\circ$ & $90^\circ$\\
\midrule
$100$  & $0.0263$ & $0.0347$ & $0.0405$ & $0.0346$\\
$200$  & $0.0137$ & $0.0171$ & $0.0133$ & $0.0194$\\
$800$  & $0.0113$ & $0.0111$ & $0.0097$ & $0.0111$\\
$3200$ & $0.0059$ & $0.0040$ & $0.0034$ & $0.0042$\\
\bottomrule
\end{tabular}\end{center}
$|\varepsilon|/\sqrt{(p+q)/n}\in[0.07,0.17]\Rightarrow\kappa\approx0.15$. Worst case over
$40$ draws: relative downward deviation of $\lambda_{\min}$ is $12.7\%$/$6.3\%$/$2.7\%$
at $n=200/600/1800$, always below $\sqrt{(p+q)/n}=17.3\%/10.0\%/5.8\%$.
(Scripts: \texttt{b11}, \texttt{b14\_g\_alpha\_kappa.py}.)

\subsection{The remainder refined: a computable one-sided bias, not a symmetric fluctuation}
\label{sec:bias}

\begin{theorem}[shape of the remainder]\label{thm:bias}
Write $H:=G^{\top}G/n=\Gr+D$ and let $(\lambda_1,v_1)$ be the smallest eigenpair of
$\Gr$. Define the second-order displacement
\begin{equation}\label{eq:biasdef}
  \beta_2:=\sum_{k\ge2}\frac{\mathbb{E}\,|\langle v_1,Dv_k\rangle|^2}{\lambda_k-\lambda_1}
  \;=\;\frac1{n^2}\sum_{k\ge2}\frac{\mathbb{E}\,|\langle v_1,\Delta v_k\rangle|^2}
        {\lambda_k-\lambda_1},
  \qquad \Delta:=G^{\top}G-n\Gr .
\end{equation}
Then, with $\varepsilon_n:=\minr G/\sqrt{n\lambda_{\min}(\Gr)}-1$,
\begin{equation}\label{eq:biasform}
  \varepsilon_n=-\frac{\beta_2}{2\lambda_1}\ \pm\ c_{\rm fluc}\,r\sqrt{\frac{D_{\rm eff}}{n}},
  \qquad
  \mathbb{E}\bigl[\lambda_{\min}(G^{\top}G)\bigr]-n\lambda_1=-n\beta_2+O(n^{-1}).
\end{equation}
The leading term is (i) one-sided and negative ---
$\mathbb{E}[\lambda_{\min}(G^{\top}G)/n]<\lambda_{\min}(\Gr)$, by concavity of
$\lambda_{\min}$ together with $\mathbb{E}H=\Gr$; (ii) of order $1/n$ in $\beta_2$ for
fixed $\Gr$ (each $\langle v_1,Dv_k\rangle$ is $O_p(n^{-1/2})$ by Lemma~\ref{lem:quad}),
hence \emph{relative} to $\lambda_1$ it is $\beta_2/(2\lambda_1)=O(1/n)$, not
$O(n^{-1/2})$; and (iii) computable from the dictionary alone, since
$\beta_2$ needs only the eigendecomposition of $\Gr$ and the second moments of $U$.
Note the two readings of the same quantity: the \emph{absolute} displacement
$n\beta_2$ is $O(1)$ in $n$, while the \emph{relative} one is $O(1/n)$. The fluctuation
term carries the near-null dimension $D_{\rm eff}$, not $p+q$.
\end{theorem}

\begin{proof}[Proof outline]
Second-order Rayleigh--Schr\"odinger expansion of $\lambda_{\min}$ at $\Gr$ under the
symmetric perturbation $D$: the first-order shift $\langle v_1,Dv_1\rangle$ has mean
$0$ because $\mathbb{E}D=0$, so the mean displacement is the sum of the second-order
terms, each of which is negative for the smallest eigenvalue since $\lambda_k-\lambda_1>0$
for $k\ge2$; the displayed value of $n\beta_2$ is that sum, and the remainder requires the
non-degeneracy condition $|\langle v_1,Dv_k\rangle|\ll\lambda_k-\lambda_1$, which is the
regime where the expansion is valid. Concavity of $\lambda_{\min}$ makes the average
strictly below the deterministic value: this is Jensen's gap in operator form, and it is
the same mechanism (a concave spectral functional, a zero-mean perturbation) that produces
the one-sided residual in the coverage certificate of the companion chapter. The symmetric
spread left over is Lemma~\ref{lem:quad} restricted to the set of directions that nearly
attain $\lambda_1$, whose dimension is $D_{\rm eff}$.
\emph{Not} proved here and not claimed: an upper bound on $\beta_2/\lambda_1$ that is
uniform in $\alpha$. Since $\lambda_1\asymp\sin^2\alpha$ (Theorem~\ref{thm:det}),
uniformity requires $\beta_2=O(\sin^2\alpha/n)$, i.e.\ the numerator must vanish at
exactly the same rate as the denominator; Theorem~\ref{thm:gammaclosed} records this as a
measured fact over a scanned parameter range, with the parameter grid stated there.
\end{proof}

\begin{center}\begin{tabular}{lcccccc}
\toprule $n$ & 100 & 200 & 600 & 1800 & 5400\\
\midrule
$1-\mathbb{E}[\lambda_{\min}(G^{\top}G)/n]/\lambda_{\min}(\Gr)$
 & $0.1342$ & $0.0840$ & $0.0344$ & $0.0102$ & $0.0052$\\
$\times n$ & $13.4$ & $16.8$ & $20.6$ & $18.4$ & $27.8$\\
\bottomrule
\end{tabular}\end{center}

\noindent $\times n$ is flat within a logarithmic drift, while $\times\sqrt n$ would have
to decrease by a factor $5$ over the table and does not ($1.34,1.19,0.84,0.44,0.38$):
the rate is $1/n$. Quantiles of $\varepsilon_n$ ($p+q=24$, $d=10$, $\alpha=5^\circ$,
$400$ trials) confirm one-sidedness --- at $n=600$ the $95\%$ quantile is $+0.0042$, so
$\minr G$ falls short of the deterministic prediction with probability $\approx0.95$:
\begin{center}\begin{tabular}{lrrrrrrrr}
\toprule $n$ & $q_{0.0025}$ & $q_{0.01}$ & $q_{0.05}$ & median & $q_{0.95}$ & $q_{0.99}$ & mean & sd\\
\midrule
$200$  & $-0.0945$ & $-0.0885$ & $-0.0766$ & $-0.0413$ & $-0.0111$ & $+0.0030$ & $-0.0432$ & $0.0201$\\
$600$  & $-0.0461$ & $-0.0436$ & $-0.0355$ & $-0.0149$ & $+0.0042$ & $+0.0135$ & $-0.0155$ & $0.0123$\\
$1800$ & $-0.0245$ & $-0.0208$ & $-0.0165$ & $-0.0058$ & $+0.0056$ & $+0.0105$ & $-0.0055$ & $0.0071$\\
\bottomrule
\end{tabular}\end{center}

\noindent And the predicted-versus-measured displacement $-n\beta_2$, with \emph{no} fitted
parameter:
\begin{center}\begin{tabular}{lccc}
\toprule $n$ & measured mean shift $\mathbb{E}[\lambda_{\min}(G^{\top}G)]-n\lambda_1$ & predicted $-n\beta_2$ & ratio\\
\midrule
$200$  & $-1.06\!\times\!10^{-1}$ & $-1.28\!\times\!10^{-1}$ & $0.825$\\
$600$  & $-1.20\!\times\!10^{-1}$ & $-1.38\!\times\!10^{-1}$ & $0.870$\\
$1800$ & $-1.35\!\times\!10^{-1}$ & $-1.43\!\times\!10^{-1}$ & $0.947$\\
\bottomrule
\end{tabular}\end{center}

\noindent The absolute shift is nearly $n$-independent, exactly as $n\beta_2=O(1)$ predicts;
agreement improves with $n$ because the perturbation-series validity condition
$|\langle v_1,\Delta v_k\rangle|\ll\lambda_k-\lambda_1$ degrades at small $n$, which is
also the source of the residual $5$--$18\%$ gap.
(Scripts: \texttt{b15\_bias\_not\_noise.py}, \texttt{b24\_ggamma.py},
\texttt{b25\_pqn\_law.py}.)

\subsection{Uniformity in $\alpha$ and the size of the constants}\label{sec:gammaclosed}

The question that decides whether the certificate needs an applicability domain is
whether the relative bias decays more slowly than $\lambda_{\min}(\Gr)\asymp\sin^2\alpha$
as $\alpha\to0$. It does not.

\begin{theorem}[$\alpha$-uniform bias; measured]\label{thm:gammaclosed}
For $d\in\{3,\ldots,24\}$, $p+q\in\{2,\ldots,14\}$, $n\in\{200,600,1200,2400\}$ and
$\alpha\in(0,30^\circ]$,
\begin{equation}\label{eq:c9}
  \Bigl|\frac{\mathbb{E}[\lambda_{\min}(G^{\top}G/n)]}{\lambda_{\min}(\Gr)}-1\Bigr|
  \le \frac{c}{n},\qquad c\le 9,
\end{equation}
while for $\alpha$ up to $89^\circ$, $c\le20$. The fluctuation satisfies
$\operatorname{sd}\bigl(\lambda_{\min}(G^{\top}G/n)/\lambda_{\min}(\Gr)\bigr)
=0.68/\sqrt n$ in $21/21$ cells, independent of $p+q$.
Consequently $\minr G\ge\sqrt{n\lambda_{\min}(\Gr)}\,(1-c/2n-0.88/\sqrt n)$ with
$0.5\%$ failure probability.
\end{theorem}

\begin{center}\begin{tabular}{lccc}
\toprule $n$ & mean relative bias, $\alpha\le20^\circ$ & theory $-(p+q)/n=-7/n$ & $\alpha=85^\circ$\\
\midrule
$200$  & $-0.0356$ & $-0.0350$ & $-0.0574$\\
$600$  & $-0.0125$ & $-0.0117$ & $-0.0315$\\
$2400$ & $-0.00313$ & $-0.00292$ & $-0.0150$\\
\bottomrule
\end{tabular}\end{center}

\noindent Three parameter scans pin the constant $c$ and kill two conjectures of ours:
\begin{itemize}\itemsep1pt
\item $\alpha$, fixed $p+q=7,d=8,n=600$: $c$ by $|\text{bias}|\!\cdot\! n$ equals
$7.79,7.20,8.58,7.15,7.80$ at $\alpha=0.1^\circ,0.5^\circ,2^\circ,10^\circ,30^\circ$ and
then rises to $14.17$ ($60^\circ$) and $19.26$ ($89^\circ$). Flat exactly where the
certificate needs to hold, worst where $\lambda_{\min}(\Gr)\approx1$ anyway.
\item $p+q$, at $n=200/600/1200$: $0.38/0.94/-0.76$ ($p+q=2$);
$5.70/8.53/6.25$ ($p+q=7$); $9.45/11.39/12.75$ ($p+q=14$).
So $c\ne p+q$: it grows sublinearly and saturates. A law of the form
$(p+q)/n$ is falsified.
\item $d$, at $p+q=7,n=600$: $3.01,4.05,5.87,8.68,11.27,9.06,8.84$ for $d=3,4,6,8,12,16,24$.
$c$ \emph{increases} with $d$ and saturates around $9$; our predicted $O(1/d)$
correction had the wrong sign.
\end{itemize}

\section{The leakage theorem and the tight worst case}\label{sec:leakage}

The first-order expansion of the correction coefficients is
$\hat\psi-\psi^\star=(G^{\top}G)^{-1}G^{\top}r$ with
$r=\operatorname{vec}(\Delta U+\Sigma)$, a residual independent of the fitted class.
Projecting on the smallest singular triple $(u_{\min},\sigma_{\min},v_{\min})$ of $G$ and
using $\|\sum_jv_jF_j\|_F/\|v\|=1$ from Definition~\ref{def:angle} gives the exact
statement we need.

\begin{theorem}[identifiability of the correction, worst case]\label{thm:tight}
For every dictionary, every $\alpha$, every $n>d$, and every $\Delta,\Sigma$,
\begin{equation}\label{eq:tight}
  \|\hat C-C^\star\|_F\le
  \frac{\|\Delta U+\Sigma\|_F}{\sigma_{\min}(G)}
  \le\frac{\|\Delta\|_F+\sigma\sqrt d}{\sin\alpha_{\min}}\,
     \sqrt{1+2\cos^2\alpha_{\min}\,\Cnorm}\,(1+o(1)),
\end{equation}
where $C^\star=\sum_j\psi_j^\star F_j$ and $\hat C=\sum_j\hat\psi_jF_j$.
The leading constant is exactly $1$ and cannot be improved: for
$\Delta:=W\,U^{\top}(UU^{\top})^{-1}$, rescaled to any prescribed
$\|\Delta\|_F$ (solvable whenever $n>d$), one has $\operatorname{vec}(\Delta U)\parallel
u_{\min}$ and \eqref{eq:tight} is attained.
\end{theorem}

\begin{corollary}[average case is a dimensional artifact]\label{cor:dilution}
If the direction of $\Delta$ is isotropic, then a \emph{single} direction contributes
$|\langle\operatorname{vec}(\Delta U),u_{\min}\rangle|\approx\|\Delta U\|_F/d$, but the correction
\emph{block} accumulates $q$ such components: by rotational invariance
$\|\mathcal P_{\mathcal F}\Delta\|_F^{2}/\|\Delta\|_F^{2}\sim\operatorname{Beta}(q/2,(d^{2}-q)/2)$,
whence, to leading order in $1/d$,
\[
  \mathbb{E}\|\hat C-C^\star\|\;\approx\;
  \frac{\mathbb{E}\|\mathcal P_{\mathcal F}\Delta\|_F}{\sin\alpha_{\min}}
  \;=\;\kappa(q)\,\frac{\sqrt q}{d}\,\frac{\|\Delta\|_F}{\sin\alpha_{\min}},
  \qquad \kappa(q):=\sqrt{\tfrac2q}\,\frac{\Gamma\!\bigl(\frac{q+1}{2}\bigr)}{\Gamma\!\bigl(\frac q2\bigr)}\;\uparrow\;1 .
\]
The average case thus vanishes as $d\to\infty$ at a $\sqrt q/d$ rate --- the $\sqrt q$ is block
aggregation, not physics ($\kappa(1)=0.798$, $\kappa(3)=0.921$, $\kappa(6)=0.960$), while the
worst-case bound \eqref{eq:tight} does not dilute at all.
Noise, by contrast, does not dilute: $\operatorname{vec}(\Sigma)$ is isotropic in $dn$
dimensions, giving $\approx\sigma/\sin\alpha_{\min}$ with no $d$.
\end{corollary}

\begin{remark}[consequence for reporting]\label{rem:reporting}
In high- or infinite-dimensional SciML (PDEs, rich dictionaries) ``random
misspecification is harmless'' is a statement about the dimension, not about the
physics. Only the worst-case normalization of Theorem~\ref{thm:tight} is a certificate.
\end{remark}

\begin{remark}[fitted constants in prior versions of this work are not free --- and one decomposition was wrong]\label{rem:fitted}
The measured ratio $\text{err}/\text{bound}=0.144$ at $d=10$, $q=3$ is a $\sqrt q$-type dilution times
the dictionary multiplier: $\kappa(3)\sqrt3/10=0.160$ against the $c_1\approx0.155$--$0.16$ fitted in
two independent designs (the adversarial table, and A's two-arm least-squares experiment; scripts
\texttt{b40}, \texttt{b41}). An earlier version of this Remark wrote that number as
$(1/d)\times\sqrt{1+2\cos^{2}\alpha\,\Cnorm}=0.10\times1.08$, which is $0.108$, not the $0.155$ it
claimed to reproduce; the missing factor is $\sqrt q$, i.e.\ block aggregation, and the $q$-scan of
the $q$-scan \texttt{b41} decides the question rather than the arithmetic: $c_1$ at $q=1,2,3,4,6$ is
$0.077$, $0.124$, $0.151$, $0.183$, $0.227$, so it is not $q$-blind, and
$c_1/(\kappa(q)\sqrt q/d)\in[0.94,0.99]$ with no trend in $q$. Both factors are computable before
seeing data ($\Cnorm$ measured $0.123$--$0.134\approx1/d$). The companion constant $c_2\approx0.063$
is $\sqrt{q/n}=\sqrt{3/600}=0.0707$ to within $13\%$ --- a Gaussian quantity over $\sqrt n$ with the
same block factor, and independent of $d$ --- not a physical constant. Hence the two-scale law
$c_1\|\Delta\|+c_2\sigma$ of an earlier draft should be read as ``one tight bound plus one dilution
factor, each carrying its own $\sqrt q$'', and $c_1,c_2$ must not be presented as universal.
\end{remark}

\section{The certificate and what it costs to use it}\label{sec:cert}

\begin{equation}\label{eq:star}
  \boxed{\ \|\hat C-C^\star\|_F\ \le\
  \frac{\|\Delta\|_F+\sigma\sqrt d}{\sin\alpha_{\min}}\cdot
  \sqrt{1+2\cos^2\alpha_{\min}\cdot\|\Le^{-1}H\|_2^2}\ }
  \tag{$\star$}
\end{equation}
Every quantity on the right is computable from the dictionary, the declared
model-form-error budget $\|\Delta\|_F$, and the noise level: there is no fitted constant,
no resampling, and no unknown spectral quantity. The sample-size factor of
Theorem~\ref{thm:gammaclosed} may be appended as
$(1-c/2n-0.88/\sqrt n)^{-1}$; Section~\ref{sec:cancellation} shows where it belongs and
where it does not.

\subsection{$(\star)$ holds with no sample-size margin, in the only test that can
detect a failure}\label{sec:adversarial}

Random-direction $\Delta$ tests of a certificate are \emph{vacuously true}: at
$p+q=7$, $n=600$, $\sigma=10^{-4}$, the ratio $\text{err}/\text{bound}$ sits in
$0.13$--$0.17$ and is flat in $n$ from $60$ to $1800$ --- a $30\times$ range in which the
remainder of Theorem~\ref{thm:gammaclosed} changes by more than $10\%$. Corollary
\ref{cor:dilution} explains it: the leakage is diluted by $1/d$, so the test never gets
close to the bound. Any claim of the form ``the certificate passes at $n=\ldots$'' made
on a random arm carries no information about the margin. The adversarial arm of
Theorem~\ref{thm:tight} is required.

\begin{center}\begin{tabular}{lrrrrrrc}
\toprule $n$ & $\alpha=2^\circ$ & $5^\circ$ & $20^\circ$ & $45^\circ$ & $85^\circ$ &
max($85^\circ$) & violations\\
\midrule
$30$   & $0.880$ & $0.879$ & $0.890$ & $0.927$ & $0.983$ & $0.992$ & $0/120$\\
$60$   & $0.883$ & $0.884$ & $0.895$ & $0.930$ & $0.988$ & $0.992$ & $0/120$\\
$100$  & $0.884$ & $0.885$ & $0.896$ & $0.931$ & $0.990$ & $0.993$ & $0/120$\\
$200$  & $0.887$ & $0.887$ & $0.896$ & $0.933$ & $0.991$ & $0.993$ & $0/120$\\
$600$  & $0.887$ & $0.888$ & $0.897$ & $0.935$ & $0.992$ & $0.993$ & $0/120$\\
$1800$ & $0.887$ & $0.886$ & $0.897$ & $0.933$ & $0.993$ & $0.993$ & $0/120$\\
\bottomrule
\end{tabular}\end{center}

\noindent $36$ cells $\times\ 120$ trials $=4320$ adversarial runs
($\operatorname{vec}(\Delta U)\parallel u_{\min}(G)$, $\|\Delta\|_F=0.05$,
$d=10,p=4,q=3$): \textbf{zero violations with the margin set to $1$}, maximum observed
ratio $0.993$. Note the counterintuitive ordering: the bound is \emph{tightest} at large
$\alpha$ ($0.983\to0.993$), not small, because at small $\alpha$ the bound itself
shrinks. (Script: \texttt{b27\_adversarial\_margin.py}.)

\subsection{Why the ratio is flat in $n$: an identity, measured}
\label{sec:cancellation}

Under the adversarial construction the achieved ratio factorizes
\begin{equation}\label{eq:factor}
  \frac{\text{err}}{\text{bound}}=q_4\cdot q_1\cdot q_2\cdot q_3,\qquad
  q_4=\sqrt{\frac{\mathrm{LB}}{\lambda_{\min}(\Gr)}},\quad
  q_1=\sqrt{\frac{\lambda_{\min}(\Gr)}{\lambda_{\min}(G^{\top}G/n)}},\quad
  q_2=\frac{\|\Delta U\|_F}{\sqrt n\,\|\Delta\|_F},\quad
  q_3=\|v_{\mathcal F}\| ,
\end{equation}
where $\mathrm{LB}=\sin^2\alpha/(1+2\cos^2\alpha\Cnorm)$ is the lower branch of
Theorem~\ref{thm:twosided} and $v_{\mathcal F}$ is the correction-block part of the
minimum right-singular direction of the design matrix: the leakage is the
\emph{amplitude} divided by $\minr G$, but only the fraction $\|v_{\mathcal F}\|$ of that
amplitude lands in the correction block. Thus $q_4$ is \emph{the slack of the dictionary
theorem}, $q_1$ \emph{the sampling shortfall} measured in Theorems~\ref{thm:bias} and
\ref{thm:gammaclosed}, $q_2$ \emph{the price of realizing $\Delta U$} with a
minimum-norm operator, and $q_3$ \emph{the block mass}. The four factors multiply to the
observed ratio per trial (script \texttt{b32}: $q_4q_1q_2q_3=0.8905$ against a measured
$0.8905$ at $\alpha=2^\circ$, $n=200$); an earlier draft of this chapter omitted $q_3$
and left an unexplained ``residual'' column of $\approx5\%$ in the table below, which is
$1-q_3$ read off column means ($q_3=0.943$--$0.950$ at $\alpha=2^\circ$,
$0.9985$--$0.9995$ at $85^\circ$; the direct per-trial measurement over $60$ trials gives
$0.9466$ (per-trial sd $0.0066$) and
$0.9993$) --- see Section~\ref{sec:notclaims}, item (h).
\begin{center}\begin{tabular}{lrrrrrrr}
\toprule $n$ & $\alpha$ & err/bnd & $q_4$ & $q_1$ & $q_2$ & $q_1q_2$ & $1-q_3$\\
\midrule
$30$   & $2^\circ$ & $0.8859$ & $0.9482$ & $\mathbf{1.0733}$ & $\mathbf{0.9190}$ & $0.9863$ & $0.0464$\\
$60$   & $2^\circ$ & $0.8908$ & $0.9484$ & $1.0453$ & $0.9475$ & $0.9904$ & $0.0473$\\
$200$  & $2^\circ$ & $0.8916$ & $0.9478$ & $1.0187$ & $0.9754$ & $0.9936$ & $0.0498$\\
$600$  & $2^\circ$ & $0.8920$ & $0.9478$ & $1.0050$ & $0.9905$ & $0.9955$ & $0.0514$\\
$1800$ & $2^\circ$ & $0.8920$ & $0.9476$ & $1.0018$ & $0.9959$ & $0.9977$ & $0.0534$\\
$30$   & $20^\circ$& $0.8945$ & $0.9531$ & $1.0656$ & $0.9269$ & $0.9871$ & $0.0433$\\
$600$  & $20^\circ$& $0.9029$ & $0.9538$ & $1.0061$ & $0.9894$ & $0.9954$ & $0.0464$\\
$1800$ & $20^\circ$& $0.9010$ & $0.9528$ & $1.0035$ & $0.9941$ & $0.9976$ & $0.0495$\\
$30$   & $85^\circ$& $0.9901$ & $0.9996$ & $1.0896$ & $0.9129$ & $0.9946$ & $0.0015$\\
$1800$ & $85^\circ$& $0.9988$ & $0.9996$ & $1.0086$ & $0.9913$ & $1.0000$ & $0.0005$\\
\bottomrule
\end{tabular}\end{center}
\noindent The last column is $q_4q_1q_2-\text{err}/\text{bound}$ under column means, which
is why it is not exactly $1-q_3$; the per-trial closure of \eqref{eq:factor} is the
statement in the text above.

\begin{theorem}[tightness is a property of $\Gr$; measured]\label{thm:cancel}
$q_4$ is independent of $n$, $\sigma$ and the dictionary scaling: at $\alpha=2^\circ$ it
is $0.9476$--$0.9484$ over $n=30\ldots1800$ ($0.08\%$ spread), and $0.9996$ at
$\alpha=85^\circ$. Moreover $q_1-1\approx3/n$ and $q_2=1/q_1\cdot(1\pm1\%)$ cell by cell,
so $q_1q_2\in[0.986,1.000]$ throughout the i.i.d.\ grid. Section~\ref{sec:notclaims},
item 3, measures the correlated (trajectory) design arm, where the product drifts down to
$0.975$: the cancellation degrades but always in the safe direction.
\end{theorem}

\begin{proof}[Mechanism]
Both $q_1$ and $q_2$ are lower-edge Marchenko--Pastur effects of the \emph{same} data
matrix $U$: $\lambda_{\min}$ of the design Gram is short by the same amount that the
minimal stretch of the map $\Delta\mapsto\Delta U$ is short, and a minimal-norm
realization of a fixed $\Delta U$ becomes exactly that much \emph{cheaper} in
$\|\Delta\|_F$. The two factors are one phenomenon read from two sides, so their product
is $\approx1$.
\end{proof}

\begin{remark}[$q_4$ is not measured --- it is a formula]\label{rem:q4analytic}
Because the upper side of Theorem~\ref{thm:twosided} is attained per sample
(Section~\ref{sec:twosided}, script \texttt{b29}), the dictionary factor is available in
closed form:
\begin{equation}\label{eq:q4formula}
  q_4=\sqrt{\frac{\mathrm{LB}}{\lambda_{\min}(\Gr)}}
  \;\approx\;\Bigl(\frac{1+\cos^2\alpha\,\Cnorm}{1+2\cos^2\alpha\,\Cnorm}\Bigr)^{1/2},
\end{equation}
and the right-hand side uses only the dictionary. With
$\Cnorm=0.12809$ (median over $300$ draws at $d=10$, $p+q=7$) it predicts
$q_4=0.94770$ against the measured $0.9478$ at $\alpha=2^\circ$, and
$0.99951$ against $0.9996$ at $\alpha=85^\circ$ --- both to $10^{-4}$
(script \texttt{b30}). The $5.2\%$ that the dictionary factor removes from $(\star)$
is thus \emph{exactly} the $\varepsilon=1$ Young step of the lower branch: no sampling,
no noise, and no dictionary scale enters it.
The claim was then checked over the whole angle range, not just at two points: $4400$
dictionaries at $d=10$ over eleven $\alpha$'s give per-sample deviation of
the closed form \eqref{eq:q4formula} from the exact factor of $0.0000$--$0.0020$ in the
mean (worst
single draw $0.0040$, again at $\alpha=45^\circ$) with zero upper-side violations, and
the fitted $q_4$ at $\alpha=20^\circ$ in the table above --- obtained from a
least-squares solve, hence not equivalent to the closed form by construction --- is
$0.9528$--$0.9538$ against a predicted $0.9525$. (Script \texttt{b31}.)
The deviation is $1.5\%$ at $d=8$ (script \texttt{b29}) and $0.4\%$ at $d=10$: smaller in
better-conditioned dictionaries, although we have not established the scaling law.
The fit-based check exists at three angles only; the other eight are checked against the
exact dictionary factor, which is the same statement one step earlier in the chain.
Proposition~\ref{prop:plane} sharpens the display in two ways at once. Since
$\lambda_{\min}(\Gr)\le\operatorname{cf}(t,\alpha)$, the dictionary-only expression
$\sqrt{\mathrm{LB}/\operatorname{cf}}$ is a \emph{proven lower bound} for $q_4$ rather than an
approximation --- so no sampling, noise or dictionary scale accounts for the looseness $q_4$
measures, which is carried entirely by the dictionary side --- and because $\operatorname{cf}$ sits
closer to $\lambda_{\min}(\Gr)$ it approximates $q_4$ to half that slack under the square root:
$\le1.1\times10^{-3}$ relative at $d=8$, against the $1.5\%$ of \eqref{eq:q4formula}'s right-hand
side at the same $d$. Theorem~\ref{thm:exact} then removes the gap altogether: substituting its root
$\lambda^{\star}$ for $\lambda_{\min}(\Gr)$ makes $q_4$ a computed quantity rather than an estimated
one, and Corollary~\ref{cor:square} shows the closed form \emph{is} exact whenever $p=q$.
Remark~\ref{rem:exactcost} states what that choice costs and why we did not take it.
\end{remark}

\begin{corollary}[the observed looseness of $(\star)$ is one dictionary scalar]
\label{cor:slackone}
Combining \eqref{eq:q4formula} with $q_3\approx\sqrt{\lambda_{\min}(\Gr)/\sin^2\alpha}$
(the dictionary-side surrogate for the block mass; it agrees with $\|v_{\mathcal F}\|$ to
$0.6\%$ in the mean and $2.1\%$ on the worst trial, script \texttt{b32}) and
$q_1q_2\in[0.986,1.000]$ gives
\begin{equation}\label{eq:slackone}
  \frac{\text{err}}{\text{bound}}
  \;\approx\;\Bigl(\frac{\mathrm{LB}}{\sin^2\alpha}\Bigr)^{1/2}
  \;=\;\bigl(1+2\cos^2\alpha\,\Cnorm\bigr)^{-1/2},
\end{equation}
which contains neither $\lambda_{\min}(\Gr)$ nor any sampled quantity.
Predictions against the fitted values of the table above:
$0.8923$ vs $0.8859$--$0.8920$ ($\alpha=2^\circ$), $0.9031$ vs $0.8945$--$0.9029$
($20^\circ$), $0.9990$ vs $0.9901$--$0.9988$ ($85^\circ$); the largest discrepancies are the
$n=30$ cells, where $q_1q_2$ is at its minimum ($0.986$--$0.987$).
In words: how loose the certificate is in practice is decided by the same
$\varepsilon=1$ Young step, and it is readable off the dictionary before the fit.
\end{corollary}

\begin{remark}[where the sample-size remainder belongs --- and where it must not appear]
\label{rem:marginplace}
The proof of Theorem~\ref{thm:cancel} is unavailable to the certificate's derivation:
the proof chain goes through the one-sided inequality
$\sigma_{\min}(G)\ge\sqrt{n\lambda_{\min}(\Gr)}(1-c/2n-0.88/\sqrt n)$, which cannot see
the $q_2$ compensation. Hence the factor $(1-c/2n-0.88/\sqrt n)^{-1}$
\textbf{must appear in the statement of the theorem and must not appear in any numerical
table.} Extrapolating it to $n=30$ predicts $\text{err}/\text{bound}\approx1.02$ (a
violation) against a measured $0.886$; including it in a results table would tell a
reader that the certificate degrades with sample size, which is the opposite of the
measurement. We flag this explicitly because our own draft made the prediction before
the measurement and was wrong in direction twice: once on the sign of the effect
(a larger margin enlarges the bound and \emph{lowers} the ratio; all $25$ cells satisfied
$\text{new}<\text{old}<\text{raw}$) and once on its size (the $c/n+1.75/\sqrt n$
extrapolation overestimates $q_1-1$ by a factor $1.6$).
\end{remark}

\section{Limitations, retracted claims, and open items}\label{sec:notclaims}

\begin{enumerate}\itemsep2pt
\item \textbf{Constants are calibrated, not derived.} $\kappa\approx0.15$ in
Theorem~\ref{thm:sampling} and $c\le9$/$c\le20$ in Theorem~\ref{thm:gammaclosed} are
empirical envelopes over scanned ranges ($d\in[3,24]$, $p+q\in[2,14]$,
$\alpha\in(0,89^\circ]$, $n\in[200,2400]$). A proof of $c(d,p+q)$ exists only as a
claim that a net argument gives \emph{some} constant; we judged the analytic derivation
not worth the effort, since it would replace a $21$-cell-supported number with an
uglier, looser bound. The honest label is ``explicit constant, empirically validated''.
\item \textbf{A regime we have not tested: $n$ comparable to $p+q$.} Our tables keep
$n> p+q$ by a factor of at least $14$. With an effectively infinite dictionary
($p+q\sim100$ at $n=600$) the remainder of Theorem~\ref{thm:sampling} is $0.19$ and the
theorem is heuristic.
\item \textbf{Gaussian $U$ is load-bearing for the tightness, not for the safety --- now
measured.} The identity $q_1q_2\approx1$ uses the lower edge of a Gaussian design, and the
hand-off prediction was that time-series designs break it. We broke it ourselves first:
columns of $U$ taken as consecutive states of $d$ independent $\operatorname{AR}(1)$ series
(unit stationary variance, so $\rho=0$ is exactly the arm of Section~\ref{sec:cancellation})
give, at $\alpha=2^\circ$,

\smallskip
\centerline{\begin{tabular}{lrrrrr}
\toprule $\rho$ & $n$ & $q_1$ & $q_2$ & $q_1q_2$ & $\mathrm{err}/\mathrm{bound}$\\
\midrule
$0.00$ & $200$ & $1.017$ & $0.977$ & $0.9935$ & $0.8923$\\
$0.50$ & $200$ & $1.026$ & $0.967$ & $0.9916$ & $0.8886$\\
$0.80$ & $200$ & $1.059$ & $0.934$ & $0.9877$ & $0.8879$\\
$0.95$ & $200$ & $1.127$ & $0.871$ & $0.9747$ & $0.8776$\\
$0.95$ & $600$ & $1.062$ & $0.930$ & $0.9855$ & $0.8852$\\
\bottomrule
\end{tabular}}
\smallskip

\noindent Both factors move as predicted ($q_1$ up, $q_2$ down) and the product \emph{leaves
$1$ downward}: the cancellation degrades from $0.9935$ to $0.9747$ at $\rho=0.95$, $n=200$,
and $q_1q_2$ stays monotone in $\rho$ and closer to $1$ at $n=600$. Every cell satisfies
$\mathrm{err}/\mathrm{bound}<1$, and the four-factor identity \eqref{eq:factor} still closes
per trial ($0.8776$ against $0.8776$). So correlation makes the certificate \emph{looser}
than the dictionary scalar \eqref{eq:slackone} predicts ($1.6\%$ at $\rho=0.95$, $n=200$),
which is the safe direction, and it does not resurrect any need for a calibrated term.
We do not claim a rate: the departure is not consistent with $(p+q)/n_{\mathrm{eff}}$ for
$n_{\mathrm{eff}}=n(1-\rho)/(1+\rho)$ (that would be $1.36$ where we measure $0.025$), and
the same ordering holds at $\alpha=20^\circ$. One diagnostic worth recording: at
$\rho=0.95$, $n=200$ the mean of the per-trial product $q_1q_2$ is $0.9747$ while the
product of the two means is $0.9814$ --- that $0.7\%$ gap \emph{is} the cancellation, since
$q_1$ and $q_2$ are negatively correlated trial by trial, and it is why every closure claim
in this chapter is stated per trial rather than from column means.
What remains genuinely untested is a \emph{learned} dictionary or an MLP-regression design,
where $q_2$ can lose its structure in a direction we cannot parameterize by $\rho$;
we still consider small-sample overconfidence there to be the disappearance of $q_2$, not a
generic small-$n$ effect (script \texttt{b34}).
\item \textbf{The direction dichotomy is exact only for the constructed angle model.}
Definition~\ref{def:angle} prescribes $F_j=\cos\alpha\,a_j+\sin\alpha\,b_j$ with $b_j$
exactly orthogonal to $\operatorname{span}\{E_i\}$. Real adaptive dictionaries (learned
corrections, kernel EDMD dictionaries) do not satisfy this; $\alpha$ then becomes a
distribution of principal angles and $\sin\alpha_{\min}$ must be estimated, which is a
different (statistical) problem, addressed in the companion chapter on coverage rather
than here.
\item \textbf{What we retracted.} (a) $\sigma_{\min}(\Gr^{1/2})=\sin\alpha/
\sqrt{1+\cos^2\alpha\Cnorm}$ as an equality --- it is an upper bound, unsafe direction
(Remark~\ref{rem:direction}); (b) the operator-norm remainder $O(d/(\sqrt n\sin^2\alpha))$
--- wrong by $7$ orders of magnitude; (c) the remainder shape $\kappa\sqrt{(p+q)/n}$ as a
\emph{prediction of the size} --- provable but loose, real fluctuation is
$0.68/\sqrt n$; (d) the conjectured law $|\text{bias}|= (p+q)/n$ and an $O(1/d)$
correction --- both falsified by the scans of Section~\ref{sec:gammaclosed}, the second
with the wrong sign; (e) the extrapolated sample-size margin as a numerical-table entry
--- see Remark~\ref{rem:marginplace}; (f) the relative-bias shape
$|\text{bias}|\propto B_n/(n\lambda_{\min}(\Gr))$ written with an extra factor $1/n$ ---
it conflated the absolute displacement $n\beta_2=O(1)$ with the relative one
$\varepsilon_n=-\beta_2/(2\lambda_1)=O(1/n)$, and the wrong form predicts a
$\alpha\to0$ divergence that Section~\ref{sec:gammaclosed} measures as flat; (g) a
reported value $\|\Le^{-1}H\|_2^2=1.06$ produced by a Cholesky half-inverse rather than
a solve (see the pitfall note in Section~\ref{sec:setting}); the correct figure at $d=8$
is $0.17$; (h) \eqref{eq:factor} as written in an earlier draft, with three factors and the
claim that it ``factorizes exactly'' --- the identity closes only after the block-mass
factor $q_3=\|v_{\mathcal F}\|$ is included, and the $\approx5\%$ column we then labelled
``residual'' is precisely the missing factor. Discovered by checking the sentence
``the observed looseness is $q_4$ alone'' against the table underneath it, not by an
experiment; (i) the conjecture that the sphere minimisation behind $\operatorname{cf}$ is attained at
the top right-singular direction $w_0$ --- falsified by an exhaustive circle search at $q=2$, after
$3200$ random-direction draws had missed it because random points do not resolve $S^2$
(Remark~\ref{rem:w0}); (j) the decomposition of the fitted $c_1\approx0.16$ as $(1/d)\times
\sqrt{1+2\cos^2\alpha\,\Cnorm}$ --- that product is $0.108$, and the missing $\sqrt q$ is block
aggregation, settled by the $q$-scan \texttt{b41} rather than by arithmetic; Corollary~\ref{cor:dilution}
carried the same omission; (k) the sentence that $\lambda_{\min}(\Gr)$ is ``obtainable to round-off''
by the root find of Theorem~\ref{thm:exact} --- it is obtainable to \emph{absolute} round-off, and
the conditioning scan \texttt{b43} measures the law $\mathrm{relerr}\times\lambda_{\min}(\Gr)\approx1.7\cdot10^{-15}$ over six
decades, so at $\alpha=0.001^\circ$ the same code is good only to $10^{-5}$ relative. The overclaim
survived three rounds of QA because every script that checked it ($\alpha\ge1^\circ$) sat in the regime
where absolute and relative agree; (l) the \emph{attribution} of that $\le3.4\%$ residual and of the
$\operatorname{cf}-\lambda_{\min}$ sign flip to ``the angle distribution''. The anisotropy design \texttt{b44} holds the
cosine spectrum fixed and switches $\Le$ from a generic Gram to $\gamma I_p$, which removes eight to
twelve orders of the gap, and holds $\Le$ fixed and switches the spectrum, which moves it by a factor
$13$--$300$: the variable is the anisotropy of $\Le$, and the flip additionally requires a second
substantial cosine. The numbers were not in question; the label was, and a wrong label sends a future
reader to the wrong relaxation. The same script's attempt to separate the \emph{alignment} of the weak
dictionary direction from its size failed --- its ``aligned'' and ``transverse'' cells are one and the
same ensemble, because the coupling's singular vectors were drawn independently of $\Le$ --- so
alignment is untested here, not ruled out;
(m) the surviving clause of (l) itself, that the flip ``\emph{additionally} requires a second
substantial cosine''. The two-cosine scan \texttt{b45} fixes the geometry and runs the two cosines as independent axes ---
a clean $2\times2$ plus two seven-point ladders, $1800$ draws --- and $\operatorname{cf}$ lands below
$\lambda_{\min}(\Gr)$ in $100/100$ draws at a second cosine of $0.087$, and again in $100/100$ at
$\operatorname{cf}=0.133$. Neither ``a second substantial cosine'' nor the substitute I had ready
(``$\operatorname{cf}$ small enough'') is the criterion. What decides is the atom structure of the
measure $\nu_v$ of Proposition~\ref{prop:side}: with $\Le$'s weak eigenvector along any single left
singular direction of $K$ the measure has one atom, $\operatorname{slack}=0$ to round-off in twelve
cells, and the gap stays on the $\operatorname{eigvalsh}$ floor; at $45^\circ$ to that basis the
measure is genuinely multi-atom and every cell flips, with
$|\operatorname{gap}/(\operatorname{cf}\cdot\operatorname{slack})|=1.00$. The spectrum scales
$\operatorname{slack}$; it does not select its sign. That run also repaired (l)'s broken control ---
drawing $\Le$ \emph{after} $K$ and pinning the weak eigenvector to $K$'s own left singular structure
--- so the two families agree: alignment is not a competing explanation, it is the atom count seen
from the other side;
(n) the reading of the family scan \texttt{b47} as evidence that no direction beats
$\operatorname{cf}_{1}$. The polar-net search \texttt{b48} replaced sampling --- a net of $96\times240$
directions per draw, $7{,}380{,}480$ family values --- and directions \emph{do} beat it, at
\emph{first} order: along $v(t)=\cos t\,v_{1}+\sin t\,u$ with $u\perp v_{1}$ one has
$B'(0)=0$ but $C'(0)=2v_{1}^{\top}K^{\top}\Le^{-2}Ku$, so
$\lambda'(0^{+})=-2\,v_{1}^{\top}K^{\top}\Le^{-2}Ku\cdot\lambda(1-\lambda)/|\partial_\lambda F|$ is
generally nonzero, and $\operatorname{cf}_{1}$ is a saddle of the family rather than its minimum
($84/84$ generic draws, $\|P_{v_{1}^{\perp}}K^{\top}\Le^{-2}Kv_{1}\|=1.5\cdot10^{-3}\ldots2.1\cdot10^{-1}$,
measured slope matching the formula to $1.7\cdot10^{-2}$ relative). The correct statement is a
criterion: $v_{1}$ is stationary iff $K^{\top}\Le^{-2}Kv_{1}\parallel v_{1}$, which holds on the
$\operatorname{slack}=0$ loci ($\dim E_{0}=q$: equal angles, isotropy; there the search finds
$0$ competitors and $\|P_{v_{1}^{\perp}}\cdot\|\le6.2\cdot10^{-16}$) and fails generically. Two
artifacts of my own are retracted with it: the first version of that search took $v_{1}$ to be an arbitrary
eigenvector of $E_{0}$, which at $\dim E_{0}=3$ is not the $C$-maximiser --- its self-check caught a
$1.4\cdot10^{-2}$ discrepancy; and a degenerate-cell ``violation'' of Proposition~\ref{prop:chord}
($15{,}793$ of them) came from normalising a round-off-sized vector, which amplifies its leftover
$v_{1}$-component so that $v(t)$ leaves the sphere, $B>1$, and the smaller root of
\eqref{eq:family} turns negative. Proposition~\ref{prop:chord} itself is \emph{not} in trouble: on
the corrected net and on $144{,}000$ ray points no family member falls below
$\lambda_{\min}(\Gr)$, and every competitor pays intercept $\bigl(B/c^{2}\le1-10^{-6}\bigr)$ for slope
$\bigl(C/(c^{2}t_{1})=1.000005\ldots1.000499\bigr)$, as the implicit-function signs predict.
(o) the attribution of $\operatorname{cf}_1$'s residual gap to the \emph{direction} $v_1$, and with it
the conjecture that $\min_v\lambda_1(v)=\lambda_1(v_\ast)=\lambda_{\min}(\Gr)$. The identity
\eqref{eq:fibre} is true --- verified as an identity on $320/320$ draws to $2.5\cdot10^{-15}$ --- but
its content is the \emph{exact} fibre crossing $\lambda_{\rm ex}(v_\ast)$, not the $n=1$ member:
$\lambda_1(v_\ast)$ stays $0.50$--$1.01$ of the gap above the certificate, and in the one
$\operatorname{dim}E_0=3$ cell even above $\operatorname{cf}_1=\lambda_1(v_1)$ on all $40$ of its draws
($\max(\operatorname{cf}_1-\lambda_1(v_\ast))/\text{gap}=-4.9\cdot10^{-5}$), which violates nothing (Proposition~\ref{prop:chord} never claims
$\lambda_1(v)$ is smaller at good $v$) but removes the motivation for the search. The residual is
truncational (Remark~\ref{rem:tight}).
(p) the use of uniform or semi-uniform direction nets as evidence about $\min_v\lambda_n(v)$, and with
it the acceptance thresholds tested by the direction-net run \texttt{b51}. Its $\min_{\rm net}\lambda_1=0.0000$ in
$7$ of $8$ cells is a property of the angle list, not of the functional: the left endpoint of a
$\operatorname{linspace}(10^{-3},\pi/2,n_t)$ list is that statistic's ceiling once the descent set has
width $t_{\rm hi}\le1.6\cdot10^{-3}$, so flatness in $n_t$ was construction, and the $>0.99$ /
$\le0.90$ thresholds of the previous draft measured the grid. Uniform nets are retired as a measurement
device for this question; grid-free routes (the analytic ray, the eigenvector block) replace them.
(q) the strength of the ray leg's gap fraction, and with it the reading of (iv) as a per-draw
guarantee. The ray scan \texttt{b55} falsified the claim exactly as it had been falsified-in-advance: on new tables
(doubled $p$, doubled $q$, new seed) the analytic two-jet ray's $\lambda_2$ value drops to $0.7387$
(median) in one $q=6$ cell, and along a $q$-ladder the per-cell medians fall from $1.0000$ at $q=2$ to
$0.6547$ at $q=8$ with $20/20$ draws below $0.90$ at the top of the ladder. Since the ranges I quoted in
(iv) were \emph{per-cell medians} and are now labelled as such, no number in this chapter changes sign,
but the sentence ``the grid-free route lifts $\lambda_2$ to $\ge0.9965$'' must be read as a statement
about $q\le3$ tables, and the ray leg is closed as a research direction rather than extended. What
survives on the same tables is the eigenvector block: $(\operatorname{cf}_1-\lambda_2(v_\ast))/\text{gap}$
prints $1.00$ at both its median and its maximum in all $20$ cells, i.e.\ the deficit is below
$5\cdot10^{-3}$ of the gap throughout (b53's tighter readout of the same quantity gave
$3.3\cdot10^{-7}$ to $5.0\cdot10^{-4}$).
\item \textbf{What this chapter does not do.} It gives no real-data validation. An
attempted cyclostationary ENSO leg was closed after three independent failure modes were
identified (a small-denominator statistic, non-replication across window choices, and a
delay-embedding basis artifact that a zero-physics control reproduced in full); the
surviving negative content --- a seasonal-normal null genuinely cannot produce the
observed numbers, and $\|A^n\|_2/\rho^n$ being two orders of magnitude more stable than
trace-normalized alternatives under basis changes --- is reported as a separate
methodological note, not as evidence for the certificate.
\item \textbf{Closed, then retracted, then closed again: the direction $w_0$.} The quantity
$\operatorname{cf}(t,\alpha)$, $t=\|\Le^{-1}H\|_2^2$, first appeared in a draft as a numerically
suggested candidate; Proposition~\ref{prop:plane} derives it as the Rayleigh minimum on an explicit
plane. The conjecture that followed --- that the minimiser over the correction sphere is the top
right-singular direction $w_0$ of $C$ --- is \emph{false}, and we found that by searching instead of
by sampling: a dense circle search beats $\mu(w_0)$ in $143$ of the $180$ exhaustively-searched $q=2$
dictionaries at $\alpha\ge20^\circ$ --- $79\%$ of them, so the failure is typical rather than
exceptional (Remark~\ref{rem:w0}, script \texttt{b37}). What replaced it is stronger than
what it replaced. Theorem~\ref{thm:exact} gives $\lambda_{\min}(\Gr)$ as the unique root of one
monotone scalar equation in $\lambda_{\max}$ of an explicit $q\times q$ matrix --- exact to
$1.8\cdot10^{-11}$ over $6480$ dictionaries, nine $(d,p,q)$ configurations,
$\alpha=1^\circ\ldots89^\circ$ --- and Corollary~\ref{cor:square} shows that in the square case
$p=q$ that root is \emph{exactly} $\operatorname{cf}(t,\alpha)$ (verified to $1.4\cdot10^{-11}$ on
seven square configurations $\times$ five angles, script \texttt{b39}).
\emph{What remains open:} no closed form is known for $\lambda_{\max}(B(\lambda))$ when $q<p$; the
measured defect of the $\operatorname{cf}$ surrogate there is $\le1.9\cdot10^{-3}$ at $(d,p,q)
=(8,4,3)$ and $\le6.9\cdot10^{-3}$ at $(6,4,3)$, both at $\alpha=45^\circ$. So the practical question
(is the closed form adequate without a root find?) is answered yes in the scanned regime, and the
theoretical one (a $q<p$ closed form) is not answered. Proposition~\ref{prop:multi} sharpens which of
the two is left: the \emph{root find} needs no equal-angle hypothesis at all, so the open question
concerns the closed form only. And the prescribed-spectrum run \texttt{b42} put a price on the other relaxation one might reach
for --- a distribution of principal angles in place of the constructed single one:
$\lambda_{\min}(\Gr)$ moves by $\le3.4\%$, while the \emph{sign} of
$\operatorname{cf}-\lambda_{\min}$ moves by a full $180^\circ$. That price is right and its label was
wrong (retractions~(l) and (m) below): the $2\times5$ design \texttt{b44} attributes both effects to the
\emph{anisotropy} of $\Le$ --- at $\Le=\gamma I_p$ the gap is identically zero for every spectrum
(Corollary~\ref{cor:iso}) --- and the two-cosine scan \texttt{b45} names the anisotropy: the sign of the gap is the sign of
the moment defect $D(\operatorname{cf})$ of Proposition~\ref{prop:side}, whose leading term is
$\operatorname{cf}\cdot\operatorname{slack}$ with
$\operatorname{slack}=c^{2}t-\lambda_{\max}(E_0^{\top}K^{\top}\Le^{-2}KE_0)$, one joint functional of
$\Le$'s spectrum and $K$'s weights, and \emph{not} a condition on the cosine spectrum --- the flip
survives a second cosine of $0.087$ and a $\operatorname{cf}$ of $0.133$ alike. An angle-distribution
theory of the smallest eigenvalue is therefore a low-yield target for two independent reasons: at
isotropy the distribution is invisible to $\lambda_{\min}(\Gr)$, and where it is visible it enters only
through $\operatorname{slack}$, a quantity that is one eigenvalue of an explicit matrix. The functional
the previous draft listed as open is thus identified. Pairing the intercept $c^{2}$ with the slope
that actually attains it --- $t_{1}$ of \eqref{eq:chord} --- turns the diagnosis into a replacement:
$\operatorname{cf}_{1}$ bounds $\lambda_{\min}(\Gr)$ from above with \emph{no} angle, isotropy or
simplicity hypothesis (Proposition~\ref{prop:chord}), and its residual is second order (measured
log-log slope $2.081$, pre-registered as $2$). Two items stay open and the first has been settled in
the direction that matters for practice while staying open as a closed form: the polar-net search \texttt{b48} settled the
question ``is $\operatorname{cf}_{1}$ the smallest member of \eqref{eq:family}?'' in the
\emph{negative} (retraction~(n)), and the ray computation \texttt{b49} sized what the negative costs and buys. Along the
eigenvalue-free ray $u\parallel P_{v_{1}^{\perp}}K^{\top}\Le^{-2}Kv_{1}$ the depth of the descent is
$\operatorname{dip}=-L_1^{2}/(4L_2)$ of Remark~\ref{rem:chord}, verified to $10^{-3}$ relative on $280$
draws, and one step on that ray removes $0.0004$--$0.50$ of $\operatorname{cf}_{1}$'s gap at the median
($0.76$ worst draw, absolute $|\operatorname{dip}|\le1.2\cdot10^{-7}$) while a $64\times64$ net on the
same draws gets at most $\sim10\%$ deeper still. Evaluating \eqref{eq:family} at the resulting direction
gives $\lambda_{\min}(\Gr)\le\operatorname{cf}_{2}$ by Proposition~\ref{prop:chord} --- \emph{no new
proof}, one $2\times2$ root (tightening it below $\operatorname{cf}_{1}$ additionally needs $L_2>0$,
observed $280/280$ and not proved). What is deliberately \emph{not} claimed
is that $\operatorname{cf}_{1}+\operatorname{dip}$ bounds anything on either side --- the third jet of
$\lambda(t)$ is unbounded and the measured drift in $|c_1|$ has no consistent sign --- and the family
minimum over the whole sphere stays without a closed form: the one-dimensional fractional programme is
open, but its remaining marginal is now measured, and we stopped on it for that reason rather than for
lack of a technique. Proposition~\ref{prop:gauss} changes the shape of that open item, not its status:
since $\lambda_2(v)$ is a certified bound for \emph{every} unit $v$ and an explicit cubic root in four
scalars, $\min_{v}\lambda_2(v)$ over any finite set of directions is again certified (and at every
direction tested by \texttt{b50}, $\lambda_2(v)\le\lambda_1(v)$; the pointwise ordering of successive
Gauss rules is measured, not proved). What
\texttt{b50} measures is that at
fixed $v=v_1$ the moment side is already exhausted ($|\lambda_2-\lambda_{\rm ex}|\le2\cdot10^{-3}$ of
the gap). The three direction runs \texttt{b51}--\texttt{b53} then settle the direction half of that alternative, and settle it
against the programme: the descent set of $v\mapsto\lambda_n(v)$ from $v_1$ has angular width
$t_{\rm hi}\approx2.5\cdot10^{-5}\ldots1.6\cdot10^{-3}$ (so the $64\times64$ net above, and every net
starting at $10^{-3}$, is blind to it by construction --- retraction~(p)), and at the fibre direction
$v_\ast$ --- the minimiser of the \emph{exact} crossing $\lambda_{\rm ex}$, not of $\lambda_1$ ---
the $n=1$ member is still $0.50$--$1.01$ of the gap short while $\lambda_2(v_\ast)$ is within
$5\cdot10^{-4}$ (Remark~\ref{rem:tight}). The ray scan \texttt{b55} then replaces the trade-off that these runs
described with the criterion of Lemma~\ref{lem:margin}, closes the direction question in the
equal-angle cells (Corollary~\ref{cor:equal}: $\min_v\lambda_1(v)=\operatorname{cf}_1$, attained at the
$C$-maximisers and \emph{not} at $v_\ast$), and falsifies the small-$q$ universality of the ray
(retraction~(q)). What remains open is therefore not ``how fast does
$\min_v\lambda_n(v)$ descend on a net'' --- that question was a grid artifact --- but two sharper ones:
whether $\inf_{\|v\|=1}\lambda_1(v)$ ever reaches $\lambda_{\min}(\Gr)$ in a generic table (the
inequality $\lambda_1\ge\lambda_{\min}(\Gr)$ is \eqref{eq:chain1}, and no draw has attained it;
The leverage gate \texttt{b59} narrows it without closing it --- on $264$ rays through $v_1$, minimised exactly and
grid-free by the stationary numerator $G$ of Remark~\ref{rem:tight}(vii), every ray infimum stays at
least $0.2357$ of the gap \emph{above} $\lambda_{\min}(\Gr)$, and $120$ random rays on a $q=4$ draw beat
neither $v_\ast$ nor that floor. The cross-lane run \texttt{b60} then replaced that sampling by the exact sphere gradient of
Remark~\ref{rem:tight}(viii) --- no grid, no finite differences, $60$ starts per draw over five cells ---
and the sphere infimum stays at least $0.3378$ of the gap above $\lambda_{\min}(\Gr)$ on all $50$ draws,
while on $38$ of them a direction strictly below $\lambda_1(v_\ast)$ clears that draw's own evaluator
floor. The infimum is therefore now \emph{searched} rather than sampled, and still open as a proof; and
Remark~\ref{rem:tight}(ix) prices the beat it finds --- $\tfrac12\kappa\theta^2$ with $\kappa$ a measured
sphere curvature and $\theta$ the angle by which $v_\ast$ misses the minimiser --- so the residual
question is to control $\theta$, not to keep hunting for a better direction), and
whether the difference $\lambda_1(v)-\lambda_2(v)$ has a closed-form leading term along $v(t)$, which
is the item \texttt{b53} did not reach. The moment question
("is $n=2$ enough?") is answered: it is enough to within $2\cdot10^{-3}$ of the gap at $n$ equal to
two, and it cannot be pushed past the atom count of $\rho_v$. The
second item is unchanged: whether the $o(\lambda)$ remainder of \eqref{eq:slack} admits a
closed-form bound, so that the residual
of $\operatorname{cf}_{1}$ could be sized without evaluating the defect. As before, no statement in
this chapter uses
$\operatorname{cf}$ or $\lambda^{\star}$, and the competing-branch explanation for the old residual
gap --- the lower branch switching to $\lambda_{\min}(\Le)/2$ --- stays ruled out:
$\min\lambda_{\min}(\Le)/2\ge1.38$ against $\max\mathrm{UB}\le0.99$.
\end{enumerate}

\section{Positioning relative to the closest work}\label{sec:position}

\begin{itemize}\itemsep1pt
\item \emph{Eigenvalue bounds for block and saddle-point matrices}
\cite{bradley2021,bergamaschi2026,bergamaschi2025,helanow2023}. Bounding the spectrum of a block
matrix by extremal eigenvalues and singular values of its blocks, with inertia and
multiplicity statements attached, is \emph{standard} in the preconditioning literature, and
that literature already expresses the real eigenvalues as bounds in terms of the extremal
roots of a family of parametric polynomials. Proposition~\ref{prop:tight} is therefore not
offered as a technique. What it offers is that the bound is \emph{attained} along the
dictionary's own fibre, and that it inverts in closed form --- neither of which a
preconditioning paper needs. This item was added in the round that ran the search, because
the availability of the technique elsewhere lowers what may be claimed for it; the
Gauss-rule side (Proposition~\ref{prop:gauss}) is searched as well, and its downgrade is the same one:
the nodes of the $n$-point fibre rule are the eigenvalues of the Jacobi matrix of the fibre measure ---
the Golub--Welsch construction \cite{golubwelsch1969,vanassche2023,webb2026,zheng2026,li2024},
including the Lanczos--Golub--Welsch route \cite{zheng2026} and the question of when
Lanczos iterations generate symmetric quadrature nodes \cite{li2024} --- and the inequality side of \eqref{eq:chain1}
is Rayleigh--Ritz at $n=1$ and Cauchy interlacing of the Krylov compression at $n\ge2$, so neither the
rule nor the bound is offered as a technique; the interlacing side is used as a standing property in
current numerical work too, where exact Ritz extraction is credited with ``the interlacing property
that makes each band energy an upper bound on the true one'' \cite{gubler2026} and a Ritz-value
interlacing theorem is proved and extended to generalised singular values as a routine tool
\cite{fang2026}. What survives on both sides is the same pair --- attainment
along the fibre and closed-form inversion --- together with the measured losses of
Remark~\ref{rem:tight}, which an interlacing statement does not supply. Small print, because it is the
kind of inference I have had to take back before: the single hit for the phrase \emph{Gauss--Stieltjes}
on a preprint server \cite{tudela2026}, where the rule serves passivity-preserving reduction rather than
certification, is \emph{not} evidence that the statement is unoccupied --- the primary literature for
this technique is in journals and monographs, and the judgement above rests on the seven Golub--Welsch
hits and on the interlacing fact, not on a low count.
\item \emph{Error bounds for Krylov approximations of matrix functions}
\cite{chen2021,xu2022,frommer2012,simunec2023}. This is the closest live
machinery to what Remark~\ref{rem:tight} does, and it cuts \emph{against} this chapter: that work
already derives both a priori and a posteriori bounds for Lanczos approximations of
$f(\mathbf{A})\mathbf{b}$ \emph{and} of the quadratic forms $\mathbf{b}^{\textsf{H}}f(\mathbf{A})\mathbf{b}$,
with bounds that resolve clustered and isolated eigenvalues of $\mathbf{A}$ and are extended to finite
precision \cite{chen2021} (abstract read), and is proposed as a practical stopping criterion
\cite{xu2022}. Since $\operatorname{cf}_1$ here is a two-moment $(A_0,B_2)$ truncation of exactly
that kind of object, ``a certificate computable from low-order spectral data'' is not offered as a
contribution. What survives is the same pair as the item above --- attainment along the fibre and
closed-form inversion --- plus the negative measurement of the \emph{direction} side: that
error has a spectrum-dependent part ($23.5\times$ shift of the median angle across $40$ subspace draws
at fixed cosines, distributions overlapping) and
an orientation part ($1.87$--$3.40$ decades at fixed spectrum), no function of
$\operatorname{spec}(A_0)=\{c_i^2\}$ explains the second, and the angle is not an iteration error at all
but the distance between the exact eigenvector's fibre block and this functional's own minimiser. The
four entries above were read entry by entry (title, authors, abstract), not inferred from a search count.
\item \emph{Principal vectors for subspace pruning in Koopman approximation}
\cite{shah2026}. They use principal angles to measure
\emph{invariance defect}, $\mathcal K(S)$ versus $S$. Our angles are between the prior
and correction \emph{blocks} and measure \emph{attributability}. Contribution framing
matters here: the claim of this note is not ``principal angles in Koopman'', which is
no longer novel, but ``model-form error leaks into the correction block with worst-case
amplification exactly $1/\sin\alpha$, constant $1$, dimension-free, and attained''.
The angle machinery itself is classical and we use it as classical: computation of the angles
between subspaces \cite{bjorckgolub1973}, the perturbation and majorisation theory of those angles
and of Ritz values \cite{knyazevargentati2002,knyazevargentati2006,drmac2000}, and the coherence
obstacle to making $\sin\alpha_{\min}$ large by dictionary design, namely the Welch bound
\cite{welch1974}. None of these is offered as a contribution here; what is offered is that the
amplification factor of the attribution error \emph{is} $1/\sin\alpha$, with constant one.
\item \emph{Eigenvector perturbation theory} \cite{daviskahan1970,wedin1972,kato1995}. A reader
trained on sin-$\Theta$ theorems will look for one here, and it is not what this note contains: the
angle we measure between the exact fibre block and the minimiser of the functional is a
\emph{model-discrepancy} angle, not an eigenvector perturbation, and it does not vanish under
refinement of the estimator (Appendix~\ref{app:direction}). The classical bounds are the right
tools for a different question, and we say so rather than letting the notation suggest a theorem we
did not prove.
\item \emph{Identifiability bounds for linear and nonlinear ODEs from solution data}
\cite{pan2026}. Equation-versus-equation distinguishability via
Hausdorff distance on solution sets and metric entropy, minimax over classes. We are
complementary: attribution of one operator inside one decomposition
$A^\star=P+C+\Delta$, i.e.\ how the same data must be split between two blocks. The
generic phrase ``identifiability theory for SciML'' is no longer available to us; the
claim must read ``decomposition-level identifiability''.
\item \emph{Learning physics from an imperfect ancestor} \cite{mousavi2026}. Structure prior from an
imperfect neural operator plus residual refinement. It
documents that the ``imperfect prior $+$ learned correction'' architecture is active, and
it contains no fidelity or identifiability theory for the correction term. This note
supplies exactly that missing piece for the linear-operator case.
\item \emph{Residual pseudospectra for a physics-informed Koopman backbone}
\cite{lorenzo2026}. Kernel EDMD with residual minimization and
pseudospectral analysis on ERA5/HadISST SST. Their certificate is spectral consistency;
ours is a leakage bound in the dictionary decomposition, with no information quantity and
no coverage interpretation. Together with the operator-identification monograph of Kutz
\emph{et al.} \cite{kutz2016} (the standard reference for identification on standardised data,
which is precisely the practice our methodological note criticizes), these fix the
boundary of what is new here.
\item \emph{The least-squares projection structure and the Frisch--Waugh--Lovell theorem}
\cite{frisch1933,lovell1963,ding2021}. Fitting the correction block after the prior block is, at the
level of the normal equations, residualisation: the estimator of $C$ is formed from the part of the
data that $P$ cannot explain, which is the content of that theorem, and \cite{ding2021} is the
standard care taken with the resulting standard errors. Our object is different in kind: FWL says
what the point estimate is, and this note bounds what an unmodelled operator $\Delta$ is \emph{forced}
to contribute to it, in the worst case, as a function of the dictionary's principal angles. The
distinction is stated because the two results are easy to conflate, and a reader who conflates them
will conclude the leakage theorem is already in the regression literature.
\item \emph{Grey-box modelling and system identification} \cite{ljung2010}. The
``known structure plus parameterised residual'' model class is the grey-box class of system
identification, where the questions of identifiability and of variance reduction are well
developed. This note's contribution is confined to one quantity in that class that the classical
theory does not price: the worst-case transfer of an \emph{unmodelled} operator into the fitted
correction block, which for the linear operator case here is exactly $1/\sin\alpha_{\min}$.
\item \emph{The chord bound of Proposition~\ref{prop:chord} against the truncated moment problem}
(\cite{kovalyov2025,fritzsche2017}; abstracts read
$2026\text{-}10\text{-}01$, not the papers). Both describe the truncated (multi- and matricial) moment
problem \emph{through the Stieltjes transform} --- as an interpolation problem solved step by step with
continued fractions, and via a Schur-type algorithm respectively. Read against them, the function
$s_v(\lambda)=B(v)^{2}/(B(v)-\lambda C(v))$ introduced in the proof of Proposition~\ref{prop:chord} is
exactly the two-moment (one-atom)
approximant of a Stieltjes transform, i.e.\ the kind of object produced by the \emph{first} step of such
an algorithm, and the inequality $\phi_v\ge s_v$ is Cauchy--Schwarz (equivalently Jensen) for a positive
measure with prescribed mass and first moment: an argument I do not claim as new. What those papers do
not do, and what this chapter does, is the \emph{direction of use}: $s_v$ lower-bounds
$f(\lambda)=\lambda_{\max}(K^{\top}(\Le-\lambda I)^{-1}K)$, whose first crossing with $1-\lambda$
\emph{is} $\lambda_{\min}(\Gr)$, so choosing $v$ to maximise the slope on the top eigenspace yields a
closed-form surrogate on a certified side of that eigenvalue --- a statement about a block Gram matrix of
a pair of subspaces, not about a moment solution set. Two things stay recorded as open rather than cited:
whether the specific two-moment bound is stated classically under the name Chebyshev--Markov--Stieltjes
(or as Markov's theorem for Stieltjes transforms, or in Wall/Gautschi's Gauss-rule error-sign treatment)
is \emph{not} verified entry by entry, so no reference is attached to Proposition~\ref{prop:chord}; and if
verification succeeds, the correct edit is to demote the inequality step of its proof to that citation
while keeping the eigensystem construction, which is the part that would then remain.
The two-atom check \texttt{b50} sharpens that caveat in the unfavourable direction, and it is recorded here rather than only in
the log: Proposition~\ref{prop:gauss} replaces the one-atom approximant by the $n$-point Gauss rule on
the \emph{same} measure, whose lower-bound property is exactly the classical sign statement
$\int h\,d\rho\ge G_n(h)$ for $h^{(2n)}\ge0$. If the name-check comes back, what is being cited then
covers the entire inequality machinery of \eqref{eq:hier}, not just its $n=1$ case; the part that could
still remain mine is the identification of the crossing of the resulting rational function with
$1-\lambda$ as a certified surrogate for $\lambda_{\min}(\Gr)$, the direction family
\eqref{eq:family}, and the $\Le^{-j}K$ moments as the computable ingredients --- the quadrature itself
would be classical.
\end{itemize}

\section{Reproduction}\label{sec:reproduce}

\noindent The manuscript, the scripts, their saved output, the climate index tables, and the full
collaboration log of this project are archived at
\url{https://github.com/m-rui001/identifiability-leakage-principal-angles} (MIT licence). The layout is
\texttt{manuscript/}, \texttt{scripts/agent\_a/} and \texttt{scripts/agent\_b/} (the numbered scripts
below), \texttt{evidence/} (the stdout each script produced, named after it), \texttt{logs/} and
\texttt{tools/}. Each number in the tables below is printed by the script named next to it, and the file
under \texttt{evidence/} is that script's saved output; the measurement scripts need numpy and scipy.
The climate index tables are also archived, but they belong to the cyclostationary ENSO attempt that
\S\ref{sec:notclaims} records as \emph{closed}: they are shipped for completeness and as evidence of
what was tried, not as validation.
Where a script ran and the claim was afterwards retracted, the retraction is in
\S\ref{sec:notclaims} and in the log, not quietly deleted.

\noindent\textbf{Geometry, sampling, the certificate.}
\begin{center}\footnotesize\begin{tabular}{@{}>{\raggedright\arraybackslash}p{5.1cm}%
>{\raggedright\arraybackslash}p{2.5cm}>{\raggedright\arraybackslash}p{5.8cm}@{}}
\toprule Script & Objects & Verifies\\
\midrule
\texttt{b11\_theorem\_remainder.py} & Props.~\ref{prop:block}--\ref{prop:square},
Thm.~\ref{thm:sampling} & block Gram to $10^{-16}$; $r(M)$ flatness; $n^{-1/2}$ rate;
$\kappa\approx0.15$\\
\texttt{b13\_two\_sided\_bound.py} & Thm.~\ref{thm:twosided} & $900$ draws, $0$
violations on either side; tightness tables\\
\texttt{b12\_certificate\_no\_constants.py} & $(\star)$, Cor.~\ref{cor:dilution} &
$24/24$ cells pass, ratio $0.117$--$0.175$; adversarial $0.876$--$0.919$\\
\texttt{b14\_g\_alpha\_kappa.py} & Thm.~\ref{thm:sampling} & $\kappa$ calibration,
dictionary-size scan\\
\texttt{b15\_bias\_not\_noise.py} & Thm.~\ref{thm:bias} & sign tails, Jensen gap, $n\beta_2$
prediction ($0.825/0.870/0.947$)\\
\texttt{b24\_ggamma.py} & Thm.~\ref{thm:gammaclosed} & $\alpha:0.25^\circ\to85^\circ$,
$\lambda_{\min}(\Gr)$ over $4$ decades, bias flat in $\alpha$\\
\texttt{b25\_pqn\_law.py} & Thm.~\ref{thm:gammaclosed} & $21$ cells: $\mathrm{sd}\sqrt
n=0.68\pm0.02$; falsifies $c=p+q$ and $O(1/d)$\\
\texttt{b26\_margin\_upgrade.py} & Remark~\ref{rem:marginplace} & random arm is
uninformative ($0.113$--$0.174$, flat in $n$)\\
\texttt{b27\_adversarial\_margin.py} & $(\star)$ & $4320$ adversarial trials, zero
violations, no margin\\
\bottomrule
\end{tabular}\end{center}

\noindent\textbf{The slack factorization.}
\begin{center}\footnotesize\begin{tabular}{@{}>{\raggedright\arraybackslash}p{5.1cm}%
>{\raggedright\arraybackslash}p{2.5cm}>{\raggedright\arraybackslash}p{5.8cm}@{}}
\toprule Script & Objects & Verifies\\
\midrule
\texttt{b28\_cancellation.py} & Thm.~\ref{thm:cancel} & $q_4$ $n$-free to $0.08\%$,
$q_1q_2\in[0.986,1.000]$; leaves the $1-q_3$ column of the table in
\S\ref{sec:cancellation}\\
\texttt{b29\_upper\_bound\_check.py} & Thm.~\ref{thm:twosided} & $2400$ per-sample
draws, $0$ violations of either side; $\lambda_{\min}/\mathrm{UB}\ge0.985$\\
\texttt{b30\_q4\_analytic.py} & Rem.~\ref{rem:q4analytic} & closed-form $q_4$ to
$10^{-4}$ of b28's measured value; $\|\Le^{-1}H\|_2^2=0.128$ at $d=10$\\
\texttt{b31\_q4\_curve.py} & \eqref{eq:q4formula} & eleven $\alpha$'s, $4400$
dictionaries: per-sample deviation $\le0.004$, $0$ upper-side violations\\
\texttt{b32\_residual\_diagnosis.py} & \eqref{eq:factor}, Cor.~\ref{cor:slackone} &
per-trial closure $q_4q_1q_2q_3=0.8905$ vs.\ measured $0.8905$;
$\|v_{\mathcal F}\|=0.9466$ (sd $0.0066$) at $2^\circ$, $0.9993$ at $85^\circ$\\
\texttt{b33\_lambda\_min\_sharp.py} & Prop.~\ref{prop:plane}, Rem.~\ref{rem:cf} & $2700$ draws:
$\lambda_{\min}(\Gr)<\operatorname{cf}$ on every one, relative gap
$\le2.2\cdot10^{-3}$; $\min\lambda_{\min}(\Le)/2=1.38$ against $\max\mathrm{UB}=0.99$\\
\texttt{b34\_correlated\_u.py} & \S\ref{sec:notclaims} item 3 & $\operatorname{AR}(1)$
trajectory design: $q_1\uparrow$, $q_2\downarrow$, $q_1q_2$ down to $0.9747$
($\rho=0.95$, $n=200$), all $\mathrm{err}/\mathrm{bound}<1$, \eqref{eq:factor} closes per
trial ($0.8776$/$0.8776$)\\
\bottomrule
\end{tabular}\end{center}

\noindent\textbf{The exact $\lambda_{\min}(\Gr)$ leg.}
\begin{center}\footnotesize\begin{tabular}{@{}>{\raggedright\arraybackslash}p{5.1cm}%
>{\raggedright\arraybackslash}p{2.5cm}>{\raggedright\arraybackslash}p{5.8cm}@{}}
\toprule Script & Objects & Verifies\\
\midrule
\texttt{b35\_mu\_reduction.py} & Thm.~\ref{thm:exact}, Rem.~\ref{rem:w0} & $3200$ draws, four $(d,p,q)$:
$0$ violations of $\lambda_{\min}(\Gr)\le\mu(w_0)$,
$|\lambda_{\min}(\Gr)/\mu(w_0)-1|\le1.0\cdot10^{-5}$; the row's original ``$0/3200$ random $w$
beating $w_0$'' was a sampling artifact, see \texttt{b37}\\
\texttt{b36\_cf\_chain.py} & Prop.~\ref{prop:plane} & $6300$ draws, $42$ cells: the chain
$\lambda_{\min}(\Gr)\le\mu(w_0)\le\operatorname{cf}<\mathrm{UB}$ holds on every draw, and
$\operatorname{cf}$ equals a brute-force plane minimisation to $4.5\cdot10^{-12}$\\
\texttt{b37\_w0\_optimality.py} & Rem.~\ref{rem:w0} & exhaustive $S^1$ search at $q=2$ vs.\ random
directions at $q=3$: $143/180$ non-square $q=2$ draws at $\alpha\ge20^\circ$ admit $\mu(w)<\mu(w_0)$
($157/840$ overall), separation up to $0.30$ rad, gap $\le9.6\cdot10^{-7}$\\
\texttt{b38\_exact\_scalar\_formula.py} & Thm.~\ref{thm:exact} & $6480$ dictionaries, nine $(d,p,q)$,
$\alpha=1^\circ\ldots89^\circ$: worst $|\lambda^\star/\lambda_{\min}(\Gr)-1|=1.8\cdot10^{-11}$,
$0$ violations; the $\lambda_{\min}(\Le)$ face never active\\
\texttt{b39\_square\_case\_cf\_exact.py} & Cor.~\ref{cor:square} & $7$ square configs $\times$ five
angles, $200$ draws: $|\lambda_{\min}(\Gr)/\operatorname{cf}-1|\le1.4\cdot10^{-11}$,
$\|\Le-C^{-\top}C^{-1}\|/\|\Le\|\le1.5\cdot10^{-15}$; rectangular defects
$1.9\cdot10^{-3}$ ($8,4,3$), $6.9\cdot10^{-3}$ ($6,4,3$)\\
\bottomrule
\end{tabular}\end{center}

\noindent\textbf{The two-constant audit, and the scope checks on $\operatorname{cf}$.}
\begin{center}\footnotesize\begin{tabular}{@{}>{\raggedright\arraybackslash}p{5.1cm}%
>{\raggedright\arraybackslash}p{2.5cm}>{\raggedright\arraybackslash}p{5.8cm}@{}}
\toprule Script & Objects & Verifies\\
\midrule
\texttt{b40\_a\_d10\_dilution\_audit.py} (+\texttt{b40b}) &
Intro, Cor.~\ref{cor:dilution} & audit of A's two-arm design, reproduced to the last digit on
his seed: the model-form constant is the projection fraction ($8\%$ over $d=6\ldots20$, Beta
mean verified to $0.5\%$ by $10^4$ draws), flat in $n$; the noise constant is $\sqrt{q/n}$
($13\%$); $c_2/c_1=d/\sqrt n$; A's ``predicted $c$'' column shown to be the same row's own fit\\
\texttt{b41\_q\_scan\_dilution.py} & Cor.~\ref{cor:dilution} & $q=1,2,3,4,6$ at fixed $d=10$:
$c_1=0.077$--$0.227$, so the dilution is $\kappa(q)\sqrt q/d$ ($c_1$ within $[0.94,0.99]$ of it,
no trend), not $q$-blind $1/d$; the $\alpha$-tilt is shown \emph{not} to discriminate\\
\texttt{b42\_multi\_angle\_exact.py} & Prop.~\ref{prop:multi}, Rem.~\ref{rem:multi} &
prescribed principal-cosine spectra in place of the constructed single angle: the root find holds
to $\le1.6\cdot10^{-12}$ (and to $\le1.3\cdot10^{-15}$ against $\lambda_{\min}=1-c_1$ at $\Le=I_p$),
\eqref{eq:exact} and \eqref{eq:multi} agree to $\le4.8\cdot10^{-13}$ on equal-angle cells, while
$\operatorname{cf}$ stays within $3.4\%$ but flips sign --- $\ge0$ on $240$ equal-angle draws,
$\le0$ on $119$--$120$ of $120$ per wide-spread cell\\
\texttt{b43\_rootfind\_conditioning.py} & Thm.~\ref{thm:exact}, Rem.~\ref{rem:multi}(iii) &
accuracy of the reduction itself, $60$ draws/cell, $\alpha=45^\circ\ldots0.001^\circ$ and
$\kappa(\Le)=1\ldots10^{12}$: $\mathrm{relerr}\times\lambda_{\min}(\Gr)\in[1.6,1.8]\cdot10^{-15}$ flat
over six decades (absolute floor, hence relative accuracy $\sim\sin^{-2}\alpha$); unguarded bracket
endpoint at $10^{-12}$ below the pole gives order-$100$ errors, guarded restores the floor to
$\lambda_{\min}\approx10^{-9}$; at $p=q$ the $\operatorname{cf}$ deviation equals that floor, so it is
implementation, not the closed form\\
\texttt{b44\_cf\_exactness\_attribution.py} & Cor.~\ref{cor:iso}, Rem.~\ref{rem:multi}(ii),
\S\ref{sec:notclaims}(l) & $2\times5$ design (generic vs.\ isotropic $\Le$, five cosine spectra,
$100$ draws/cell) $+\,18$ arbitrary-$K$ cells: at $\Le=\gamma I_p$ the gap is at the floor
($\le2.1\cdot10^{-15}$) for \emph{every} spectrum; the flip is resolution-checked
($9.1\cdot10^{5}$, $1.1\cdot10^{4}$ above $100\epsilon\|\Gr\|_2$, $0$ unresolved); the sign
condition this row first attached to that flip is retracted by \texttt{b45}\\
\end{tabular}\end{center}

\noindent\textbf{The one-pole reading and its replacement leg.}
\begin{center}\footnotesize\begin{tabular}{@{}>{\raggedright\arraybackslash}p{5.1cm}%
>{\raggedright\arraybackslash}p{2.5cm}>{\raggedright\arraybackslash}p{5.8cm}@{}}
\toprule Script & Objects & Verifies\\
\midrule
\texttt{b45\_resolvent\_one\_pole\_\allowbreak diagnostic.py} & Prop.~\ref{prop:side}, Rem.~\ref{rem:side},
Cor.~\ref{cor:iso}, \S\ref{sec:notclaims}(m) & ${\sim}6650$ draws: the sign of
$\operatorname{cf}-\lambda_{\min}(\Gr)$ is the sign of $D(\operatorname{cf})$, $0$ exceptions in $500$
resolved draws, and the gap equals $D/(1+f')$ to $\le1.1\cdot10^{-4}$ relative; $D\equiv0$ at
$\Le=\gamma I_p$ ($15$ cells, $\le1.3\cdot10^{-15}$) and for $ZK=0$, $Z\ne0$
($|t-1|\le1.3\cdot10^{-15}$); flip in $100/100$ draws at a second cosine of $0.087$ and at
$\operatorname{cf}=0.133$, absent in all $12$ single-atom cells\\
\texttt{b46\_true\_slope\_surrogate.py} & Prop.~\ref{prop:chord}, Rem.~\ref{rem:chord},
\S\ref{sec:notclaims} open items & $2100$ draws on \texttt{b45}'s cells with the pole moved to the
true slope: flip $100/100\to0/100$, $98/100\to0/100$, and $0/100$ in all sixteen cells;
$\operatorname{cf}\le\operatorname{cf}_1$ in $100/100$ wherever $\operatorname{slack}$ clears the
floor, coin-flip where $E_0=\mathbb R^q$; $\operatorname{med}|\operatorname{gap}_1|/
\operatorname{med}|\operatorname{gap}|=0.001$--$0.31$ where the linear term dominates, $1.00$--$1.10$
where it vanishes; log-log slope of $|\operatorname{gap}_1|$ vs $\operatorname{cf}_1$ $=2.081$
(pre-registered $2$, falsifier $1$); nulls $|t_1-1/\gamma|\le2.0\cdot10^{-15}$,
$\max|D_1|\le1.3\cdot10^{-15}$\\
\texttt{b47\_theorem\_probe.py} & Prop.~\ref{prop:chord} & $1600$ draws: the Jensen step tested
directly --- $\phi_v-\allowbreak s_v\ge0$ at $\lambda=0.3,0.7,0.95$ of $\operatorname{cf}_1$ for the
chosen $v_1$ and for $200$ random directions per draw, worst $-5.6\cdot10^{-16}$ (isotropy cells, gap
$\equiv0$); the family $\lambda(v)$ of \eqref{eq:family} above $\lambda_{\min}(\Gr)$ in
$320{,}000$/$320{,}000$ tests including the capped branch, $0$ resolved sign flips of
$\operatorname{gap}_1$; records the degeneracy $\dim E_0=3$ at equal angles, where a single returned
eigenvector misses $\max_{E_0}C$ by up to $6.6\cdot10^{-1}$\\
\end{tabular}\end{center}

\noindent\textbf{Searching the family \eqref{eq:family} instead of sampling it.}
\begin{center}\footnotesize\begin{tabular}{@{}>{\raggedright\arraybackslash}p{5.1cm}%
>{\raggedright\arraybackslash}p{2.5cm}>{\raggedright\arraybackslash}p{5.8cm}@{}}
\toprule Script & Objects & Verifies\\
\midrule
\texttt{b48\_family\_min\_search.py} & Prop.~\ref{prop:chord}, Rem.~\ref{rem:chord},
\S\ref{sec:notclaims}(n) & search instead of sample: polar net $96\times240$ per draw,
$7{,}380{,}480$ family values $+$ $144{,}000$ ray points, $\max|\,\|v\|-1\,|\le7.8\cdot10^{-16}$.
$0$ members below $\lambda_{\min}(\Gr)$, but $\operatorname{cf}_1$ is \emph{not} the family minimum:
the first variation along $v(t)=\cos t\,v_1+\sin t\,u$ is
$-2\,v_1^{\top}K^{\top}\Le^{-2}Ku\cdot\lambda(1-\lambda)/|\partial_\lambda F|$, negative in $84/84$
generic draws (measured/predicted relative error $\le1.7\cdot10^{-2}$), while the $\dim E_0=q$ cells
give $0$ competitors and $\|P_{v_1^{\perp}}K^{\top}\Le^{-2}Kv_1\|\le6.2\cdot10^{-16}$; the best ray
removes $0.07$--$0.50$ of $\operatorname{cf}_1$'s gap (worst draw $0.58$), each competitor buying slope
$C/(c^{2}t_1)\in[1.000005,1.000499]$ with intercept $B/c^{2}\le1-10^{-6}$\\
\texttt{b49\_two\_jet\_dip\_formula.py} & Rem.~\ref{rem:chord},
\S\ref{sec:notclaims} open item 1 & sizes that descent in closed form before searching for it:
$280$ saddle draws, $\sim1.6\cdot10^{6}$ family values ($504{,}000$ on the steepest rays, the rest on a
$64\times64$ net over the same draws), $\max|\,\|v\|-1\,|\le4.4\cdot10^{-16}$,
$|\lambda(v_1)-\operatorname{cf}_1|\le3.7\cdot10^{-16}$. The two-jet
$\operatorname{dip}=-L_1^{2}/(4L_2)$ matches the measured ray depth to
$5.9\cdot10^{-6}\ldots1.0\cdot10^{-3}$ relative per cell, $t_{\mathrm{meas}}/t^{*}\in[0.993,1.008]$,
$L_2>0$ in $280/280$, $0$ members below $\lambda_{\min}(\Gr)$ and $0$ draws where the predicted gain
exceeds the gap; one ray step removes $0.0004$--$0.50$ of the gap at the median (worst $0.7642$,
$|\operatorname{dip}|\le1.2\cdot10^{-7}$ absolute), reproducing \texttt{b48}'s $0.07$--$0.50$ on a
different seed by a different method. The net goes $\lesssim10\%$ deeper (wins in $32/40$, $31/40$,
$7/40$, $2/40$ and $0/40$ cells, median $\operatorname{ray}/\operatorname{net}=0.906$--$1.009$, and it
does not contain the steepest direction at $5.6^{\circ}$ --- sampled, not exhaustive). The equal-angle
cell is $40/40$ stationary, where the formula's hypothesis $\|P_{v_1^{\perp}}K^{\top}\Le^{-2}Kv_1\|>0$
simply fails\\
\bottomrule
\end{tabular}\end{center}

\noindent\textbf{Replacing the one-pole chord by a Gauss rule on the same measure.}
\begin{center}\footnotesize\begin{tabular}{@{}>{\raggedright\arraybackslash}p{5.1cm}%
>{\raggedright\arraybackslash}p{2.5cm}>{\raggedright\arraybackslash}p{5.8cm}@{}}
\toprule Script & Objects & Verifies\\
\midrule
\texttt{b50\_gauss\_moment\_\allowbreak hierarchy.py} & Prop.~\ref{prop:gauss}, Rem.~\ref{rem:gauss},
\eqref{eq:cubic} & $280$ draws, seven cells: the rules integrate $1,y,\ldots,y^{2n-1}$ exactly to
$\le7\cdot10^{-15}$ relative (\emph{checked before} any side claim --- two earlier runs of this script
mis-indexed the Stieltjes recurrence and are void), $\lambda_1=\operatorname{cf}_1$ to
$\le7.2\cdot10^{-16}$, and $\lambda_2,\lambda_3,\lambda_{\rm ex},\lambda^{\star}\ge
\lambda_{\min}(\Gr)$ in $280/280$. Ordering $\lambda_1\ge\lambda_2$ in $280/280$,
$\lambda_2\ge\lambda_3$ in $239/239$ finite pairs, pointwise $G_1\le G_2\le G_3\le\phi_v$ in
$240/240$. $\lambda_2(v_1)$ removes $0.26$--$0.9998$ of $\operatorname{cf}_1$'s residual gap
(cell medians $0.50\ldots0.999$) against $0.0004$--$0.72$ for \texttt{b49}'s ray; the explicit cubic
\eqref{eq:cubic} equals the eigenvalue-plus-bisection root to $2.5\cdot10^{-16}$ on $280/280$ and
reproduces a hand two-atom measure. Hierarchy terminates at the atom count: $39$ of $80$ $n=3$ rules
break down on the two near-isotropic cells ($\Le$ has two eigenvalues), and there
$\lambda_2=\lambda_{\rm ex}$ exactly\\
\bottomrule
\end{tabular}\end{center}

\noindent\textbf{Direction versus moment order: what a net can and cannot see.}
\begin{center}\footnotesize\begin{tabular}{@{}>{\raggedright\arraybackslash}p{5.1cm}%
>{\raggedright\arraybackslash}p{2.5cm}>{\raggedright\arraybackslash}p{5.8cm}@{}}
\toprule Script & Objects & Verifies\\
\midrule
\texttt{b51\_net\_lambda2\_\allowbreak direction\_moment.py} & Prop.~\ref{prop:gauss} on a net &
$320$ draws, eight cells, same seed and call order as \texttt{b52}/\texttt{b53}: gates pass ($0$ side
violations, $\max|\|v\|-1|=2.2\cdot10^{-16}$, moment exactness $\le6.6\cdot10^{-15}$ by two independent
routes, $|\lambda_1(v_1)-\operatorname{cf}_1|\le1.7\cdot10^{-16}$, $0$ root-filter rejections), but
$\min_{\rm net}\lambda_1$ removes $0.0000$ of the gap in $7$ of $8$ cells --- and that column is
\emph{retracted as evidence} by \texttt{b52}, retraction~(p). Five defects were caught in a pre-run
audit (three moment matrices where four are indexed, a non-unit $q=2$ perpendicular frame, an
RNG-consuming print, a dead \texttt{nanmin} line, and a false nesting premise); none of them survived
into the run\\
\texttt{b52\_boundary\_layer\_\allowbreak reconciliation.py} & Retraction~(p), $t_{\rm hi}$ &
Same $320$ draws: the \texttt{b50} ray reproduces, and \texttt{b51}'s own angle list on that same ray
returns $0.0000$ --- the net is blind, not wrong. The descent set of $v\mapsto\lambda_n(v)$ from $v_1$
sits inside $t<t_{\rm hi}=2.5\cdot10^{-5}\ldots1.6\cdot10^{-3}$ with $t_{\rm hi}(\lambda_2)\equiv
t_{\rm hi}(\lambda_1)$; a $\operatorname{geomspace}(10^{-8},\pi/2,300)$ list lifts $\lambda_2$ to
$\ge0.9994$ in six of the seven ray cells, $0.9164$ in the $q=4$ one (residual to the analytic ray
$\le8.4\cdot10^{-2}$, typically $3\cdot10^{-6}$),
and rotating the azimuth by $0.2$ rad retains $0.916$--$0.960$ of the descent\\
\bottomrule
\end{tabular}\end{center}

\noindent\textbf{The exact fibre direction, and the margin that orders the family at it.}
\begin{center}\footnotesize\begin{tabular}{@{}>{\raggedright\arraybackslash}p{5.1cm}%
>{\raggedright\arraybackslash}p{2.5cm}>{\raggedright\arraybackslash}p{5.8cm}@{}}
\toprule Script & Objects & Verifies\\
\midrule
\texttt{b53\_exact\_direction\_\allowbreak tightness.py} & Prop.~\ref{prop:tight}, Rem.~\ref{rem:tight},
\eqref{eq:fibre} & Same $320$ draws: \eqref{eq:fibre} as an identity to $\le2.5\cdot10^{-15}$ on
$320/320$ plus the $p$-block row to $\le2.2\cdot10^{-15}$; the intercept defect
$c^{2}-B(v_\ast)\le4.6\cdot10^{-6}$ in the seven $\operatorname{dim}E_0=1$ cells and $\le3\cdot10^{-15}$
in the degenerate one;
$\lambda_2(v_\ast)-\lambda_{\min}(\Gr)\le5\cdot10^{-4}$ of the gap (median per cell) with $0$ side
violations; $\lambda_1(v_\ast)-\lambda_{\min}(\Gr)$ median $0.50$--$1.01$ of the gap --- the refutation
of the $n=1$ conjecture; $\operatorname{angle}(v_\ast,v_1)$ median $1.3\cdot10^{-5}\ldots1.3\cdot10^{-3}$
(a second computation of $t_{\rm hi}$); $\operatorname{angle}(v_\ast,\text{ray})\le5\cdot10^{-4}$ while
$\lambda_1(\text{ray})$ removes only $0.0003$--$0.50$ and $\lambda_2(\text{ray})$ $\ge0.9965$ (six ray
cells; the $q=4$ one $0.9222$), all of them per-cell \emph{medians}; their dispersion is added by
\texttt{b55}, together with the $q$-trend that retracts them as a guarantee\\
\texttt{b55\_fibre\_vs\_chord\_\allowbreak margin.py} & Lem.~\ref{lem:margin}, Cor.~\ref{cor:equal},
\S\ref{sec:notclaims}(q) & $550$ draws in $20$ cells (b53's $320$, plus doubled-$p$ and doubled-$q$
tables on a new seed): the criterion \eqref{eq:margin} has $0$ sign disagreements and all three side
conditions hold on all $550$; \eqref{eq:firstorder} reproduces $\lambda_1(v_\ast)-\operatorname{cf}_1$
at ratio $0.9996$--$1.0003$ per cell, and along the great circle $v(\vartheta)$ ($24{,}000$ points,
$|T|$ up to $1.5\cdot10^{3}$, $\max|\,\|v\|-1\,|\le3.3\cdot10^{-16}$) the sign never flips, the $1\%$
band reaching $|T|\approx3.6\cdot10^{2}$ at $q=3$ and $2.0\cdot10^{1}$ at $q=4$ with departure
$\le5.5\cdot10^{-3}$ for $|T|\le1$. Equal-angle cells: $\max\|A_0-c^{2}I_q\|\le5.8\cdot10^{-15}$,
$B(v)=c^{2}$ with $0$ exceptions over $960{,}000$ grid points, $0$ points below
$\operatorname{cf}_1-\text{floor}$, and $\lambda_1(v_\ast)>\operatorname{cf}_1$ in $24/24$.
The ray leg fails its own registered test: $\lambda_2(\text{ray})$ medians $1.0000$, $0.9995$, $0.9715$,
$0.8932$, $0.9034$, $0.6547$ at $q=2,3,4,5,6,8$ ($20/20$ draws below $0.90$ at $q=8$), and b53's own
$q=4$ cell has per-draw minimum $0.7804$ with $13/40$ below $0.90$. Two defects of my own checks are
recorded in the paired \texttt{.txt} rather than removed: a discriminant compared without its $B^{2}$
factor, and a registered count relation on the wrong side of the iff\\
\bottomrule
\end{tabular}\end{center}

\noindent\textbf{The ray boundary as an algebraic set.}
\begin{center}\footnotesize\begin{tabular}{@{}>{\raggedright\arraybackslash}p{5.1cm}%
>{\raggedright\arraybackslash}p{2.5cm}>{\raggedright\arraybackslash}p{5.8cm}@{}}
\toprule Script & Objects & Verifies\\
\midrule
\texttt{b57\_descent\_boundary\_cubic.py}, \texttt{b57b\_implied\_outlier.py} &
Lem.~\ref{lem:cubic}, Rem.~\ref{rem:tight}(vi) &
$470$ draws in thirteen cells (b53's $320$, two deterministic mix tables, a five-point
near-equal-angle ladder): the cross term $v_1^{\top}K^{\top}\Le^{-1}Ku$ vanishes to
$\le2.1\cdot10^{-15}$, $B$ and $C$ along the ray match their closed forms to $\le7.1\cdot10^{-15}$,
$\Phi(\tan\theta)$ matches the margin to $\le4.7\cdot10^{-14}$ and the cubic vanishes at its own root
to $\le3.1\cdot10^{-12}$; the two routes to the boundary agree in $\lambda$ to $\le5.1\cdot10^{-16}$
apart from two draws pinned at the grid's left endpoint, which \texttt{b57b} shows is a bracket
artefact ($0$ of $40$ draws have more than one descent block). $L_1=\frac{\partial r}{\partial t}A_1$
to $\le4.6\cdot10^{-14}$ relative, the two-jet shift holds to $2\cdot10^{-5}$ where measurable,
$\operatorname{angle}(v_\ast,v_1)/\theta_{\rm hi}$ equal to $0.5000$ in both deterministic tables and
$0.493$--$0.551$ on the sampled cells. Ladder:
slopes falsified for $x_{\rm hi}$ and $P$, $P\le0$ in $13$--$18$ of $30$ near-equal-angle draws.
Three defects of my own checks are recorded in the \texttt{.txt}: a dip ratio registered against a
prediction of the wrong sign, a descent predicate at $10^{-12}\times$gap that sits below the measured
$\lambda$-noise floor, and a $10^{-10}$ left endpoint doing to a geometric grid what
retraction~(p) found in a linear one\\
\bottomrule
\end{tabular}\end{center}

\noindent\textbf{The same family, ray-level and sphere-level (continued).}
\begin{center}\footnotesize\begin{tabular}{@{}>{\raggedright\arraybackslash}p{5.1cm}%
>{\raggedright\arraybackslash}p{2.5cm}>{\raggedright\arraybackslash}p{5.8cm}@{}}
\toprule Script & Objects & Verifies\\
\midrule
\texttt{b59\_setlevel\_gridfree.py} & Lem.~\ref{lem:cubic} as a \emph{set} statement,
Rem.~\ref{rem:tight}(vii), \S\ref{sec:notclaims} open item~1 &
$791{,}736$ ray points ($264$ rays $\times$ $2999$ angles, eight cells): \eqref{eq:raycubic} has
$13{,}757$ sign exceptions, all with $|\lambda_1-\operatorname{cf}_1|\le4.6\cdot10^{-16}$, none above
$10^{-15}$, all at $\tan\theta\le2.2\cdot10^{-7}$, worst one $1.38$ times the measured disagreement of
the two independent root routes ($3.33\cdot10^{-16}$) --- so no set-level exception at my evaluator's
precision. $B(x),C(x)$ match direct evaluation to $3.9\cdot10^{-12}$; the
stationary numerator $G$ is verified exact-against-exact, $9.6\cdot10^{-16}$ relative to $-F_x/F_\lambda$,
sign agreement $5221/5221$, finite differences demoted to context (floor $6\cdot10^{-11}$ to
$5\cdot10^{-10}$ near a stationary point). No ray's grid undercuts the exact ray min by more than
$10\times$ that ray's own measured floor ($0/264$; median excess $8.6\cdot10^{-9}$ of gap), every ray
minimum sits $\ge0.2357$ of the gap above $\lambda_{\min}(\Gr)$ ($0/264$ below), and $0/120$ random
$q=4$ rays beat $v_\ast$. Three measuring-stick defects in the \texttt{.txt} and the script's
REGISTRATION HISTORY: a band six orders under the evaluator's error, a finite difference used as the
reference for an identity, and a candidate list missing $x=0$\\
\addlinespace
\texttt{b60\_sphere\_gradient.py} & Rem.~\ref{rem:tight}(viii), \S\ref{sec:notclaims} open item~1 &
Exact sphere gradient $\operatorname{proj}_{T_vS^{p-1}}(2\lambda_B A_0v+2\lambda_C B_2v)$: verified
against the closed-form ray derivative of (vii) ($6.8\cdot10^{-15}$ relative to $|\lambda|/x$) and the
root against \texttt{np.roots} ($8.9\cdot10^{-15}$). $50$ draws $\times$ $60$ starts: the sphere minimum
is below $\lambda_1(v_\ast)$ on $50/50$, claimed for the $38$ whose beat exceeds $10\times$ that draw's
floor (max $2.1\cdot10^{-6}$ of gap, $=2955$ floors, angle $\le4.4\cdot10^{-5}$ rad);
$0/50$ below $\lambda_{\min}(\Gr)$ beyond floor, deficit $\ge0.3378$ of gap; converged-start spread
$\le7.1\cdot10^{-8}$ of gap ($0/50$ over $10^{-6}$, single basin); $0/50$ grid undercuts, nested
$\le0.0625$ floors; $19$ stops above $10^{-6}|\lambda|$ in gradient, so only function values are claimed\\
\bottomrule
\end{tabular}\end{center}

\noindent\textbf{The sphere minimum read as a law (continued).}
\begin{center}\footnotesize\begin{tabular}{@{}>{\raggedright\arraybackslash}p{5.1cm}%
>{\raggedright\arraybackslash}p{2.5cm}>{\raggedright\arraybackslash}p{5.8cm}@{}}
\toprule Script & Objects & Verifies\\
\midrule
\texttt{b61\_geodesic\_law.py} & Rem.~\ref{rem:tight}(ix), \S\ref{sec:notclaims} open item~1 &
$\lambda_1(v_\ast)-\lambda_1(v_{\rm b})=d_1\theta+\tfrac12\kappa\theta^2$ with each coefficient from two
routes ($\kappa$: Richardson second difference of values vs centred difference of (viii)'s gradient;
ratio med $1.0001$ in every cell). Gates $1.7\cdot10^{-14}$, $2.8\cdot10^{-17}$; $0/50$ negative curvature
at a reported minimum; $0/50$ with linear term $\ge10\%$ of the beat. Law holds with leverage on
$10/50$ --- ratio $1.0000$ ($0.9914$--$1.0008$), residual $\le0.31$ floor --- and the other $40$ are
reported as \emph{no leverage}, not as passes; the pooled slope there is $1.631$ not $2$ because $\kappa$
and $\theta$ anticorrelate within a cell. Retracts [B-CLAIM-62]; the \texttt{.txt} records the v1 floor
being inflated by a truncation term rather than noise, and the fix can only remove passes\\
\texttt{b62\_newton\_step.py} & (ix)'s Newton reading: tested and \emph{stopped}, no claim added;
\S\ref{sec:notclaims} open item~1 &
$\theta=\|\nabla_{\!S}\lambda_1(v_\ast)\|/\kappa$ and $\text{beat}=\|\nabla_{\!S}\lambda_1(v_\ast)\|^2/2
\kappa$, coefficients at $v_\ast$ only. Where the step clears its own floor the ratios are $1.0000$
($0.9992$--$1.0039$) and $\text{beat}/(g^{\ast2}/2\kappa^{\ast})\le1.0037$; $0/50$ violations of pure
Newton, of the cubic-corrected step or of the value form; $\kappa^{\ast}>0$ on $50/50$ and no draw has
$\|\nabla_{\!S}\lambda_1(v_\ast)\|$ exactly zero ($\min4.5\cdot10^{-10}$). It is recorded as stopped
because the recast raised the number of measurable draws from $10/50$ to $11/50$ --- ten of them the
same cell as the row above --- which is this row's own pre-registered condition for dropping it, and
because the two-route floor on $g^{\ast}$ ($\max1.4\cdot10^{-6}$) exceeds $|g^{\ast}|$ itself
($\mathrm{med}\,2.5\cdot10^{-7}$): the bottleneck is measuring $g^{\ast}$, not the law. Fired as
registered: on $24/50$ draws the descent started at $v_\ast$ stops short of the six-start best by
${>}10$ evaluator floors, i.e.\ by $4.3\cdot10^{-11}$ of the gap at the median, $7.5\cdot10^{-9}$ at
worst --- \texttt{b60}'s single-basin statement therefore holds at the gap scale, not at the
evaluator's\\
\bottomrule
\end{tabular}\end{center}

\noindent\textbf{What controls the angle, swept by design (continued).}
\begin{center}\footnotesize\begin{tabular}{@{}>{\raggedright\arraybackslash}p{5.1cm}%
>{\raggedright\arraybackslash}p{2.5cm}>{\raggedright\arraybackslash}p{5.8cm}@{}}
\toprule Script & Objects & Verifies\\
\midrule
\texttt{b63\_theta\_small\_\allowbreak cosine\_sweep.py} & \S\ref{sec:notclaims} open item~1:
designed sweep, both registered directions falsified &
Two families, eight swept values, five draws each; leverage precondition met at all sixteen points
($5/5$ testable at $15$ of $16$, $3/5$ at the sixteenth, whose largest stopping
gradient is instead $6.10\cdot10^{-7}$ --- which is why it is $3/5$; the fifteen full ones top out at
$1.73\cdot10^{-8}$). M1 fires: $\theta$ is not non-decreasing in the two-smallest-cosine ratio (Spearman $+0.548$,
three sign flips, $\mathrm{med}\,\theta$ from $8.96\cdot10^{-5}$ at $r=10^{-4}$ to $2.91\cdot10^{-4}$ at
$r=0.999$ with a return to $6.39\cdot10^{-5}$ midway). M2 fires: $\theta$ is largest where the two
\emph{largest} cosines coincide ($4.56\cdot10^{-3}$ at $d=10^{-4}$ against $3.12\cdot10^{-4}$ at $d=0.40$,
single-point max $7.58\cdot10^{-2}$). The \texttt{.txt} records my a-priori reason as wrong --- it assumed
$\Le^{-1}$ weights the small cosines, but $A_0$ contains no $\Le$ --- and one header label as wrong
(\texttt{lambda\_1(v*)} printed where the code evaluates $\lambda_-$)\\
\texttt{b63b\_spectrum\_vs\_\allowbreak subspace.py} & \S\ref{sec:notclaims} open item~1: spectrum vs
subspace at fixed spectrum; reframes $\theta$ as model discrepancy &
Two cosine spectra, $40$ subspace draws each, all $80/80$ leverage-bearing. Identities checked before use:
$\max\|A_0-M^{\top}M\|=3.33\cdot10^{-15}$, $\max\|B_2-M^{\top}\Le^{-1}M\|=3.89\cdot10^{-15}$,
$\operatorname{spec}(A_0)=\{c_i^2\}$ to $1.95\cdot10^{-14}$ (so $\rho_A=(c_1/c_2)^2$, hence fixed at fixed
cosines). S1 holds as a designed contrast: median $\theta$ $4.912\cdot10^{-3}$ (coincident top pair)
against $2.090\cdot10^{-4}$ (spaced), $23.5\times$, with overlapping ranges
($2.685\cdot10^{-5}$ vs $9.324\cdot10^{-4}$). S2 fires: $1.87$--$3.40$ decades of $\theta$ at fixed
spectrum. S3 is \emph{withdrawn as void}, not fired --- a Spearman against a regressor whose within-group
range is $6.9\cdot10^{-15}$ and $1.8\cdot10^{-14}$ (relative $5.6\cdot10^{-15}$); the surviving statement is the invariance one. D2: $\|\nabla_{\!S}\lambda_1(v_\ast)\|
=8.0\cdot10^{-6}$--$8.5\cdot10^{-4}$, $\ge3.9\cdot10^{8}\times$ its rounding scale, $0/80$ draws
stationary --- so $\theta$ is a model-discrepancy angle, not a convergence error. The history block also
withdraws a per-draw overclaim of my own and records that reading a within-group slope as a refutation
of \texttt{b61} broke \texttt{b61}'s own recorded caution\\
\bottomrule
\end{tabular}\end{center}

\noindent Run with \texttt{PYTHONIOENCODING=utf-8 python B/<script>}; each writes its
table to the paired \texttt{.txt}. Expected numbers are the ones printed in this chapter;
the Monte-Carlo standard error on the bias tables is $\pm0.49$ ($n=200$) and $\pm0.34$
($n=600$) in units of $|\text{bias}|\!\cdot\!n$, which is \emph{smaller} than the
observed spread across dictionary sizes --- that residual spread is carried by
$\kappa(\Le)$, see Remark~\ref{rem:exact}, and is not sampling noise.

\appendix

\section{The direction leg, in full}\label{app:direction}
This appendix is the measurement ledger behind Remark~\ref{rem:dirleg}; it is kept complete
rather than summarised because every figure in it is what retired a claim of mine, and because
the script identifiers are the pointers used by Section~\ref{sec:reproduce}.

\begin{remark}[what \texttt{b51}--\texttt{b55} measured: the residual is truncational, not directional]
\label{rem:tight}
Proposition~\ref{prop:tight} was written by \texttt{b53} to test a stronger conjecture of mine, which
it refuted: that $\min_v\lambda_1(v)=\lambda_1(v_\ast)=\lambda_{\min}(\Gr)$, \emph{i.e.}\ that the
$n=1$ member already attains the certificate at the right direction. Scripts
\texttt{b51\_net\_lambda2\_direction\_\allowbreak moment.py}, \texttt{b52\_boundary\_layer\_reconciliation.py}
and \texttt{b53\_exact\_direction\_tightness.py} run on the same $320$ dictionary draws (identical seed
and call order, eight cells, $q=2,3,4$) and report, in this order:
\emph{(i) \texttt{b51} is a measurement of its own grid.} A $64\times64$ polar net gives
$\min_{\rm net}\lambda_1=0.0000$ of the gap in $7$ of $8$ cells, contradicting the analytic ray of
\texttt{b50}; the net result is not wrong (its own gates pass: $0$ side violations in $320$,
$\max|\|v\|-1|=2.2\cdot10^{-16}$, $0$ root-filter rejections, moment exactness
$\le6.6\cdot10^{-15}$) --- it is blind, for the reason in (ii).
\emph{(ii) The descent set is a boundary layer.} Reconciling the two on the same draws,
\texttt{b52} brackets the directions along which $\lambda_n$ descends from $v_1$ inside
$t<t_{\rm hi}\approx2.5\cdot10^{-5}\ldots1.6\cdot10^{-3}$, with $t_{\rm hi}(\lambda_2)\equiv
t_{\rm hi}(\lambda_1)$: the width is a property of the direction functional, not of $n$. Every grid
whose angle list starts at $10^{-3}$ contains $v_1$ and nothing else in the descent set, so its
minimum is constant in the grid density \emph{by construction}; replacing the list by
$\operatorname{geomspace}(10^{-8},\pi/2,300)$ lifts $\lambda_2$ to $\ge0.9994$ of the gap in six of the
seven cells that carry a ray, and the grid-free route (analytic ray, $n=2$) to $\ge0.9965$ --- the
$q=4$ cell being the single exception at $0.9164$ (grid) and $0.9221$ (ray). The descent is wide in the azimuth
$u$: rotating it by $0.2$ rad retains $0.916$--$0.960$ of the descent --- only the polar resolution failed.
\emph{(iii) At $v_\ast$ the identity holds and the $n=1$ member fails.} \texttt{b53} checks
\eqref{eq:fibre} as an identity (no root finding involved) to
$\|K^{\top}(\Le-\lambda I)^{-1}Ky-(1-\lambda)y\|/((1-\lambda)\|y\|)\le2.5\cdot10^{-15}$ on $320/320$,
together with the $p$-block row on the same vectors ($\le2.2\cdot10^{-15}$). It also measures the
intercept at $v_\ast$: $0<c^{2}-B(v_\ast)\le4.6\cdot10^{-6}$ in the seven $\operatorname{dim}E_0=1$
cells, and $|c^{2}-B(v_\ast)|\le3\cdot10^{-15}$ --- machine zero --- in the
$\operatorname{dim}E_0=3$ cell, whose top eigenspace is degenerate. Against that,
$\lambda_1(v_\ast)$ sits $0.50$--$1.01$ of the gap \emph{above}
$\lambda_{\min}(\Gr)$ (median, per cell) --- not below it, and not erratically, so this is truncation
and not branch error --- while $\lambda_2(v_\ast)$ is within $3.3\cdot10^{-7}$, $8.7\cdot10^{-7}$,
$\ldots$, $5.0\cdot10^{-4}$ of it, with $0$ side violations in $320/320$.
The one systematic exception is the degenerate cell, and \eqref{eq:family} explains it rather than
contradicting it: $\operatorname{cf}_1$ is defined as the member with the largest intercept \emph{and,
among those, the largest slope} \eqref{eq:chord}, and at fixed $B$ the root of \eqref{eq:family}
decreases in $C$ --- now the proved $\partial r/\partial t<0$ of Lemma~\ref{lem:margin}, not a reading
of the quadratic; when $\operatorname{dim}E_0=1$ saturating $B=c^{2}$ pins $v=v_1$ up to sign, but when
$\operatorname{dim}E_0=3$ it leaves a two-sphere of directions, $v_\ast$ among them without any reason
to maximise $C$. Measured: $\lambda_1(v_\ast)>\operatorname{cf}_1$ in $40/40$ draws of that cell and in
$0$ of $280$ elsewhere. That cell is the equal-angle one, and \texttt{b55} shows the degeneracy is not
an accident of the draw: $K=\Le^{1/2}M$ with $M=Q_1\operatorname{diag}(c,\ldots,c)Q_2^{\top}$ gives
$A_0=K^{\top}\Le^{-1}K=M^{\top}M=c^{2}I_q$ \emph{exactly} ($\max\|A_0-c^{2}I_q\|\le5.8\cdot10^{-15}$
over $12$ draws at $q=3$ and $12$ at $q=2$), so $\delta=0$ identically and Corollary~\ref{cor:equal}
applies --- see (v).
\emph{(iv) The two knobs are not interchangeable.} The ray direction $u$ of \texttt{b49} is within
$5\cdot10^{-4}$ rad of $v_\ast$ (two cells: identical to machine precision), yet
$\lambda_1(\text{ray})$ removes only $0.0003$--$0.50$ of the gap while $\lambda_2(\text{ray})$ removes
$0.9965$--$1.0000$ \emph{as a per-cell median} (six of the seven ray cells; the $q=4$ cell gives
$0.9222$). Those are medians, not guarantees: on the same draws the $q=4$ cell's per-draw minimum is
$0.7804$ with $13$ of $40$ below $0.90$ \texttt{(b55)}, and across a $q$-ladder of new tables
\texttt{(b55)} the median falls $1.0000$, $0.9995$, $0.9715$, $0.8932$, $0.9034$, $0.6547$ at
$q=2,3,4,5,6,8$ --- $20/20$ draws below $0.90$ at $q=8$. The ray leg is therefore a small-$q$
phenomenon and retraction~(q) retires the strength I had attached to it; on the very same tables
$\lambda_2(v_\ast)$ still removes $1.00\cdot10^{+00}$ of the gap (median \emph{and} max, $20/20$ cells).
Direction was already essentially right; what bought the gap was the order of the
rule. And $\operatorname{angle}(v_\ast,v_1)$, computed independently of \texttt{b52} from the
eigenvector, has median
$1.3\cdot10^{-5}\ldots1.3\cdot10^{-3}$ --- a second, unrelated computation of $t_{\rm hi}$, which is
the corroborating (not the proving) half of (ii): same draws, same scaling.
\emph{(v) The margin is exact, and it is one number.} Lemma~\ref{lem:margin} turns the trade-off of
(iii)--(iv) into a criterion, and \texttt{b55} tests it on $550$ draws in $20$ cells (b53's eight, plus
$p$- and $q$-doubled tables on a new seed): the sign of
$t-t_1-\frac{\delta}{c^2}\frac{1-t_1\mu}{\mu}$ disagrees with the sign of
$\lambda_1(v_\ast)-\operatorname{cf}_1$ in $0$ of them, each side condition of
\eqref{eq:margin} holds on all $550$, and the collapse \eqref{eq:firstorder} predicts the
\emph{size} to a relative $4\cdot10^{-4}$ (per-cell medians $0.9996$--$1.0003$ for the ratio of
$\lambda_1(v_\ast)-\operatorname{cf}_1$ to its first-order value). That agreement is a statement about
$|T|$, not about the law, so the same script drives $|T|$ up along the great circle
$v(\vartheta)=\cos\vartheta\,v_\ast+\sin\vartheta\,w$, $w\perp v_\ast$: over $24{,}000$ points with
$|T|$ spanning $10^{-8}$ to $1.5\cdot10^{3}$ the sign never flips, while the first-order ratio drifts
to $1.67$, its departure from $1$ being $\le5.5\cdot10^{-3}$ for $|T|\le1$. So the criterion carries,
the quantitative form only within a measured regime, and
$\min_{\|v\|=1}\lambda_1(v)=\operatorname{cf}_1$ in the equal-angle cells by Corollary~\ref{cor:equal}
($0$ competitors below $\operatorname{cf}_1-\text{floor}$ on a $20{,}000$-point circle at $q=2$ and on
$60{,}000$ points at $q=3$, while $\lambda_1(v_\ast)>\operatorname{cf}_1$ in $24$ of $24$ draws).
The upshot for the open item it was meant to settle: retraction~(n) established that
$\operatorname{cf}_1$ is a saddle, not a minimum, so a descent direction exists; (i)--(v) establish
that chasing it cannot recover the residual at $n=1$, because the $n=1$ chord loses
$0.5$--$1.0$ of the gap \emph{even at $v_\ast$}, the direction that attains the exact fibre crossing
of Proposition~\ref{prop:tight} --- which is a different minimiser from $\min_v\lambda_1(v)$, and
whether that infimum reaches $\lambda_{\min}(\Gr)$ is now the open question (Section~\ref{sec:notclaims},
open item~1). The residual of $\operatorname{cf}_1$ is
therefore ordered by the moment side, and the sentence of Remark~\ref{rem:gauss} that attributes it to
``the fixed direction'' is corrected here --- see retraction~(o).
\emph{(vi) \texttt{b57} measures Lemma~\ref{lem:cubic}, and the near-degenerate cells refute the
interpolation I hoped for.} Two independent routes to the boundary of $\{\lambda_1\le\operatorname{cf}_1\}$
on a ray --- the leading positive root of $\Phi$ and a $200$-step bisection of
$\lambda_1(v_\theta)-\operatorname{cf}_1$ over a geometric grid of $3000$ angles on
$[10^{-10},\pi/2)$ --- are compared on $470$ draws in thirteen cells, and the disagreement is reported
in $\lambda$, multiplying by the measured $|\partial\lambda_1/\partial\theta|$ at the boundary. Every
cell satisfies the bookkeeping identity for that product to $0$ digits, and the mismatch is at most
$5.1\cdot10^{-16}$ (cell medians $5\cdot10^{-17}$ to $1.2\cdot10^{-16}$) except for two draws, in the
$2$--$45$--$88$ and $q=4$ cells, whose bisected angle came out at $1.0\cdot10^{-10}$ --- the grid's own
left endpoint, so the predicate fired on a round-off dip and bracketed nothing: the same left-endpoint
artefact as retraction~(p), in a grid I had chosen to avoid it. Read that way, the $\theta$-level
disagreement of $1.4\cdot10^{-6}$ in the shallowest cell is $2.4\cdot10^{-16}$ in $\lambda$, and no
draw contradicts \eqref{eq:margin}; the first-order root shift of the lemma is verifiable only where it
exceeds that floor, and there it holds to a relative $2\cdot10^{-5}$ (four cells).
The equal-angle falsifier fires nowhere: $0$ of $120$ random rays leave $\operatorname{cf}_1$, with
deepest apparent dip $2\cdot10^{-16}$ against a gap of $4.8\cdot10^{-5}$.
Driving the two top cosines together instead ($\cos(20\pm d)$, $d=10^{-3}$ to $5$ degrees, $150$
draws) does \emph{not} collapse the layer: $x_{\rm hi}$ has log-log slope $-0.28$ and $P$ slope $+0.71$
against my pre-registered $0$ (both falsified; $g$ and $D$ pass at $0.02$ and $1.01$), $P\le0$ in
$13$--$18$ of $30$ draws at $d\le0.1^\circ$ with no upper crossing on those rays, and
$\lambda_1(v_\ast)$ then recovers $0.986$--$0.993$ of the gap. The mechanism is the one already flagged
at the definition of $t_1$: as the top eigenvalue of $K^{\top}\Le^{-1}K$ becomes multiple the data stop
selecting $v_1$, so $g$, $P$ and every statement referring to them become solver-chosen. A continuity
statement in $\|K^{\top}\Le^{-1}K-c^{2}I_q\|$ therefore cannot be written with $v_1$ in it; the
parameter that controls this regime is the top-eigenvalue separation, not the norm defect.

\emph{(vii) \texttt{b59} tests \eqref{eq:raycubic} as a \emph{set} statement and turns the ray into a
grid-free inner problem; it narrows open item~1 without settling it.} On $264$ rays through $v_1$
(eight cells, $12$ draws each, rays = the steepest descent direction plus two random ones) and
$2999$ angles per ray --- $791{,}736$ evaluations --- the equivalence \eqref{eq:raycubic} has
$13{,}757$ points where the two sides disagree in sign, all of them with
$|\lambda_1(v_\theta)-\operatorname{cf}_1|\le4.6\cdot10^{-16}$, none above $10^{-15}$, and all at
$\tan\theta\le2.2\cdot10^{-7}$, i.\@ where the two quantities differ only by round-off. That bound is
not impressive by itself, so it is measured against my own evaluator: the note's root formula
$\bigl((B+C)-\sqrt{\operatorname{disc}}\bigr)/(2C)$ subtracts two $O(1)$ numbers to produce $\lambda\approx3.7\cdot10^{-3}$
and so loses about $2.5$ digits, while the rationalised form
$\lambda_-=2B(1-B)/\bigl((B+C)+\sqrt{\operatorname{disc}}\bigr)$ is an independent route through the product of the
roots; those two routes disagree by up to $3.33\cdot10^{-16}$ on the same cloud. The worst exception is
$1.38$ times that noise floor, which is what ``no set-level exception'' means here --- a statement
\emph{weaker} than \eqref{eq:raycubic}, and written weaker on purpose. My first registered band for this
test was $10^{-11}$ of the local gap ($\approx1.8\cdot10^{-17}$), six orders below the evaluator's own
error, and it reported the iff as falsified; the fix was the second root route, not the band.
With $x=\tan\theta$ the two moments are rational, $B(x)=(c^{2}+ax^{2})/(1+x^{2})$ and
$C(x)=(t_1c^{2}+2gx+bx^{2})/(1+x^{2})$, so a ray minimiser of $\lambda_1$ is a root of the cleared
stationary numerator
$G(x,\lambda)=\hat C\lambda^{2}-(\hat B+\hat C)\lambda+\hat B(1-2B)$ with
$\hat B=-2Dx$, $\hat C=2\bigl(g(1-x^{2})+(b-t_1c^{2})x\bigr)$, and $\partial\lambda/\partial x=-F_x/F_\lambda$
on the $\lambda_-$ branch ($F_\lambda=-\sqrt{\operatorname{disc}}$). One caveat got me: differentiating $C$ itself
produces $b-t_1c^{2}$, whereas the combination $C-t_1B$ of the lemma carries $b-t_1a$ --- the two are not
interchangeable, and I conflated them on paper first. $G$ is validated exact-against-exact, against the
derivative of the closed-form root rather than a finite difference, to a relative
$9.6\cdot10^{-16}$ (scale $|\lambda|/x$) on all $264$ rays with the sign matching in $5221$ of $5221$
places where the derivative clears the difference floor; the finite-difference route is reported as
context only, because its own floor is $6\cdot10^{-11}$ to $5\cdot10^{-10}$ near a stationary point
where it is the reference that is unusable, not $G$. Minimising over $\{\lambda_1\le\operatorname{cf}_1\}$
then needs no grid: the candidates are the two endpoints $x=0$ and $x=\infty$, the edges of the descent
intervals (which are single intervals on $264/264$ rays), and the bisected roots of $G$. The registered
test --- a grid may not undercut the exact ray minimum by more than ten times that ray's own measured
two-route floor --- passes $0/264$, the median grid excess being $8.6\cdot10^{-9}$ of the gap and the
worst raw undercut $4.1\cdot10^{-6}$ of the gap in the cell whose gap is smallest, i.\@ where the floor
allows it; the earlier $103$ apparent violations were my candidate list omitting
$x=0$, which is exactly the minimiser on a ray whose descent set is empty.
For open item~1 this gives numbers inside the family it searches: every one of the $264$ ray minima sits
strictly \emph{above} $\lambda_{\min}(\Gr)$, the smallest deficit being $0.2357$ of the gap (per-cell
minima $0.2357$--$1.0000$; the equal-angle cell attains $1.0000$ because there the ray infimum is
$\operatorname{cf}_1$ itself), and $0$ rays fall below the floor. $54$ of $264$ rays undercut
$\lambda_1(v_\ast)$ strictly, but only $24$ by more than $10^{-12}$ in $\lambda$, and the largest beat
outside the equal-angle cell is $4.7\cdot10^{-6}$ of the gap ($2.8\cdot10^{-9}$ in $\lambda$); the
equal-angle beats, up to $8.7\cdot10^{-3}$ of the gap, are the already-known
$\lambda_1(v_\ast)>\operatorname{cf}_1$ of retraction~(o), not a new phenomenon. On one $q=4$ draw,
$120$ random rays give $0$ that beat $v_\ast$, the best ray sitting $0.3009$ of the gap above
$\lambda_{\min}(\Gr)$ against $\lambda_1(v_\ast)$'s $0.2741$.
\emph{What this does not buy, plus a wording I had to take back.} The first version of this sentence
claimed that rays through $v_1$ sweep only a two-plane family rather than the sphere. That is wrong as
stated: every unit $v$ decomposes as $\cos\theta\,v_1+\sin\theta\,u$ with $u\perp v_1$, so \emph{the union
over $u$ of these rays is the sphere}, and
$\inf_{u\perp v_1}\min_\theta\lambda_1(v_{\theta,u})=\min_{\|v\|=1}\lambda_1(v)$ holds exactly. What
\texttt{b59} samples is three directions $u$ per draw (the steepest plus two random), and $120$ on one
$q=4$ table --- a Monte-Carlo sample of a $(p{-}1)$-dimensional set, not a two-plane restriction. The
open item is therefore narrowed by evidence, not by geometry: within the sampled planes the descent set
never reaches $\lambda_{\min}(\Gr)$ and $v_\ast$ is never beaten, while $G$ makes the inner problem exact,
which is what a genuine search over $u$ needs in order to say anything.
\emph{(viii) \texttt{b60} runs that genuine search, and $v_\ast$ loses.} Since the union over $u\perp v_1$
of the rays above \emph{is} the sphere, $\min_{\|v\|=1}\lambda_1(v)$ is a $(p{-}1)$-dimensional problem,
and $G$ makes it grid-free pointwise: with $F_\lambda=-\sqrt{\operatorname{disc}}$ on the $\lambda_-$
branch, $\lambda_B=-(1-\lambda-2B)/F_\lambda$ and $\lambda_C=-(\lambda^2-\lambda)/F_\lambda$, so the
gradient of $\lambda_1(v)=\lambda_-(B(v),C(v))$ is $2\lambda_B A_0v+2\lambda_C B_2v$ projected to
$T_vS^{p-1}$ --- no finite differences. That formula is checked \emph{against} the closed-form ray
derivative of (vii) ($6.8\cdot10^{-15}$ relative to the natural scale $|\lambda|/x$, gate $10^{-8}$) and
the root against \texttt{np.roots} on the same quadratic ($8.9\cdot10^{-15}$). Retraction with
backtracking from $60$ starts per draw, five cells $\times$ ten draws:
\emph{the sphere minimum is strictly below $\lambda_1(v_\ast)$ on all $50$ draws}. The clause of (i)--(v)
that called $v_\ast$ ``a different minimiser from $\min_v\lambda_1(v)$'' was until this script a logical
distinction --- the two minimisers could still have coincided --- and it is now a measured one. The claim
is restricted to the $38$ of $50$
whose beat exceeds ten times that draw's own measured two-route floor, the largest being
$2.1\cdot10^{-6}$ of the gap ($2955$ floors), and the winning direction sits within
$4.4\cdot10^{-5}$ rad of $v_\ast$ (median $<5\cdot10^{-7}$). So $v_\ast$ is near-optimal, not optimal, and
the margin is three orders of magnitude smaller than the effect (vii) removed.
\emph{What the search does not reach}: $0$ of $50$ draws put the sphere minimum below
$\lambda_{\min}(\Gr)$ beyond its own floor, the smallest deficit over all $50$ being $0.3378$ of the gap
(median $0.91938$); $0$ of $50$ draws show two converged starts disagreeing by more than $10^{-6}$ of the
gap (max $7.1\cdot10^{-8}$, i.\@ single basin); a $200\,000$-point random direction grid never undercuts
the optimiser by more than ten floors ($0$ of $50$); and the sphere minimum never exceeds the $3$-ray
minimum of (vii) by more than $0.0625$ floors (max $1.1\cdot10^{-13}$ of the gap), as monotonicity
requires. Nineteen reported minima still had $\|\nabla_S\lambda_1\|>10^{-6}|\lambda|$ at stopping, so the
statements above compare \emph{function values at concrete directions}, not stationarity. Open item~1 is
thus answered in one part and sharpened in the other: the sphere infimum does not touch
$\lambda_{\min}(\Gr)$, and $v_\ast$ is not its minimiser --- what remains is not whether a better
direction exists but \emph{which printable quantity controls the angle between $v_\ast$ and it}, since the
per-cell beat spans three orders of magnitude ($1.3\cdot10^{3}$ to $1.7\cdot10^{6}$ floors) while the
per-cell median angles stay in $[0,1.05\cdot10^{-5}]$ rad.
\emph{(ix) \texttt{b61} turns that residual into a law, and my own registered claim is what it refutes.}
At the sphere minimiser $v_{\rm b}$ let $u$ be the geodesic direction of $v_\ast$ and
$\theta=\angle(v_{\rm b},v_\ast)$; then $\lambda_1(v_\ast)-\lambda_1(v_{\rm b})=d_1\theta+\tfrac12\kappa\theta^2$
with \emph{both} coefficients measured by two independent routes --- $\kappa$ by a Richardson-corrected
second difference of values and by a centred difference of the gradient of (viii), whose ratio is
$1.0001$ (median, every cell) --- and the expansion is then checked against the directly measured beat
over the whole arc. Gates $1.7\cdot10^{-14}$ and $2.8\cdot10^{-17}$. The linear term never reaches $10\%$
of the beat ($0/50$; median $1.7\cdot10^{-3}$), so the stopping-criterion caveat of (viii) does not
contaminate it, and the curvature is positive at all $50$ reported minima ($0/50$ negative beyond its own
floor). Where the quadratic term stands ten times above its own floor --- $10$ of $50$ draws, all in the
near-equal-angle cell, $\theta\le7.3\cdot10^{-5}$ --- the ratio
$\mathrm{beat}/(\tfrac12\kappa\theta^2)$ is $1.0000$ (median), $0.9914$--$1.0008$, residual $\le0.31$ of
the floor. \emph{So the fibre direction's suboptimality for $\lambda_1$ is quadratic in the angle it is
off by, with a curvature I can print} ($\kappa\in[0.150,0.926]$ measured). The claim this refutes is mine:
[B-CLAIM-62] of the previous round asserted that the beat and the angle are not controlled by one
quantity, and one quantity controls them. The other $40$ draws are reported as \emph{no leverage} --- the
floor there is dominated by the linear coefficient's own error times $\theta$, which exceeds
$\tfrac12\kappa\theta^2$ --- so the law is claimed on $10/50$, not on $50/50$. Two warnings travel with
it. \emph{(a) A pooled exponent is not the law's exponent:} on those $10$ draws
$\log\mathrm{beat}$ versus $\log\theta$ has slope $1.631$, not $2$, because within a cell $\kappa$ and
$\theta$ are anticorrelated ($\kappa\in[0.744,0.915]$ falling as $\theta$ grows); the identity is
per-draw and the regression is across-draws. \emph{(b) $\theta$ is exactly zero on $10$ of the $50$ draws},
and the equal-angle cell's median $\theta$ is $0$ --- the same draws whose \texttt{b60} beats sat inside
the evaluator noise, which is the consistency one would want: there $v_{\rm b}=v_\ast$ to machine
precision, as (vii) already predicted from the ray infimum being $\operatorname{cf}_1$ itself. What
open item~1 asks now is narrower than either half of (viii): $\kappa$ is measured and the descent is
$\tfrac12\kappa\theta^2$, so the whole question is \emph{what controls $\theta=\angle(v_\ast,\arg\min_v\lambda_1(v))$},
which (viii) bounded above by $4.4\cdot10^{-5}$ rad on its own draws, and \texttt{b61} by
$7.3\cdot10^{-5}$ on mine. A Newton reading of that $\theta$ was tested and stopped: formally
$\theta=\|\nabla_{\!S}\lambda_1(v_\ast)\|/\kappa$ and $\text{beat}=\|\nabla_{\!S}\lambda_1(v_\ast)\|^2/2\kappa$,
and both hold to $0.4\%$ on the draws where the step clears its own floor ($\theta$ ratio
$0.9992$--$1.0039$, $\text{beat}$ ratio $\le1.0037$), but \texttt{b62} raised the number of measurable
draws from $10/50$ to $11/50$ --- ten of them the same cell \texttt{b61} could already see --- so the
recast is not offered as a result. Its useful by-product is a diagnosis: the measured two-route floor
on $\|\nabla_{\!S}\lambda_1(v_\ast)\|$ ($\max1.4\cdot10^{-6}$ in $\lambda$, largest route disagreement
$3.1\cdot10^{-7}$) exceeds the quantity itself ($\mathrm{med}\,2.5\cdot10^{-7}$), and relative to $g^{\ast}$
that floor is $1.6\cdot10^{-3}$ in the cell where the step is resolvable but $129$ in the equal-angle cell
--- so the bottleneck is measuring that gradient, not the law. One
registered falsifier did fire: on $24/50$ draws the descent started at $v_\ast$ stops short of the
six-start best by more than ten copies of that draw's evaluator floor ($\approx4\cdot10^{-19}$), which
is $4.3\cdot10^{-11}$ of the gap at the median and $7.5\cdot10^{-9}$ at worst --- three to five orders
below the $10^{-6}$-of-gap criterion \texttt{b60} used for the same question. Read as registered, the
single-basin statement therefore holds at the gap scale and not at the evaluator scale.

Open item~1 was then run as a \emph{designed} sweep instead of a search (\texttt{b63},
\texttt{b63b}), and both registered directions were falsified. $\theta$ is not monotone in the ratio of
the two smallest cosines (Spearman $+0.548$, three sign flips over eight swept values, with $5/5$
leverage-bearing draws at every point), and it is instead \emph{largest} where the two \emph{largest}
cosines coincide ($\mathrm{med}\,\theta=4.56\cdot10^{-3}$ at $d=10^{-4}$ against $3.12\cdot10^{-4}$ at
$d=0.40$). Separating spectrum from subspace at fixed spectrum ($40$ draws per spectrum, $80/80$
leverage-bearing, the first design in this thread that lost no draw to resolution) leaves a two-factor
answer and no control variable: the cosine spectrum shifts the median $\theta$ by $23.5\times$
($4.912\cdot10^{-3}$ against $2.090\cdot10^{-4}$) while the two distributions \emph{overlap}
($\min=2.685\cdot10^{-5}$ below $\max=9.324\cdot10^{-4}$), and the subspace orientation alone spreads
$\theta$ over $1.87$--$3.40$ decades at fixed spectrum. The candidate mechanism
$\rho_A=\lambda_1(A_0)/\lambda_2(A_0)$ is then inert \emph{by construction}, not by measurement: with
$M:=\Le^{-1/2}K$ --- the matrix whose singular values are the cosines of Prop.~\ref{prop:multi}, so this
factorisation is not a new lemma --- $A_0=K^{\top}\Le^{-1}K=M^{\top}M$ and $\operatorname{spec}(A_0)=\{c_i^2\}$
(checked to $1.95\cdot10^{-14}$, and to $3.33\cdot10^{-15}$ as an identity in $M$), hence
$\rho_A=(c_1/c_2)^2$ cannot vary at fixed cosines. My within-spectrum Spearman against $\rho_A$ was
therefore a correlation against rounding noise (within-group ranges $6.9\cdot10^{-15}$ and
$1.8\cdot10^{-14}$) and is withdrawn
as void rather than reported as a fired falsifier; what the same data do license is the invariance
statement --- $\operatorname{spec}(A_0)$ fixed, $\theta$ spreading $1.87$--$3.40$ decades, so no function
of $A_0$'s spectrum explains the within-spectrum spread. The last number is the one that changes the
question: $v_\ast$ is \emph{not} a critical point of $\lambda_1$
($\|\nabla_{\!S}\lambda_1(v_\ast)\|=8.0\cdot10^{-6}$--$8.5\cdot10^{-4}$, at least $3.9\cdot10^{8}$
times its own rounding scale, on $80/80$ draws). So $\theta$ is the angle between the exact eigenvector's
fibre block and \emph{this functional's own} minimiser --- a model-discrepancy angle, not a convergence
error --- and it should not be read as evidence that the descent fails. In this design the only
legitimate test of \texttt{b61}'s quadratic law remains its per-draw ratio
$\mathrm{beat}/(\tfrac12\kappa\theta^2)$ with $\kappa$ measured by the two independent routes, and that
has not been run; the within-group slopes ($+1.984$, residual geometric s.d.\ $0.091$; $+1.539$,
$0.618$) are reported only to show that they mean nothing on their own, per warning (a) above.
\end{remark}

\begin{remark}[what \texttt{b50} measured, and what the hierarchy does not buy]
\label{rem:gauss}
Script \texttt{b50\_gauss\_moment\_\allowbreak hierarchy.py} runs this on $280$ dictionary pairs in seven cells
($q=2,3,4$, four cosine spectra, two near-isotropic families; $40$ draws each,
$\operatorname{dim}E_0$ listed
per cell), after checking on every draw that the rules do integrate $1,y,\ldots,y^{2n-1}$ exactly
(max relative deviation $\le7\cdot10^{-15}$) --- a rule that integrates no polynomial exactly is not a
Gauss rule, and a side violation under it would be my bookkeeping rather than the theorem.
Result: $\lambda_1=\operatorname{cf}_1$ to $\le7.2\cdot10^{-16}$; $\lambda_2,\lambda_3\ge
\lambda_{\min}(\Gr)$ in $280/280$ (per-draw floor $100\epsilon\|\Gr\|_2$); $\lambda_1\ge\lambda_2$ in
$280/280$ and $\lambda_2\ge\lambda_3$ in $239/239$ finite pairs, and the pointwise ordering
$G_1\le G_2\le G_3\le\phi_v$ in $240/240$ test points --- reported as measured, since the error
formula certifies each rule separately and does not order successive $n$.
The size is what makes it a proposition and not an exercise: $\lambda_2(v_1)$ removes $0.26$--$0.9998$
of the residual gap $\operatorname{cf}_1-\lambda_{\min}(\Gr)$ (cell medians
$0.50$, $0.63$, $0.79$, $0.81$, $0.87$, $0.998$, $0.999$), where the steepest-descent ray of
\texttt{b49} removes $0.0004$--$0.72$; one extra scalar moment buys an order of magnitude more than
optimising the direction did. The cubic form \eqref{eq:cubic} was checked against the
eigenvalue-plus-bisection value on all $280$ draws, max difference $2.5\cdot10^{-16}$, and against a
two-atom measure computed by hand.
Two limits, both stated by the numbers rather than assumed. \emph{(i) The hierarchy terminates:}
$\rho_v$ carries at most $p$ atoms, and where $\Le$ has few distinct eigenvalues the rule stops
early --- in the two near-isotropic cells $\lambda_2$ is already \emph{exact} for $\phi_v$ and $39$ of
$80$ $n=3$ rules break down ($h_2=0$). \emph{(ii) Higher $n$ is not where the residual lives:} write
$\lambda_{\rm ex}(v_1)$ for the crossing of $\phi_{v_1}$ itself, which the hierarchy attains once $n$
reaches the number of atoms of $\rho_{v_1}$. Over the same
draws $|\lambda_2(v_1)-\lambda_{\rm ex}(v_1)|$ is at most $2.0\cdot10^{-3}$ of the gap
$\operatorname{cf}_1-\lambda_{\min}(\Gr)$, while $\rho_{v_1}$ carries $2$--$7$ atoms, so the moment
truncation is negligible and
the rest of the gap is the \emph{fixed direction} --- \emph{this last clause is superseded}: at the
exact minimiser $v_\ast$ of Proposition~\ref{prop:tight} the $n=1$ chord still loses $0.50$--$1.01$ of
that gap, so what is left is the order of the rule rather than the direction; see
Remark~\ref{rem:tight} and retraction~(o). Since \eqref{eq:hier} holds for every unit $v$,
$\min_v\lambda_2(v)$ over any finite set of directions is again a certified bound, and
$\lambda_{\min}(\Gr)\le\min(\lambda_2(v_1),\min_v\lambda_2(v))$; what remains open is how fast that
minimum descends, which is the one-dimensional fractional programme of
Section~\ref{sec:notclaims}~(1) with the moment side now essentially free. What the hierarchy does
\emph{not} do is rehabilitate $\operatorname{cf}$: these draws include the cells where
$\operatorname{cf}<\lambda_{\min}(\Gr)$ ($195$ of $280$), the no-angle model for which
Proposition~\ref{prop:chord} claims only $\lambda_{\min}(\Gr)\le\operatorname{cf}_1$.
\end{remark}

\begin{thebibliography}{99}
\small

\bibitem{bunch1978}
J.~R. Bunch, C.~P. Nielsen and D.~C. Sorensen, \emph{Rank-one modification of the symmetric
eigenproblem}, Numer. Math. 31 (1978) 31--48, \doi{10.1007/BF01396012}.

\bibitem{cuppen1980}
J.~J.~M. Cuppen, \emph{A divide and conquer method for the symmetric tridiagonal eigenproblem},
Numer. Math. 36 (1980) 177--195, \doi{10.1007/BF01396757}.

\bibitem{gu-eisenstat1995}
M. Gu and S.~C. Eisenstat, \emph{A divide-and-conquer algorithm for the symmetric tridiagonal
eigenproblem}, SIAM J. Matrix Anal. Appl. 16 (1995) 172--191, \doi{10.1137/S0895479892241287}.

\bibitem{stor2013}
N. Jakov\v{c}evi\'{c} Stor, I. Slapi\'{c} and J.~L. Barlow, \emph{Accurate eigenvalue decomposition of
arrowhead matrices and applications}, arXiv:1302.7203 [math.NA] (2013).

\bibitem{golubwelsch1969}
G.~H. Golub and J.~H. Welsch, \emph{Calculation of Gauss quadrature rules}, Math. Comp. 23 (1969)
221--230, \doi{10.1090/S0025-5718-69-99647-1}.

\bibitem{gautschi1996}
W. Gautschi, \emph{Orthogonal polynomials: applications and computation}, Acta Numerica 5 (1996)
45--119, \doi{10.1017/S0962492900002622}.

\bibitem{gautschi2004}
W. Gautschi, \emph{Orthogonal Polynomials: Computation and Approximation}, Oxford University Press,
2004, \doi{10.1093/oso/9780198506720.001.0001}.

\bibitem{szego1939}
G. Szeg\"{o}, \emph{Orthogonal Polynomials}, Amer. Math. Soc. Colloquium Publication 23, 1939,
\doi{10.1090/coll/023}.

\bibitem{daviskahan1970}
C. Davis and W.~M. Kahan, \emph{The rotation of eigenvectors by a perturbation.~III}, SIAM J. Numer.
Anal. 7 (1970) 1--46, \doi{10.1137/0707001}.

\bibitem{wedin1972}
P.-\AA{} Wedin, \emph{Perturbation bounds in connection with singular value decomposition}, BIT 12
(1972) 99--111, \doi{10.1007/BF01932678}.

\bibitem{kato1995}
T. Kato, \emph{Perturbation Theory for Linear Operators}, Classics in Mathematics, Springer, 1995,
\doi{10.1007/978-3-642-66282-9}.

\bibitem{bjorckgolub1973}
{\AA}. Bj\"{o}rck and G.~H. Golub, \emph{Numerical methods for computing angles between linear
subspaces}, Math. Comp. 27 (1973) 579--594, \doi{10.2307/2005662}.

\bibitem{knyazevargentati2002}
A.~V. Knyazev and M.~E. Argentati, \emph{Principal angles between subspaces in an $A$-based scalar
product: algorithms and perturbation estimates}, SIAM J. Sci. Comput. 23 (2002) 2008--2040,
\doi{10.1137/S1064827500377332}.

\bibitem{knyazevargentati2006}
A.~V. Knyazev and M.~E. Argentati, \emph{Majorization for changes in angles between subspaces, Ritz
values, and graph Laplacian spectra}, SIAM J. Matrix Anal. Appl. 29 (2006) 15--32,
\doi{10.1137/060649070}.

\bibitem{drmac2000}
Z. Drmac, \emph{On principal angles between subspaces of Euclidean space}, SIAM J. Matrix Anal.
Appl. 22 (2000) 173--194, \doi{10.1137/S0895479897320824}.

\bibitem{welch1974}
L. Welch, \emph{Lower bounds on the maximum cross correlation of signals}, IEEE Trans. Inform.
Theory 20 (1974) 397--399, \doi{10.1109/TIT.1974.1055219}.

\bibitem{frisch1933}
R. Frisch and F.~V. Waugh, \emph{Partial time regressions as compared with individual trends},
Econometrica 1 (1933) 387--407, \doi{10.2307/1907330}.

\bibitem{lovell1963}
M.~C. Lovell, \emph{Seasonal adjustment of economic time series and multiple regression analysis},
J. Amer. Statist. Assoc. 58 (1963) 993--1010, \doi{10.1080/01621459.1963.10480682}.

\bibitem{ding2021}
P. Ding, \emph{The Frisch--Waugh--Lovell theorem for standard errors}, Statist. Probab. Lett. 168
(2021) 108945, \doi{10.1016/j.spl.2020.108945}.

\bibitem{ljung2010}
L. Ljung, \emph{Perspectives on system identification}, Annu. Rev. Control 34 (2010) 1--12,
\doi{10.1016/j.arcontrol.2009.12.001}.

\bibitem{kutz2016}
J.~N. Kutz, S.~L. Brunton, B.~W. Brunton and J.~L. Proctor, \emph{Dynamic Mode Decomposition:
Data-Driven Modeling of Complex Systems}, SIAM, 2016, \doi{10.1137/1.9781611974508}.

\bibitem{bradley2021}
S. Bradley and C. Greif, \emph{Eigenvalue bounds for double saddle-point systems},
arXiv:2110.13328 [math.NA] (2021).

\bibitem{bergamaschi2025}
L. Bergamaschi, A. Martinez, J.~W. Pearson and A. P\"{o}tschka, \emph{Eigenvalue bounds for
preconditioned symmetric multiple saddle-point matrices}, arXiv:2506.02816 [math.NA] (2025).

\bibitem{bergamaschi2026}
L. Bergamaschi, M. Bergamaschi and J.~W. Pearson, \emph{Eigenvalue bounds for preconditioned
symmetric multiple saddle-point matrices with block-triangular preconditioners}, arXiv:2608.24203
[math.NA] (2026).

\bibitem{helanow2023}
C. Helanow and J. Ahlkrona, \emph{Theoretical results on a block preconditioner used in ice-sheet
modeling: eigenvalue bounds for singular power-law fluids}, arXiv:2307.14688 [math.NA] (2023).

\bibitem{vanassche2023}
W. Van Assche, \emph{A Golub--Welsch version for simultaneous Gaussian quadrature},
arXiv:2309.11864 [math.CA] (2023).

\bibitem{webb2026}
M. Webb and G. Maierhofer, \emph{Stable Hermite transforms via the Golub--Welsch algorithm},
arXiv:2604.02041 [math.NA] (2026).

\bibitem{zheng2026}
B. Zheng, \emph{A finite-precision Lanczos--Golub--Welsch route to probability-table construction in
resonance self-shielding}, arXiv:2603.27715 [math.NA] (2026).

\bibitem{li2024}
W. Li, Z. Han and S. Zhu, \emph{When Lanczos iterations generate symmetric quadrature nodes?},
arXiv:2401.17757 [math.NA] (2024).

\bibitem{chen2021}
T. Chen, A. Greenbaum, C. Musco and C. Musco, \emph{Error bounds for Lanczos-based matrix function
approximation}, arXiv:2106.09806 [math.NA] (2021).

\bibitem{xu2022}
Q. Xu and T. Chen, \emph{A posteriori error bounds for the block-Lanczos method for matrix function
approximation}, arXiv:2211.15643 [math.NA] (2022).

\bibitem{frommer2012}
A. Frommer, K. Kahl, Th. Lippert and H. Rittich, \emph{$2$-norm error bounds and estimates for
Lanczos approximations to linear systems and rational matrix functions}, arXiv:1204.5052 [math.NA]
(2012).

\bibitem{simunec2023}
I. \v{S}imu\v{n}ec, \emph{Error bounds for the approximation of matrix functions with rational Krylov
methods}, arXiv:2311.02701 [math.NA] (2023).

\bibitem{gubler2026}
M. Gubler, T. Park, A. Bussy, J. Laute, N. Marzari, I. Timrov and L. Grigori, \emph{Randomized block
Davidson eigensolvers for plane-wave density-functional theory}, arXiv:2608.24529 [math.NA] (2026).

\bibitem{fang2026}
K. Fang and Z. Jia, \emph{The refined joint bidiagonalization method and an implicitly restarted
algorithm for large GSVD computations}, arXiv:2608.04360 [math.NA] (2026).

\bibitem{shah2026}
D. Shah and J. Cort\'{e}s, \emph{A unified algebraic framework for subspace pruning in Koopman
operator approximation via principal vectors}, arXiv:2603.29001 [math.OC] (2026).

\bibitem{pan2026}
Y. Pan and H. B\"{o}lcskei, \emph{Recovering governing equations from solution data: identifiability
bounds for linear and nonlinear ODEs}, arXiv:2606.27285 [cs.LG] (2026).

\bibitem{mousavi2026}
S.~M. Mousavi, T. Kadeethum, N. Bouklas and S. Goswami, \emph{Learning physics from an imperfect
ancestor}, arXiv:2609.24947 [cs.LG] (2026).

\bibitem{lorenzo2026}
P. Lorenzo-Sanchez, M.~J. Colbrook and A. Navarra, \emph{Residual pseudospectra reveal a
physics-informed Koopman backbone for tropical Pacific variability and ENSO prediction},
arXiv:2606.09369 [math.NA] (2026).

\bibitem{kovalyov2025}
I. Kovalyov, \emph{Multidimensional moment problem and Stieltjes transform}, arXiv:2501.05860
[math.CA] (2025).

\bibitem{fritzsche2017}
B. Fritzsche, B. Kirstein, T. Schr\"{o}der and C. M\"{a}dler, \emph{On the truncated matricial
Stieltjes moment problem $\mathsf{M}[[\alpha,\infty);(s_j)_{j=0}^m,\leq]$}, arXiv:1707.06000
[math.CV] (2017).

\bibitem{tudela2026}
M. Tudela-Pi and I. Colombaro, \emph{Passive two-plateau relaxation from Tricomi confluent
hypergeometric kernels}, arXiv:2604.11464 [math-ph] (2026).

\end{thebibliography}

\end{document}
